\section{The genuine native history of an actual original normal phase}

\begin{theorem}[An ordinary upper test internally produces its whole native geometric history]
\label{hir:g1-proper-normal-upper-native-test}
\uses{hir:g1-proper-cylinder-kept-original-lower-class,hir:g1-proper-cylinder-section-local-class,hir:g1-proper-normal-native-whole-input,hir:native-section-geometric-history-existence}
\lean{Hironaka.HironakaG1ProperNormalPhaseNative.exists_upper_native_test_history}
\leanok
Ordinary original domain data for $Q$ and a GIVEN genuine test $T$
of the actual upper $P=\operatorname{properNormalMarked}(Q)$ internally
produce a genuine native zero-plane history of the SAME length,
including every empty chronological stage. It has the actual
geometry-only Lift, all prescribed restricted centre/projection
maps, genuine native testhood, full transversality and ENTIRE
marked-record pullback compatibility for the actual section record
$V=\mathcal C(\operatorname{bar}(\operatorname{sectionMarked}(Q)))$.
Its initial native whole record pulls back under the genuine original
zero-plane equivalence to $\mathcal C(\operatorname{bar}(Q))$.
No native Lift, native test, source presentation or scalar/ideal
comparison certificate is a premise. No native canonicality follows
from this test-only assertion.
\end{theorem}
\begin{proof}
The ordinary proper-normal/section theorem gives the full local
class $P\sim S$ for the genuine section record $S$. The held
ordinary full-boundary bar and coefficient class theorems give
$S\sim\operatorname{bar}(S)\sim V$, with their actual native marks.
Thus the GIVEN actual $P$ test is an actual $V$ test, internally.
The genuine lower metadata gives the tight domain of $V$ and its
proper section carrier, supported on the ACTUAL regular zero plane.
The complete pulled boundary is transverse to that plane by the
held genuine section geometry, with no empty-boundary specialization.

Apply the held ordinary native-test history existence theorem to
this actual $V$ and this actual test. Its induction supplies the
native strict endpoint, actual native manifold and tower, all
restricted centres/projections and the geometry-only Lift. The
Lift's held length, square, geometry and whole-record laws give
each displayed field. The ENTIRE initial input comparison is
exactly the previously proved genuine zero-plane input theorem,
at the original whole ideal and true double native mark. Every
scalar action here is the literal inclusion pullback, retaining
all holomorphic germ representatives rather than a support or
radical substitute. This uses no numerical CAN comparison.
\end{proof}

\begin{theorem}[The actual original zero phase internally gives a native canonical lower tail]
\label{hir:g1-proper-normal-zero-phase-native-history}
\uses{hir:g1-proper-normal-upper-native-test,hir:g1-proper-normal-original-zero-phase,hir:native-section-geometric-history-existence,hir:native-original-same-bound-history}
\lean{Hironaka.HironakaG1ProperNormalPhaseNative.exists_original_normal_phase_native_history}
\leanok
Assume positive ORIGINAL base dimension $n>0$, ordinary original $Q$
of bound $k$, and the actual ORIGINAL StepOne canonical upper history
and TRUE zero/nonspecial endpoint branch of the preceding zero-phase
theorem. Keep its ACTUAL clock $r$, head $H$, tail $S$, head point,
mask $O$, and prescribed normal-contact lower $B$.

The actual regular native support at the head and a genuine native
history of the SAME restricted tail $S|_O$ are internally produced.
The native input is the ENTIRE inclusion pullback of the actual
historical lower $B$, including its proper carrier, complete ordered
boundary and TRUE double native mark. It has the genuine bound-$k$
domain, ORIGINAL bound-$k$ native CAN and testhood. The actual
geometry-only tail Lift gives the exact length and projection
square, ENTIRE transformed marked-record identity, whole cosupport
preimage and the ORIGINAL full value formula
\[
 I_B(S|_O,z)=\operatorname{sameIndexMap}(1,k)
              \bigl(I_{B_{\rm native}}(S_{\rm native},z)\bigr)
\]
at every native endpoint over the entire actual lower cosupport.
There is no supplied head atlas, native Lift, domain/CAN, scalar,
whole-ideal, clock or desired value packet.
\end{theorem}
\begin{proof}
The actual original canonical-upper theorem supplies genuine $P$
testhood and hence head testhood. Apply the preceding ordinary test
factory to that ACTUAL head. It internally constructs its strict
native zero-plane subset, closed regular of dimension $n$, together
with its actual Lift and FULL boundary transversality. D24 internally
constructs the actual normal lower domain and CAN on $S|_O$.

The initial proper normal contact is contained in the actual zero
plane: its literal normal coordinate vanishes, so the one-component
normal projection is zero. Monotonicity of the genuine strict
transform then puts $B.N$ inside the derived head plane. Restrict
that ACTUAL plane to the actual $O$. The held open-section theorem
gives its true regular native atlas and closedness. The whole lower
boundary is the literal restriction of the current transformed
INITIAL bar boundary. Its original boundary list equals $V.E$;
the held chronological boundary law and head Lift geometry therefore
derive transversality to the FULL actual open-section boundary,
including every born and empty occurrence.

Actual lower CAN gives its genuine test. Apply the held ordinary
native-test existence theorem again, now to this ACTUAL historical
lower $B$, this DERIVED current native plane and this SAME restricted
tail. It constructs the true native tail Lift internally. Native
domain follows from the held actual-domain theorem. The held
original SAME-bound history theorem consumes the derived Lift and
D24's actual lower CAN to prove ORIGINAL native CAN and the exact
same-index full-value formula. Ambient dimension is $n+1$ and
native dimension is $n$, so the actual dimension shift is one.
The Lift's length, square, whole-record and cosupport laws supply
the other entire chronological assertions. No coefficient is
commuted with a controlled transformation: the input is precisely
the ACTUAL lower formed from the transformed INITIAL bar at the
actual head, and all scalar maps are its genuine native inclusion.
\end{proof}

This is a genuine native CAN extraction for the actual ORIGINAL
zero-phase lower. Positive native dimension is explicit; the
zero-native-dimensional empty-history branch remains governed by
the held support-bound/zero-dimensional theorems. Positive original
levels keep their actual selected boundary-member lower, and the
general evaluator keeps its actual order-companion phase. Neither
is replaced by this zero-phase normal lower. No fixed initial
kept-$R$ CAN, original base-$\operatorname{bar}(Q)$ CAN at a later
head, or independent ambient-graph/HS-family comparison is inferred.

\begin{lemma}[The upper bar cosupport lies in the strict normal section]
\label{hir:g1-proper-normal-bar-in-section}
\uses{hir:g1-proper-normal-upper-native-test,hir:g1-proper-cylinder-section-local-class}
\lean{Hironaka.HironakaG1PositiveNormalPhase.bar_cosupport_subset_strict_zeroPlane}
\leanok
For a marked ideal $Q$ in the original invariant domain on $U$, let
$P=\operatorname{properNormalMarked}(Q)$. Along any actual test tower
$T$ of $P$, the cosupport of $\operatorname{bar}(P)_T$ is contained
in the strict transform of the zero section over $U$.
\end{lemma}
\begin{proof}
Write $V=\mathcal C(\operatorname{bar}(\operatorname{sectionMarked}(Q)))$.
Local weak equivalence identifies the cosupports of
$\operatorname{bar}(P)_T$ and $V_T$. Indeed, $\operatorname{bar}(P)$
and $P$ are locally weakly equivalent, as are $P$ and the section
record, the section record and its bar, and that bar and $V$.
The native lift along the very same upper test dynamically contains
the carrier of $V_T$, since the initial carrier of $V$ lies in the
zero section. Its endpoint is the strict transform of that section.
\end{proof}

\begin{lemma}[Cosupport of the positive selected-member lower]
\label{hir:g1-positive-normal-member-cosupport}
\uses{hir:g1-proper-normal-bar-in-section}
\lean{Hironaka.HironakaG1PositiveNormalPhase.positive_lower_cosupport}
\leanok
Let $H$ be an actual test prefix of $P$ and let $p$ be a point of its
final space. In the positive branch of Step~1, on the actual mask
$O=\operatorname{maskA}(P,H,p)$, the cosupport of the lower marked
ideal is the trace on $O$ of
\[
 \operatorname{cosupp}(\operatorname{bar}(P)_H)
 \ \cap\!\!\bigcap_{\substack{D\text{ original at }H\\p\in D}}\!D.
\]
Every point of that cosupport lies in the strict normal section and
has the same original boundary membership as $p$.
\end{lemma}
\begin{proof}
The selected-member construction takes the coefficient ideal on the
intersection of the upper carrier with the selected boundary
members. Its cosupport is the upper cosupport intersected with those
members, and open restriction takes the inverse image under
$O\hookrightarrow H$. The selected indices describe exactly the
original members containing $p$. Membership in the mask excludes
all original members missing $p$. The upper cosupport lies in the
strict normal section by the preceding lemma.
\end{proof}

\begin{theorem}[The actual head of a positive normal phase]
\label{hir:g1-positive-normal-phase-head}
\uses{hir:g1-positive-normal-member-cosupport,hir:g1-proper-normal-upper-native-test}
\lean{Hironaka.HironakaG1PositiveNormalPhase.original_normal_positive_phase_head}
\leanok
Let $T$ be a canonical upper Step~1 history of $P$ for the original
depth-$k$ lower evaluator, and let $y$ be a nonspecial cosupport point
of positive original boundary level. Split $T$ at the start of the
actual run of $y$ at that level, and write $H$ and $S$ for its head and
tail and $p=S(y)$. The head zero section has a closed, regular
strict transform of dimension $n$, transverse to the full transformed
bar boundary. The point $p$ belongs to that section and to the mask
of the selected-member lower ideal. At $p$ the lower ideal is on its
cosupport. The actual restricted tail is canonical for that lower
ideal in its original ambient evaluator, with $y$ in its final
cosupport.
\end{theorem}
\begin{proof}
The canonical-tail theorem gives mask membership, lower domain,
canonicality and endpoint cosupport at the actual head and tail.
The regime at the head has the same positive level as at $y$, and
$p$ lies in the upper bar cosupport there. The selected members all
contain $p$, so the preceding cosupport identity puts the masked
point over $p$ in the initial lower cosupport. Upper canonicality
implies testhood. Apply the upper native-history theorem to the
actual head to obtain the regular strict zero section and its
transversality. The full transformed boundary agrees with the bar
boundary because the original ordered boundary lists agree and
boundary transformation depends only on the tower and that list.
\end{proof}

\begin{theorem}[The positive lower has a section-supported weak class]
\label{hir:g1-positive-normal-member-section-class}
\uses{hir:g1-proper-cylinder-section-local-class,hir:g1-separate-carrier-original-member-stage-local,hir:g1-proper-normal-upper-native-test}
\lean{Hironaka.HironakaG1PositiveNormalPhase.positive_lower_section_class}
\leanok
At any actual test prefix $H$ of $P$, use the selected original
members through $p$ and the positive mask of $P$ at $p$.
On this same open set, the positive lower of $P$ is locally weakly
equivalent to the member cut of the transformed bar of the
section-supported record. The carrier of this second cut lies in
the strict normal section on the mask. On the native manifold of
that section, its inclusion pullback has cosupport equal to the
inverse image of the original positive lower cosupport.
\end{theorem}
\begin{proof}
The padded and section-supported records begin in the same local
weak class. The class is preserved by forming their bars and by
transforming along the common actual test prefix. Their original
boundary lists agree, so their cuts select the same original
boundary members. The member-cut comparison for separate parent
carriers identifies the resulting local weak classes on the same
mask. The carrier of the second cut is contained in its transformed
bar carrier. The latter is the strict transform of a subset of the
initial zero section, hence lies in its strict transform. The upper
native history makes this section regular, including after restriction
to the mask. The section cut is well formed by the member-restriction
and open-restriction formulas. Consequently its native inclusion
pullback preserves cosupport, and the weak-class equality identifies
that cosupport with the trace of the original lower cosupport.
\end{proof}

The section-supported member cut can have a different carrier dimension.
Its local weak class supplies test and cosupport comparisons. Comparison
of the native recursive values requires separate dimension and
positive-branch analysis.

\begin{theorem}[The actual native test tail of a positive normal phase]
\label{hir:g1-positive-normal-native-test-tail}
\uses{hir:g1-positive-normal-phase-head,hir:g1-positive-normal-member-section-class,hir-disjoint-native-original-open-lifts,hir:native-section-actual-native-test,hir:native-section-finite-whole-records}
\lean{Hironaka.HironakaG1PositiveNormalPhase.original_normal_positive_phase_native_test_tail}
\leanok
In the setting of the positive normal phase theorem, take the member cut of
the section-supported bar on the same mask. On the strict normal section
over that mask, its native inclusion pullback has a test history whose
centres and blowup maps come from the actual restricted upper tail.
The history uses the ordered boundary left after removing the selected
members. Its final marked record is exactly the inclusion pullback of
the transformed section cut, with the whole ideal and its mark. The
given endpoint lifts to a point of the final native cosupport.
\end{theorem}
\begin{proof}
The upper canonical history is a test of the padded record, so its
zero-section lift is available for the whole history. Cut that lift at the
actual positive-phase clock and restrict it to the upper mask. This gives
a lift with the whole head boundary pulled back to the mask. The head
section is transverse to this boundary. The section-supported member
cut retains an ordered sublist: the removed slots are precisely the
selected original members. This sublist is SNC and transverse to the
head section. Forget the removed slots in the lift, retaining every
subsequent exceptional slot and the very same blowups.

The ambient canonical lower tail is a test of the padded member cut.
Its weak-class comparison with the section member cut makes this same
tail a test of the section cut. The cut is supported on the head section,
and its ideal has holomorphic germs. The native test and whole-record
transform formulas apply to the reduced boundary lift and the actual
section cut. The centre thinness needed in the lift comes from the
original upper history. The given endpoint belongs to the whole strict
normal section by the upper cosupport containment. Weak equivalence
puts it in the transformed section cut's cosupport on the mask, and
the native cosupport formula transfers that membership to its lift.
\end{proof}

\begin{lemma}[Padded positive special points stay special on the section]
\label{hir:g1-positive-normal-special-of-padded}
\uses{hir:g1-proper-cylinder-section-local-class,hir:g1-positive-normal-member-cosupport}
\lean{Hironaka.HironakaG1PositiveNormalPhase.positive_normal_special_of_padded}
\leanok
At a point of positive original level at an actual upper test prefix,
speciality for the padded record implies speciality for the section
record.
\end{lemma}
\begin{proof}
The two bar transforms have the same cosupport and original boundary.
The section transform's carrier lies in the padded transform's carrier
by monotonicity of strict transforms. At positive level a point special
for the padded record is therefore special for the section record: near
that point, the member intersection of the smaller carrier is contained
in the member intersection of the larger carrier, which is already
in the common cosupport.
\end{proof}

\begin{lemma}[The section positive mask is contained in the padded mask]
\label{hir:g1-positive-normal-section-mask-subset}
\uses{hir:g1-positive-normal-special-of-padded}
\lean{Hironaka.HironakaG1PositiveNormalPhase.positive_normal_section_mask_subset}
\leanok
At a point of positive original level at an actual upper test prefix,
the Step~1 mask for the section record is contained in the padded
record's mask at that point.
\end{lemma}
\begin{proof}
The boundary exclusion conditions agree. The section's forbidden
special set at the anchor level contains the padded forbidden set
there. Taking interiors gives the inclusion.
\end{proof}

\begin{theorem}[The native lower domain on the section's nonspecial mask]
\label{hir:g1-positive-normal-native-domain-on-section-mask}
\uses{hir:g1-positive-normal-section-mask-subset,hir:g1-proper-normal-upper-native-test,bcalc:member-restriction-domain}
\lean{Hironaka.HironakaG1PositiveNormalPhase.positive_normal_native_domain_on_section_mask}
\leanok
Let $Q$ have original invariant bound $k$. At an actual upper test
prefix, suppose the point has positive original level and is nonspecial
for the section record. Its section mask contains the point. Over this
mask the native inclusion pullback of the section's selected-member
lower record is in the original invariant domain of bound $k-1$.
\end{theorem}
\begin{proof}
The original level agrees for the padded and section records. The
positive member-restriction domain theorem for the section record
gives membership in its mask and the claimed lower domain in the
ambient space. The section lower carrier lies in the strict normal
section on this mask. The upper native lift makes the section regular
and transverse to the entire transformed boundary. Transversality
persists on the mask and for the retained ordered boundary sublist.
The native domain theorem for the supported record now transfers
well-formedness, the dimension bound, the true mark and the
nonzero-along condition to the native inclusion pullback.
\end{proof}

\begin{theorem}[The lower invariant domain on a positive member mask]
\label{hir:g1-selected-domain-on-positive-mask}
\uses{bcalc:member-locus-centre,bcalc:member-restriction-domain}
\lean{Hironaka.HironakaG1PositiveNormalSpecialStep.selected_member_nonzero_failure_locus,Hironaka.HironakaG1PositiveNormalSpecialStep.isInvDomain_selected_iff_special_locus_empty,Hironaka.HironakaG1PositiveNormalSpecialStep.isInvDomain_selected_pullback_val_iff}
\leanok
On a mask that meets only original boundary members through the
anchor, fix the member cut at that anchor. The locus where its
coefficient ideal vanishes identically near a point of the cut
carrier is exactly the original special locus at the anchor's
level. For a positive anchor level the cut belongs to invariant
domain $k-1$ exactly when this locus is empty. A nonempty locus
excludes every recursive bound. On any open subwindow the
restriction of the same cut belongs to the lower invariant
domain exactly when the subwindow misses this locus.
\end{theorem}
\begin{proof}
A point of the selected carrier lies on the same original
members as the anchor. Away from their member locus, the
coefficient ideal is nonzero along their intersection. At a
point of the special locus the incident members agree with
the anchor's by the mask condition and equality of counts.
The selected carrier near this point is contained in the
cosupport of the cut. Coherence extends each germ of the cut
ideal to nearby stalks, where the positive mark makes that
germ vanish on the carrier. Thus the nonzero-along condition
fails there. The positive member cut already has the required
dimension bound and mark. Open restriction preserves
well-formedness, the dimension bound and nonvanishing along
the carrier, which gives the subwindow criterion.
\end{proof}

\begin{lemma}[The transformed selected record after a positive special step]
\label{hir:g1-selected-domain-after-special-step}
\uses{bcalc:member-locus-centre,bcalc:member-restriction-domain,hir:g1-selected-domain-on-positive-mask}
\lean{Hironaka.HironakaG1PositiveNormalSpecialStep.isInvDomain_selected_geoTransform_of_special_step}
\leanok
Let an invariant-domain record of bound $k$ have an ordered sublist of
boundary members and an anchor on a positive number of them. Suppose
every point lies on only members through the anchor. Form the selected
lower record at the anchor. Blow up the member locus at the anchor's
level by a blowup admissible for the selected record. Its transformed
carrier avoids the exceptional fibre over the centre, and the
transformed record belongs to the invariant domain of bound $k-1$.
\end{lemma}
\begin{proof}
At a point of the centre, the maximal-level member locus agrees
locally with the carrier intersected with the members through that
point. These members also pass through the anchor, so a point of the
selected carrier in this neighbourhood belongs to the centre. Pulling
back the local containment shows that the closure of the inverse
image of the selected carrier minus the centre avoids the fibre.

Each remaining point projects to the selected carrier off the centre.
The positive mask criterion gives nonvanishing of the selected
coefficient ideal there, and this condition transfers along the blowup
off the centre. Taking a
nonempty member cut gives the dimension bound $k-1$, which persists
under the blowup. Admissibility gives well-formedness, and the positive
mark stays fixed.
\end{proof}

\begin{theorem}[The native positive special step on the padded mask]
\label{hir:g1-positive-normal-native-special-step}
\uses{hir:g1-proper-normal-upper-native-test,hir:g1-positive-normal-member-cosupport,hir:g1-positive-normal-member-section-class,hir:g1-selected-domain-after-special-step,hir:native-section-original-positive-lower-record,hir:native-section-original-member-locus,hir:native-section-original-invariant-domain,hir:native-section-finite-whole-records,bcalc:member-locus-centre,bcalc:member-locus-level}
\lean{Hironaka.HironakaG1PositiveNormalSpecialStep.positive_normal_native_special_step}
\leanok
At an actual upper test prefix take a point of positive original level $s+1$ that is
nonspecial for the padded record and special for the section record.
On the padded mask, restrict the transformed section bar to the native
strict section with its full boundary. The point lifts to the member
locus of the native bar at level $s+1$, counting the pullbacks of the
original boundary members. This entire level locus is an admissible
centre. Its image is admissible for both the ambient section bar and the
padded positive lower record on the mask. A proper ambient blowup
along this image restricts on the strict section to the proper native
blowup along the native centre. The native transformed record is the
literal pullback of the ambient transformed section bar. Its cosupport
has empty member locus at the same level, counting the strict transforms
of the original members.

The native pullback of the selected section record is exactly the
member cut of the native bar at the lifted anchor. The same native
blowup is admissible for this cut, and its whole controlled transform
is the pullback of the ambient section cut transform. The ambient
section and padded cut transforms have the same local weak class.
The transformed native cut lies in invariant domain $k-1$ and its
carrier lies over the complement of the native centre.
The inclusion maps its cosupport onto the entire transformed padded
cut cosupport, which lies over the complement of the ambient centre.
\end{theorem}
\begin{proof}
Nonspeciality for the padded record puts the anchor in its mask. The
mask excludes every original boundary member missing the anchor. Thus
the count everywhere on it is at most $s+1$. The original boundary and
its level agree for the padded and section records. Restriction to the
mask preserves the bar domain. Its carrier lies in the strict normal
section, which is regular and transverse to the full transformed
boundary by the upper test lift. The native domain theorem applies to
the literal inclusion pullback. The native member-locus identity sends
the anchor to the native centre and the count is preserved by the two
inclusions. The anchor lies on a native boundary member, so the
native dimension is positive by simple normal crossings.

The maximal-level centre theorem gives admissibility both on the section
and for the ambient locus with the same level. Since the ambient bar is
supported on the section, its locus is exactly the image of the native
centre. Each point of this centre has the anchor's original membership,
by the mask condition and equality of the counts. Thus the padded lower
cosupport formula gives the centre containment for the member cut.
Normal crossings persist for the ordered boundary sublist of the cut.

The invariant domains give empty interior of the two bar cosupports.
The ambient centre admits a proper blowup, and the native projection of
its strict section is a genuine proper blowup along the native centre.
The whole-record transform identity applies to this one-step lift. The
one-step member-locus theorem clears the specified native level.

Restriction to the mask and pullback to the native section each
commute with the full selected-member construction. The ambient
section and padded cuts have the same local weak class, so the padded
test is also a test for the ambient section cut. Retain its ordered
boundary sublist in the one-step lift. The native test and whole-record
identities for this lift give admissibility and the entire transformed
cut. The mask membership condition transfers to the native bar, so
the positive special-step lemma gives the lower invariant domain and
the avoidance of the exceptional fibre by the selected carrier.
The section cut remains supported on the strict section. Its native
cosupport is the inclusion preimage, and its ambient cosupport equals
the padded cut cosupport by local weak equivalence. The projection
square transfers avoidance of the centre to this full padded cosupport.
\end{proof}

\begin{theorem}[The positive section and padded run lengths]
\label{hir:g1-positive-normal-section-run}
\uses{hir:g1-positive-normal-special-of-padded,hir:regime-run}
\lean{Hironaka.HironakaG1PositiveNormalChronology.positive_normal_section_run,Hironaka.HironakaG1PositiveNormalChronology.positive_normal_section_run_cutoff,Hironaka.HironakaG1PositiveNormalChronology.positive_normal_section_restart_centre}
\leanok
Fix a positive original level on an actual padded test history. The
section record's run at that level has length at most the padded
record's run. The lengths agree precisely when all the preceding
points counted by the padded run are nonspecial for the section.
For a strictly shorter section run, the preceding stage is
section-special and padded-nonspecial at the fixed level. If the
endpoint is section-nonspecial, the next actual blowup centre
contains the projection of the endpoint to that preceding stage.
\end{theorem}
\begin{proof}
All prefixes are padded tests. The section and padded counts agree at
each prefix, and positive padded speciality implies section
speciality. Hence each of the preceding points counted by the section
run satisfies the padded regime. If the section run were longer, the
preceding point immediately before the padded run would satisfy the
padded regime, contrary to maximality. The same maximality argument
for the section shows that its run has the full padded length when
none of those preceding points is section-special. The converse
follows from the defining property of the section run.
At the preceding stage the section regime fails and the padded
regime holds, forcing section speciality by equality of levels.
The consecutive splits give the actual intervening blowup and
the exact projection of that point. Speciality is unchanged
off an admissible centre. Since the section point is nonspecial
after this step, its projection must belong to the centre.
\end{proof}

\begin{theorem}[The values at a section-special positive endpoint]
\label{hir:g1-positive-normal-native-special-values}
\uses{hir:g1-positive-normal-section-run,hir:g1-proper-normal-upper-native-test,hir:g1-proper-normal-bar-in-section,hir:native-section-original-step-one-chronology,hir:step-one-value-masked-basic}
\lean{Hironaka.HironakaG1PositiveNormalChronology.positive_normal_native_special_values}
\leanok
Take a padded cosupport point on an actual padded test history, at
positive original level $s+1$. Suppose it is nonspecial for the
padded record and special for the section record. Its lift to the
literal native history has original Step~1 value $(s+1,\top)$
for the native section record. The padded Step~1 value at the point
has the same first component and a finite second component. The
native run at this level has length at most the padded run.
\end{theorem}
\begin{proof}
The point lies on the strict section by upper cosupport containment.
The native chronology for the section-supported record identifies the
section level, speciality and run with their native counterparts.
The comparison of run lengths applies. The native value now follows
from the special branch of Step~1. For the padded record, its
nonspecial cosupport point reads the lower value at its actual mask
and tail. This lower value is finite by the Step~1 domain lemma.
\end{proof}

\begin{theorem}[A level-one example with different original Step One branches]
\label{hir:g1-positive-normal-special-example}
\uses{hir:g1-positive-normal-native-special-values,hir:g1-proper-cylinder-section-local-class}
\lean{Hironaka.HironakaG1PositiveNormalSpecialExample.exists_positive_padded_nonspecial_section_special,Hironaka.HironakaG1PositiveNormalSpecialExample.exists_positive_normal_stepOne_value_discrepancy,Hironaka.HironakaG1PositiveNormalSpecialExample.exists_strict_positive_normal_section_run}
\leanok
Take the whole complex line as carrier, its origin as the sole old
divisor, the vanishing ideal of the origin and mark one. At the empty
tower the padded origin is a nonspecial cosupport point of level one.
The same point is special on the section. The lower coordinate of the
padded Step~1 value is finite and the native lower coordinate is top.
Consequently any comparison preserving finiteness of lower
coordinates fails from padded to native. A comparison obtained by
mapping the lower coordinate via $\operatorname{WithTop.map}$ fails in
either direction.
Appending a single empty-centre blowup at this example gives
section run zero and padded run one.
\end{theorem}
\begin{proof}
The padded carrier is the plane, the section carrier has normal
coordinate zero, and the common bar cosupport is the origin. The
section carrier meets the old divisor only at the origin. The padded
carrier meets the old divisor in a full normal line. Every
neighbourhood of the origin contains a nonzero point of that line,
so the padded member locus does not contain the origin. Apply the
actual native value comparison to the empty tower. The two different
lower coordinates exclude the stated comparisons, since
$\operatorname{WithTop.map}$ preserves both top and finiteness.
The empty-centre blowup is the identity map. The preceding
padded point satisfies its regime and the preceding section
point fails its regime, yielding the two run lengths.
\end{proof}

\begin{theorem}[The shared positive tail on the smaller section mask]
\label{hir:g1-positive-normal-shared-tail}
\uses{hir:g1-positive-normal-phase-head,hir:g1-positive-normal-member-section-class,hir:g1-positive-normal-native-domain-on-section-mask,hir:native-section-finite-whole-records}
\lean{Hironaka.HironakaG1PositiveNormalSharedTail.original_normal_positive_shared_tail}
\leanok
Suppose the anchor at the start of the padded positive run is
nonspecial for the section record. Restrict the padded lower record
and its actual tail to the section's smaller positive mask. The
restricted padded tail is canonical. The section record on that same
mask has the same local weak class as the padded cut and gives an
actual native test tail. Its whole transformed record is the native
pullback of the transformed section cut, its native input bound is
$k-1$, and the given endpoint lifts to the native cosupport.
\end{theorem}
\begin{proof}
The anchor belongs to both masks. Exact nested open restriction
identifies the padded selected record and transports its canonicality
to the section mask. Weak equivalence transfers the test condition
to the section cut. Restrict the actual strict-section geometry
and retain the ordered boundary sublist of this cut. The native
test and whole-record formulas apply to this very lift. The
native domain on the section mask and the cosupport formula give
the remaining assertions.
\end{proof}

\begin{theorem}[Positive powers in a fixed normal chart]
\label{hir:g1-fixed-normal-chart-powers}
\uses{hir:g1-proper-cylinder-section-local-class}
\lean{Hironaka.HironakaG1NormalWeightedClass.markedPullback_power_properNormal,Hironaka.HironakaG1NormalWeightedClass.sectionMarked_restrictEquiv_iff,Hironaka.HironakaG1NormalWeightedClass.restrictEquiv_commonPower_properNormal_locWPEquiv,Hironaka.HironakaG1NormalWeightedClass.commonPower_properNormal_locWPEquiv,Hironaka.HironakaG1NormalWeightedClass.power_properNormal_cosupp,Hironaka.HironakaG1NormalWeightedClass.commonPower_properNormal_test_iff}
\leanok
Literal zero-section pullback takes the whole marked power of a
padded record to the whole source power. On section-supported
records, restricted equivalence is equivalent to restricted
equivalence of the original source records. Restricted equivalence
of two positive source powers therefore gives local weak equivalence
of the two padded powers with an arbitrary common additional
ambient boundary list. Equality of whole source powers is a special
case and gives agreement of test histories. The cosupport of a
positive padded power is the zero-section image of the source
cosupport.
\end{theorem}
\begin{proof}
Coherent pullback commutes with powers. Pulling back along the zero
section gives back the source ideal, carrier, ordered boundary and
actual mark. The projection and zero section transport both
restricted inclusions. Lift the two powered restricted records to
section-supported records and use the padded-to-section weak
equivalence. Transitivity gives the comparison, and prepending a
common ambient list preserves it. Power invariance of the cosupport
gives the last assertion.
\end{proof}

\begin{theorem}[Current companions of the supported normal history]
\label{hir:g1-actual-normal-order-companion}
\uses{hir:g1-proper-normal-controlled-carrier-history,hir:native-original-companion-record}
\lean{Hironaka.HironakaG1NormalOrderCompanion.current_companionOf_on_kept_carrier,Hironaka.HironakaG1NormalOrderCompanion.exists_original_current_order_companion_history}
\leanok
On an actual padded test history, restrict the current record to the
strict carrier of the supported normal record. Computing a current
companion on this carrier gives the whole companion of the
supported record. The associated native test of the original source
has the same length. Transport through its section identifies the
whole current companion at every order and reads the exact order
mask. At a native cosupport point of positive current order both
corresponding masked companions have the claimed invariant domain,
maximal order, equal true marks and the corresponding cosupport
points.
\end{theorem}
\begin{proof}
Controlled carrier restriction gives restricted equivalence at the
current stage, which determines the companion on the kept carrier.
The literal native test transports this supported record back to
the original source on the same stages. Whole native pullback
commutes with the companion, and order transport identifies its
mask. The positive-order guard and the order bound on the open mask
give domain and maximal order. The companion formula gives the same
mark, and equality of the anchor order with the selected order gives
both cosupport memberships.
\end{proof}

\begin{theorem}[Normal flags and recursive depth padding]
\label{hir:g1-normal-flag-depth-padding}
\uses{hir:original-family-actual-depth-stability,hir:iterated-relative-normal-original-canonical-history}
\lean{Hironaka.HironakaG1NormalValueMaps.sameIndexMap_add,Hironaka.HironakaG1NormalValueMaps.depthPadMap_commutes_sameIndexMap,Hironaka.HironakaG1NormalValueMaps.depthPadMap_normal_increment,Hironaka.HironakaG1NormalValueMaps.depthPadMap_iteratedValueEmbedding,Hironaka.HironakaG1NormalValueMaps.paddedCanonicalNormalFlag,Hironaka.HironakaG1NormalValueMaps.paddedCanonicalNormalFlag_comparisons}
\leanok
Same-depth ambient dimension shifts add under composition.
Recursive depth padding commutes with these shifts, with the
one-coordinate normal embedding and with the whole finite-flag
embedding. Thus an actual normal-flag value comparison persists
when both recursive depth bounds are enlarged. Both histories stay
canonical. On the mapped cosupport, strict order, weak order and
equality are reflected. When the map covers the whole target
cosupport, its image of the maximum locus is the target maximum
locus.
\end{theorem}
\begin{proof}
The component maps commute by recursion on their value types. The
same-depth dimension map adds normal dimensions on its monomial
coordinate, and the coefficient branch follows from recursion on
the lower value. Induct along the finite flag using commutation
for one coordinate. Apply actual depth stability on both towers
to transport their comparison. Strict monotonicity of the finite
flag embedding, together with cosupport coverage, gives the
assertions about comparisons and the full maximum sets.
\end{proof}

\begin{theorem}[Maximum loci and canonicality on actual same-bound lifts]
\label{hir:g1-same-bound-native-maxima-canonicality}
\uses{hir:native-original-same-bound-history,hir:native-section-finite-whole-records}
\lean{Hironaka.HironakaG1NormalCentreTransport.Lift.original_same_bound_max_loci,Hironaka.HironakaG1NormalCentreTransport.Lift.original_same_bound_canonical_iff,Hironaka.HironakaG1NormalCentreTransport.Lift.zero_native_isEmptyTowers,Hironaka.HironakaG1NormalCentreTransport.Lift.zero_native_canonical,Hironaka.HironakaG1NormalCentreTransport.Lift.zero_native_cosupports,Hironaka.HironakaG1NormalCentreTransport.Lift.zero_native_max_loci}
\leanok
For an actual native lift of a supported invariant-domain record
and positive native dimension, an ambient canonical prefix has
exact image and preimage formulas for its maximum locus. The same
ambient and native histories are canonical at the same bound
exactly together, including empty steps. In native dimension zero,
both histories have entirely empty centres. For a supported
invariant-domain record their current cosupports and maximum loci
are empty.
\end{theorem}
\begin{proof}
The complete ambient cosupport lies on the current strict section.
Same-bound value comparison is strictly monotone there. It
preserves and reflects every comparison used to define the
maximum, which gives both set identities. Induct on the lift.
The native maximum centre maps onto the ambient maximum centre,
and an empty native centre gives an empty ambient centre. In
dimension zero discreteness and thinness force emptiness of the
native centres. Zero-dimensional invariant-domain pullback has
empty cosupport, and the whole cosupport formula transfers this to
the ambient record.
\end{proof}

\begin{theorem}[The literal native lower record at a nonspecial positive endpoint]
\label{hir:g1-positive-normal-aligned-native-tail}
\uses{hir:g1-positive-normal-clock-equality,hir:g1-positive-normal-own-native-tail,hir:native-section-original-step-one-chronology}
\lean{Hironaka.HironakaG1PositiveNormalAlignedTail.positive_normal_aligned_native_tail_records,Hironaka.HironakaG1PositiveNormalAlignedTail.positive_normal_native_tail_records_of_nonspecial,Hironaka.HironakaG1PositiveNormalAlignedTail.positive_normal_section_head_nonspecial_of_endpoint}
\leanok
At a section-nonspecial padded cosupport endpoint at the
specified positive level, the literal native history has the
same run length as the padded history. The section anchor at
that head is also nonspecial. Its own positive
mask is the exact native preimage of the section mask at the
common head. The native lift of the open restricted ambient
tail is isomorphic to the original native tail restricted to
its own mask. Along the actual inclusions at the head and
endpoint, the native lower record and its whole transform are
the pullbacks of the corresponding section lower records.
\end{theorem}
\begin{proof}
Use persistence of positive speciality and native section
chronology to identify the runs. Split the actual lifted history
at that clock. The native
mask formula gives identical restrictions, and the open-window
comparison gives the tower isomorphism. The whole positive
lower-record head and tail formulas then apply to the literal
inclusions. Their projection square identifies the given endpoint.
\end{proof}

\begin{theorem}[Canonicality and values for a supported lift at different bounds]
\label{hir:g1-supported-native-unequal-depth}
\uses{hir:g1-same-bound-native-maxima-canonicality,hir:g1-normal-flag-depth-padding,hir:g1-positive-normal-shared-tail}
\lean{Hironaka.HironakaG1NormalUnequalDepth.original_family_depth_canonical_iff,Hironaka.HironakaG1NormalUnequalDepth.Lift.original_padded_bound_comparison,Hironaka.HironakaG1NormalUnequalDepth.Lift.original_adjacent_bound_comparison,Hironaka.HironakaG1NormalUnequalDepth.original_positive_shared_tail_adjacent_bound}
\leanok
Enlarging the bound of the original invariant preserves and
reflects canonicality on the same actual history, given the
tighter domain. For a supported native lift, the ambient value
at an enlarged bound is obtained by the constant-depth ambient
dimension map followed by recursive depth padding. The
canonicality equivalence holds in both directions. In particular,
on the positive shared tail the ambient section cut at bound
$k$ and its literal native lift at bound $k-1$ satisfy this
equivalence and the corresponding value formula on cosupport.
\end{theorem}
\begin{proof}
Induct over the actual tower. At each canonical prefix the
strict depth-padding map preserves and reflects the maximum
locus and hence the next canonical step. Include the empty
centre alternative at each step. Combine this with the
same-bound comparison of the supported whole records. For
the shared tail, the exact section mask and its member cut
give the tight ambient domain. Apply the adjacent-bound
comparison to that cut and its native pullback.
\end{proof}

\begin{theorem}[Comparing actual selector graphs in the same source chart]
\label{hir:g1-actual-selector-graph-comparison}
\uses{hir:g1-fixed-normal-chart-powers}
\lean{Hironaka.HironakaG1NormalSelectorGraph.graphShiftEquiv,Hironaka.HironakaG1NormalSelectorGraph.graphShift_preimage_selectorCarrier,Hironaka.HironakaG1NormalSelectorGraph.pullback_power_selectorMarked,Hironaka.HironakaG1NormalSelectorGraph.markedPullback_power_selectorSection,Hironaka.HironakaG1NormalSelectorGraph.restrictEquiv_pullback_commonPower_selectorMarked,Hironaka.HironakaG1NormalSelectorGraph.locWPEquiv_pullback_commonPower_selectorMarked}
\leanok
Two holomorphic selectors on the same source open define
actual graphs related by a triangular biholomorphism fixing
the source coordinates. It carries the whole powered padding
of a fixed source record and pulls back every ordered extra
boundary slot. Each graph section recovers the original
source power with those boundary traces. If two source powers
are restricted equivalent at the common mark and carrier,
the two padded powers on their respective graphs are
restricted equivalent and locally weakly equivalent after
transport by this same biholomorphism.
\end{theorem}
\begin{proof}
Add the difference of the selectors to the last coordinate.
The inverse subtracts that difference, and the graph equations
are identified with their true degrees. Pullback commutes
with powers and respects the source projection. The graph
sections are retractions onto the source and recover all
fields of their whole powered records. Comparing the two
source powers along the proper carrier gives the claimed
restriction inclusions after graph transport. The corresponding
local test and cosupport agreement follows from restricted
equivalence of the coherent records on their closed carrier.
\end{proof}

\begin{theorem}[Exact masked companions at the actual common order head]
\label{hir:g1-actual-normal-order-chronology}
\uses{hir:g1-actual-normal-order-companion,hir:regime-run}
\lean{Hironaka.HironakaG1NormalOrderChronology.exists_original_native_order_chronology}
\leanok
An ordinary padded test gives an actual test of the original
native record of equal length. On the kept ambient carrier
and this native test the order clocks agree at every prefix.
At every split, including the selected order-run head, the
actual tail projections commute with the head inclusion and
identify membership in the two head masks. At positive
order the whole companion computed on the kept masked
ambient record pulls back to the exact masked native
companion, with the retained ordered boundary and true branch mark.
\end{theorem}
\begin{proof}
Transport the literal supported lift to the original native
source. Actual native order and the tail projection square
identify each prefix order and its regime. The corresponding
clocks identify the physical splits. Order comparisons then
identify the transported masks and the companion pullback
commutes with restriction to them. This gives equality of
the entire masked records on the same retained geometry.
\end{proof}

\begin{theorem}[Mark-one member records and two coordinate trace orders]
\label{hir:g1-positive-normal-mark-one-trace-orders}
\uses{hir:g1-positive-normal-aligned-native-tail,hir:step-one-value-masked-basic}
\lean{Hironaka.HironakaG1PositiveNormalNonspecialExample.positive_mark_one_single_member_lower_eq,Hironaka.HironakaG1PositiveNormalNonspecialExample.positive_mark_one_single_member_nmOrd,Hironaka.HironakaG1PositiveNormalNonspecialExample.positive_mark_one_single_member_fst,Hironaka.HironakaG1PositiveNormalNonspecialExample.positive_normal_mark_one_single_member_fst,Hironaka.HironakaG1PositiveNormalNonspecialExample.positive_normal_mark_one_native_fst,Hironaka.HironakaG1PositiveNormalNonspecialExample.actual_coordinate_pair_selected_trace_order,Hironaka.HironakaG1PositiveNormalNonspecialExample.padded_coordinate_triple_selected_trace_order,Hironaka.HironakaG1PositiveNormalNonspecialExample.not_isSpecial_of_mark_one_single_member_trace_order,Hironaka.HironakaG1PositiveNormalNonspecialExample.coordinate_pair_and_triple_nonspecial}
\leanok
At the empty history and mark one, selecting the sole old
member through the anchor restricts the carrier to its
intersection with that member. The original ideal is retained
and the selected member is deleted from the boundary. The
first entries of the lower invariant compute the trace orders
on the selected padded carrier and the native section
intersection. In the coordinate models, the trace order of
$(x,y^2)$ on $x=0$ is two, and that of $(x,y^2,t)$ on $x=0$
is one. A nonzero trace order in this singleton-member
setting makes the anchor nonspecial, so both coordinate
models give nonspecial anchors.
\end{theorem}
\begin{proof}
The mark-one homogenization and coefficient operations
retain the whole ideal, and the singleton member selection
identifies the whole lower record, including its ordered
boundary and mark. Open restriction and pullback along the
native inclusion give the stated first-entry formulas. In
the coordinate calculations the selected carrier annihilates
the first generator. The nonvanishing linear normal
generator gives order one in the three-coordinate model.
In the two-coordinate model the square of the remaining
linear coordinate has exactly order two.
A special anchor would make all sections of the ideal
vanish identically near the anchor along the selected
carrier. Its ideal would then satisfy every order threshold,
which is incompatible with a nonzero trace order.
\end{proof}

\begin{theorem}[The original canonicality test for a positive repair]
\label{hir:g1-positive-normal-repair-canonicality-test}
\uses{hir:g1-positive-normal-native-special-step,hir:centre-function}
\lean{Hironaka.HironakaG1PositiveNormalRepairCentres.RepairCentreComparison,Hironaka.HironakaG1PositiveNormalRepairCentres.repair_centre_criterion,Hironaka.HironakaG1PositiveNormalRepairCentres.padded_lower_isInvDomain,Hironaka.HironakaG1PositiveNormalRepairCentres.positive_normal_repair_geometry}
\leanok
The actual padded lower record is in bound $k$. The positive
repair has an admissible centre containing the anchor, and
its transformed padded cosupport avoids the exceptional
inverse image. For any such repair the surviving cosupport
is the inverse image of the original cosupport minus the
centre and projects onto that entire difference. The step
is canonical precisely when the original head value is
constant on the centre and smaller at the projection of
every surviving cosupport point. Under this condition the
original value is unchanged at each surviving point.
\end{theorem}
\begin{proof}
The selected padded lower domain follows from positive
member restriction. The native special step provides the
actual repair and exceptional avoidance. Off the centre the
transform has the same cosupport as the original record
under the blowup map. Surjectivity gives the complete
projected difference. A nonempty centre is the head
maximum locus exactly under the stated value comparisons,
since the cosupport is partitioned by the centre and that
projected difference. The original invariant's off-centre
property at a canonical step gives the value equality.
\end{proof}


\begin{theorem}[The native lower record at the section's positive run]
\label{hir:g1-positive-normal-own-native-tail}
\uses{hir:g1-positive-normal-section-run,hir:native-section-original-positive-lower-record,hir:native-section-original-positive-lower-tail}
\lean{Hironaka.HironakaG1PositiveNormalOwnTail.positive_normal_own_native_tail_records}
\leanok
At a positive section-nonspecial point of a padded test cosupport, cut
at the section's own run. The corresponding native run has the same
length. The native own mask identifies the literal open lift, and
the head and endpoint inclusions give the whole lower-record pullbacks.
\end{theorem}
\begin{proof}
Use the section chronology to identify its run with the native run.
Split the actual upper lift and restrict at the section mask. The
positive lower head and tail formulas give the asserted records.
\end{proof}

\begin{theorem}[Centres at adjacent steps of an actual native lift]
\label{hir:g1-positive-normal-actual-lift-centres}
\uses{hir:g1-positive-normal-section-run,hir:g1-proper-normal-upper-native-test}
\lean{Hironaka.HironakaG1PositiveNormalActualCentre.Lift.adjacent_centres,Hironaka.HironakaG1PositiveNormalActualCentre.positive_normal_native_restart_centre}
\leanok
At adjacent splits of a literal lift the centres on the two towers
correspond by both image and preimage under the section inclusion.
The tower projections identify the preceding anchor on both sides.
When the section run at a specified positive level is shorter and the
endpoint is section-nonspecial, the earlier section-special anchor
lies in both actual centres under this identification.
\end{theorem}
\begin{proof}
Induct through the paired steps of the lift. The preceding section
inclusion is injective and identifies the projection of the endpoint.
Apply the actual ambient centre test and pull its incidence back.
\end{proof}

\begin{theorem}[Adjacent bounds after the actual selected repair]
\label{hir:g1-positive-normal-repair-depth}
\uses{hir:g1-positive-normal-native-special-step,hir:g1-supported-native-unequal-depth}
\lean{Hironaka.HironakaG1PositiveNormalRepairDepth.isInvDomain_of_supported_nativePullback,Hironaka.HironakaG1PositiveNormalRepairDepth.RepairDepthData,Hironaka.HironakaG1PositiveNormalRepairDepth.RepairDepthData.compare_continuation,Hironaka.HironakaG1PositiveNormalRepairDepth.RepairDepthData.exists_native_continuation,Hironaka.HironakaG1PositiveNormalRepairDepth.positive_normal_repair_depth_data}
\leanok
The ambient section cut and its native pullback after the same
selected repair both satisfy the bound $k-1$. For actual paired
continuations, ambient canonicality at depth $k$ corresponds to
native canonicality at depth $k-1$. The value comparison uses the
dimension shift and depth padding. Given a canonical ambient
section continuation, the native canonical continuation exists
with the same length and its exact transformed record.
\end{theorem}
\begin{proof}
Use the whole repaired record and its strict-section support. A
vanishing ambient germ along this carrier would force the generators
of its native pullback to vanish, so the native domain reflects to
the ambient cut. Construct the continuation lift from the ambient
section test, and apply the adjacent-bound comparison.
\end{proof}

\begin{theorem}[Actual selector graph changes over source biholomorphisms]
\label{hir:g1-normal-source-chart-initialization}
\uses{hir:g1-actual-selector-graph-comparison}
\lean{Hironaka.HironakaG1NormalGraphInitialization.baseLiftEquiv,Hironaka.HironakaG1NormalGraphInitialization.pullback_baseLift_selectorMarked,Hironaka.HironakaG1NormalGraphInitialization.selectorChartEquiv,Hironaka.HironakaG1NormalGraphInitialization.selectorChartEquiv_comp,Hironaka.HironakaG1NormalGraphInitialization.locWPEquiv_selectorChartEquiv_commonPower}
\leanok
A source biholomorphism and independent selectors on its two opens
induce an actual ambient graph change. It preserves the whole padded
record under exact pullback and composes along changes of source
charts. When transported source powers are restricted-equivalent
on the proper carrier, the corresponding padded records have the
same tests and stage cosupports with the full ordered extra boundary.
\end{theorem}
\begin{proof}
Lift the base map with the normal coordinate fixed and translate by
the difference of the selectors. Pullback along the base map commutes
with the powered ideal, graph equation and ordered boundary. The
normal translations add under composition. Apply the common-power
comparison after this literal chart change.
\end{proof}

\begin{theorem}[Order companion tails at the actual common split]
\label{hir:g1-normal-actual-order-companion-tail}
\uses{hir:g1-actual-normal-order-chronology}
\lean{Hironaka.HironakaG1NormalOrderTail.IsOrderCompanionTailTest,Hironaka.HironakaG1NormalOrderTail.original_native_order_companion_tail_of_canonical}
\leanok
On the literal native tower of an upper padded test, native
order-monotone canonicality of the transported original record
makes its positive order-companion tail a test on its own mask.
At the same split the supported ambient companion tail is a test
under order-monotone canonicality of that supported record.
\end{theorem}
\begin{proof}
Transport the actual upper lift to the native source and identify
the common order run and the position of its anchor in the mask.
Apply order-monotone canonicality on each stated input record to
its own restricted companion tail.
\end{proof}

\begin{theorem}[Two distinct finite lower entries on a single normal test]
\label{hir:g1-positive-normal-joint-example}
\uses{hir:g1-positive-normal-mark-one-trace-orders,hir:native-section-original-step-one-chronology}
\lean{Hironaka.HironakaG1PositiveNormalJointExample.proper_source_generators,Hironaka.HironakaG1PositiveNormalJointExample.section_source_generators,Hironaka.HironakaG1PositiveNormalJointExample.proper_source_trace_order,Hironaka.HironakaG1PositiveNormalJointExample.section_source_trace_order,Hironaka.HironakaG1PositiveNormalJointExample.sourceRecord_domain,Hironaka.HironakaG1PositiveNormalJointExample.joint_positive_nonspecial_first_entries}
\leanok
For one actual marked source $(x,y^2)$ of mark one and its boundary
$x=0$, the padded and its native section records have nonspecial
anchors of level one over the same origin and empty history. Their
selected masked lower invariant first entries are one for the
padded record and two for the native section.
\end{theorem}
\begin{proof}
Identify the entire padded ideal with the three coordinate generators
and the section ideal with the original pair. The genuine native
inclusion intertwines the empty histories and transfers speciality.
The trace order formulas on the actual selected carriers give the
two entries on their own masks.
\end{proof}


\begin{theorem}[Persistence of positive speciality and equality of nonspecial runs]
\label{hir:g1-positive-normal-clock-equality}
\uses{hir:g1-positive-normal-section-run,hir:step-one-value-masked-basic}
\lean{Hironaka.HironakaG1PositiveNormalRestartCarrier.selected_carrier_strictAlong_of_level_eq,Hironaka.HironakaG1PositiveNormalRestartCarrier.positive_special_step_of_level_eq,Hironaka.HironakaG1PositiveNormalRestartCarrier.positive_normal_section_run_eq_of_nonspecial}
\leanok
When the original level is unchanged along an upper test, each old
member incidence agrees with its strict transform, and the selected
lower carrier is exactly the strict transform of the old cut carrier.
At a blowup this preserves positive speciality at all points on the
transformed upper carrier which retain the previous original level.
On the padded cosupport at a positive section-nonspecial endpoint of
that level, the section and padded run lengths agree.
\end{theorem}
\begin{proof}
Each strict-transform incidence implies the old incidence. Equality of
the counts makes the two finite sets of incident slots equal. Strict
transform commutes with the intersection with these selected members.
Near the preceding special point the old cut lies in the cosupport.
Away from the exceptional inverse image the transformed cut still lies
in the transformed cosupport, which is closed. This contains the whole
strict transform over that open neighbourhood. At a hypothetical
strict cutoff, the adjacent step therefore preserves speciality, though
the start of the section run is nonspecial and has the same level.
\end{proof}


\begin{theorem}[The native test on its own positive lower mask]
\label{hir:g1-positive-normal-own-tail-test}
\uses{hir:g1-positive-normal-clock-equality,hir:g1-positive-normal-own-native-tail,hir:g1-positive-normal-shared-tail}
\lean{Hironaka.HironakaG1PositiveNormalOwnTailTest.positive_normal_own_native_tail_test}
\leanok
At a section-nonspecial endpoint of positive level on an original
padded Step~1 canonical history, the native lower record on its own
mask is in bound $k-1$. The actual native tail restricted to this
mask tests that record, and its marked endpoint belongs to the
transformed lower cosupport.
\end{theorem}
\begin{proof}
Use the common positive cut and the ambient section lower test. Its
actual restricted native lift identifies the complete lower record
with the native own record via the window isomorphism. Transport the
test, the domain and the endpoint cosupport along that isomorphism.
\end{proof}

\begin{theorem}[Tight adjacent depths at the unrepaired positive cut]
\label{hir:g1-positive-normal-section-domain}
\uses{hir:g1-positive-normal-repair-depth,hir:g1-supported-native-unequal-depth}
\lean{Hironaka.HironakaG1PositiveNormalSectionDomain.positive_normal_section_depth_data,Hironaka.HironakaG1PositiveNormalSectionDomain.positive_normal_selected_section_isInvDomain,Hironaka.HironakaG1PositiveNormalSectionDomain.positive_normal_selected_section_canonical_iff}
\leanok
At a positive section-nonspecial head of an actual upper test, the
ambient section cut and its literal native pullback both have
invariant domain $k-1$. On any actual paired lift of the selected
tail, their canonicalities agree at respective recursion bounds $k$
and $k-1$, with the exact dimension-shift and depth-padding formula
on the current cosupport.
\end{theorem}
\begin{proof}
Construct the section lower domain and the closed transverse native
carrier at the mask. The whole cut is supported there. Reflect the
tight native domain to the ambient cut and apply the supported
adjacent-bound comparison on the supplied geometric lift.
\end{proof}

\begin{theorem}[An ambient companion test from the native canonical history]
\label{hir:g1-normal-order-companion-native-can}
\uses{hir:g1-normal-actual-order-companion-tail,hir:g1-actual-normal-order-chronology}
\lean{Hironaka.HironakaG1NormalOrderCanTransport.supported_order_companion_tail_of_native_canonical}
\leanok
Given an ordinary padded upper test and its genuine normal lift,
order-monotone CAN of the literal original native history gives
TEST of the supported ambient companion on the restricted tail at
the common positive order run.
\end{theorem}
\begin{proof}
The native CAN supplies the native companion test on its own mask.
The split, order-mask transport and whole companion record comparison
identify this test on the actual open lift of the kept ambient tail.
Reflect TEST along this lift.
\end{proof}

\begin{theorem}[Graph changes in finite normal codimension]
\label{hir:g1-normal-graph-finite-block}
\uses{hir:g1-normal-source-chart-initialization}
\lean{Hironaka.HironakaG1NormalGraphFlags.blockChartMap_comp,Hironaka.HironakaG1NormalGraphFlags.blockChartEquiv,Hironaka.HironakaG1NormalGraphFlags.pullback_power_blockMarked,Hironaka.HironakaG1NormalGraphFlags.markedPullback_power_blockSection,Hironaka.HironakaG1NormalGraphFlags.blockGraph_closed_submanifold,Hironaka.HironakaG1NormalGraphFlags.restrictEquiv_pullback_commonPower_blockMarked,Hironaka.HironakaG1NormalGraphFlags.locWPEquiv_pullback_commonPower_blockMarked}
\leanok
For an arbitrary finite number of normal coordinates, a source
biholomorphism and holomorphic graph sections induce an ambient
biholomorphism compatible with composition. The whole padded powers
and genuine extra boundary list pull back along the ambient map.
Restriction to the graph recovers the source power with the extra
boundary traces. If the source powers agree in restricted equivalence,
the ambient records have the same local tests and stage cosupports.
\end{theorem}
\begin{proof}
Translate the normal coordinate vector by the difference of the
two graph sections after transporting its source. Use source
coordinates for the carrier and this translation to obtain a
closed regular graph. Restriction removes every normal generator,
while pullback along the ambient translation transports each of them
at the true mark. The resulting restricted equivalence on the closed
graph gives the full local weak class comparison.
\end{proof}

\begin{theorem}[Finite positive Step One values obstruct the standard maps]
\label{hir:g1-positive-normal-finite-obstruction}
\uses{hir:g1-positive-normal-joint-example,hir:g1-normal-flag-depth-padding}
\lean{Hironaka.HironakaG1PositiveNormalFiniteObstruction.joint_finite_stepOne_obstruction,Hironaka.HironakaG1PositiveNormalFiniteObstruction.sameIndexMap_fstEntry,Hironaka.HironakaG1PositiveNormalFiniteObstruction.depthPadMap_fstEntry,Hironaka.HironakaG1PositiveNormalFiniteObstruction.familyDimensionMap_fstEntry,Hironaka.HironakaG1PositiveNormalFiniteObstruction.joint_finite_stepOne_dimension_depth_obstruction}
\leanok
The common mark-one source $(x,y^2)$ has two actual nonspecial
Step~1 values at level one. For all positive lower recursion depths
they have finite lower values with first entries one and two. Any
lower value comparison preserving that entry fails to identify
the Step~1 values through the induced map, including dimension
shift, depth padding and their composition.
\end{theorem}
\begin{proof}
Unfold the two finite branches on their own actual masks at the
empty history. The first entries are the two computed selected
trace orders. The displayed standard lower maps preserve that
entry, so equality would equate one and two.
\end{proof}

\begin{theorem}[The native top stratum and padded maxima]
\label{hir:g1-positive-normal-branch-centres}
\uses{hir:g1-positive-normal-native-special-values,hir:g1-same-bound-native-maxima-canonicality}
\lean{Hironaka.HironakaG1PositiveNormalBranchCentres.positive_normal_branch_max_loci,Hironaka.HironakaG1PositiveNormalBranchCentres.canonical_step_mismatch_of_native_maximum}
\leanok
At a section-special and padded-nonspecial anchor of common
highest positive old-member level, the literal native maximum set
is exactly the top-valued stratum at that level. The ambient
padded anchor is maximal exactly when its finite lower value
bounds all padded lower values at that level. A witness to strict
dominance gives distinct maximum sets. With actual admissible
centre blowups and native prefix CAN, such a witness gives native
CAN and failure of padded CAN for the corresponding steps.
\end{theorem}
\begin{proof}
Identify the whole upper cosupport under the section inclusion.
Compare Step~1 pairs lexicographically on their highest first
coordinate. The native special anchor attains the largest second
coordinate. A strict padded competitor excludes its image from
the padded maximum set, so the nonempty image of the native
maximum cannot be an ambient canonical centre.
\end{proof}

\begin{theorem}[Supported original Step One values and canonicality]
\label{hir:g1-supported-step-one-values-can}
\uses{hir:g1-supported-native-unequal-depth,hir:g1-same-bound-native-maxima-canonicality}
\lean{Hironaka.HironakaActualNativeRecursiveValues.Lift.original_step_one_offset_values,Hironaka.HironakaOriginalFamilyDepthStability.original_step_one_depth_stability,Hironaka.HironakaG1NormalCentreTransport.Lift.max_loci_of_value_comparison,Hironaka.HironakaG1SupportedStepOne.Lift.original_same_bound_step_one_values,Hironaka.HironakaG1SupportedStepOne.Lift.original_same_bound_step_one_max_loci,Hironaka.HironakaG1SupportedStepOne.Lift.original_same_bound_step_one_canonical_iff,Hironaka.HironakaG1SupportedStepOne.Lift.original_step_one_with_depth}
\leanok
A supported maximal-order record and its actual native lift have
simultaneous Step~1 canonicality at a common tight bound or any
larger ambient bound. Their whole Step~1 values, including the
special and both nonspecial branches, compare by the induced map
on the lower invariant. At the common bound, the maximum loci are
exact images and preimages under the section inclusion.
\end{theorem}
\begin{proof}
Apply recursive transport at the offset bound, then cancel its
injective depth padding. Induct over the actual lift to reflect
canonicality, using the full supported cosupport and strict
monotonicity of the value map on each canonical prefix. Arbitrary
larger ambient bounds follow by Step~1 depth stability.
\end{proof}


\begin{theorem}[Step One on the literal original native history]
\label{hir:g1-literal-supported-step-one}
\uses{hir:g1-supported-step-one-values-can,hir:g1-proper-normal-upper-native-test}
\lean{Hironaka.HironakaG1ProperNormalOriginalTest.supported_normal_maximal_order,Hironaka.HironakaG1LiteralSupportedStepOne.exists_literal_supported_step_one_history}
\leanok
For a maximal-order source $Q$, an ordinary padded upper test gives
an actual native lift of the whole supported normal record. Its
transport to the original source is a test of $Q$ of the same length.
If the supported ambient record is Step~1 canonical, the transported
history is Step~1 canonical for $Q$, and all its Step~1 values on the
current cosupport give the supported values by the depth and dimension
map.
\end{theorem}
\begin{proof}
The normal powers vanish on the original section, so the supported
record has the original order along that carrier. Construct its
literal native lift from the upper test. The supported comparison
supplies canonicality and values on the raw native tower. Moving
through the native zero-section biholomorphism identifies the entire
source record with $Q$ and transports the tower and all Step~1 values.
\end{proof}


\begin{theorem}[Independent one-vertex source presentations]
\label{hir:g1-independent-one-vertex-source}
\uses{hir:g1-normal-graph-finite-block}
\lean{Hironaka.HironakaG1NormalIndependentSource.native_singleVertex_mem_iff,Hironaka.HironakaG1NormalIndependentSource.restrictEquiv_of_oneVertexPresentations,Hironaka.HironakaG1NormalIndependentSource.HasOneVertexPresentation.pullback,Hironaka.HironakaG1NormalIndependentSource.locWPEquiv_blockMarked_of_oneVertexPresentations}
\leanok
Suppose two source records have actual one-vertex presentations of
the same analytic set at every point of their common carrier, with
the same original mark and ordered boundary. The source records are
restricted-equivalent. With genuine biholomorphic source charts this
gives the local weak comparison of their graph records in arbitrary
finite normal codimension, including the ordered extra boundary.
\end{theorem}
\begin{proof}
With one vertex the compression exponent is one. The standard
basis condition identifies the sole generator ideal with the
whole analytic vanishing stalk. The separately chosen charts
therefore yield the same stalk. Use the chart transport and apply
the finite normal graph comparison to the resulting equivalence.
\end{proof}

\begin{theorem}[The next normal contact in the finite example]
\label{hir:g1-finite-example-chosen-normal-contact}
\uses{hir:g1-positive-normal-joint-example,hir:g1-positive-normal-finite-obstruction}
\lean{Hironaka.HironakaG1NormalFiniteRecodingExample.padded_selected_normal_tangent,Hironaka.HironakaG1NormalFiniteRecodingExample.padded_selected_normal_contact_cosupport,Hironaka.HironakaG1NormalFiniteRecodingExample.padded_selected_normal_contact_fst,Hironaka.HironakaG1NormalFiniteRecodingExample.padded_selected_normal_contact_at_origin}
\leanok
For the mark-one joint example $(x,y^2)$, the normal coordinate is
a genuine tangent direction of the padded selected lower record.
Along that contact, all trace cosupport thresholds agree with the
section trace. The padded lower first entry at the origin is one,
and the lower record formed using this explicit normal contact has
first entry two.
\end{theorem}
\begin{proof}
The marked ideal contains the linear normal generator on the
actual positive cut. The vertical tangent differentiates it to
one. Mark-one coefficient and homogenization keep the ideal.
The normal generator vanishes on the selected section and the
other surviving generator is the square of the remaining coordinate.
\end{proof}

\begin{theorem}[Exact finite lower values on the two section histories]
\label{hir:g1-supported-positive-exact-finite-lower-values}
\uses{hir:g1-supported-step-one-values-can,hir:g1-positive-normal-clock-equality}
\lean{Hironaka.HironakaG1PositiveNormalExactFiniteValues.original_positive_lower_finite_values,Hironaka.HironakaG1PositiveNormalExactFiniteValues.positive_normal_section_exact_finite_values}
\leanok
A supported Step~1 canonical segment at a positive nonspecial
cosupport point gives two actual canonical lower tails on their
own masks, one in the ambient section record and one in the
native record. Their finite lower values agree under the
same-bound dimension shift. For an original section history,
the cut is also at the padded positive clock.
\end{theorem}
\begin{proof}
Transport supported upper canonicality to the literal native lift.
Read the two finite branches from their own canonical lower data
at the common cut. Injectivity of finite coercion into the
extended value type gives the lower-value formula.
\end{proof}

\begin{theorem}[Supported companion Step One from original native canonicality]
\label{hir:g1-order-companion-step-one-native-can}
\uses{hir:g1-normal-order-companion-native-can,hir:g1-supported-step-one-values-can}
\lean{Hironaka.HironakaG1NormalCompanionValue.StepOneCompanionReads,Hironaka.HironakaG1NormalCompanionValue.supported_step_one_companion_of_native_canonical}
\leanok
Full canonicality on the literal original native history supplies
supported ambient Step~1 canonicality on the actual restricted
positive order companion tail at any larger ambient recursion
bound. The exact endpoint inclusion identifies the values on the
native companion cosupport by the depth and dimension map.
\end{theorem}
\begin{proof}
Read native companion canonicality from the full native invariant.
Transport its source and endpoint through the actual split and
open-window isomorphisms. Apply supported Step~1 canonicality
reflection and value comparison to that compound lift.
\end{proof}

\begin{theorem}[Both comparisons on the same selected repair]
\label{hir:g1-same-special-repair-comparisons}
\uses{hir:g1-positive-normal-repair-depth,hir:g1-positive-normal-repair-canonicality-test}
\lean{Hironaka.HironakaG1NormalSpecialStepValues.positive_normal_special_repair_values}
\leanok
One actual selected repair carries the padded maximum criterion
and both repaired section domains at the tight adjacent bound.
The surviving native cosupport maps, through its inclusion and
the blowup projection, onto the entire pre-repair padded
cosupport outside that same selected centre.
\end{theorem}
\begin{proof}
Compare the two blowup descriptions over the same centre by the
actual isomorphism. The pulled-back whole marked record transports
exceptional avoidance. Apply the cosupport projection identity
and the repaired section depth comparison on that same lift.
\end{proof}

\begin{theorem}[Centre transport from a comparison of realized branches]
\label{hir:g1-normal-centre-branch-comparison}
\uses{hir:g1-positive-normal-branch-centres,hir:g1-supported-step-one-values-can}
\lean{Hironaka.HironakaG1NormalCentreComparison.LowerBranchComparison.le_iff,Hironaka.HironakaG1NormalCentreComparison.LowerBranchComparison.max_loci,Hironaka.HironakaG1NormalCentreComparison.canonical_step_iff_of_max_loci_image,Hironaka.HironakaG1NormalCentreComparison.Lift.canonical_iff_of_prefix_max_loci,Hironaka.HironakaG1NormalCentreComparison.positive_normal_max_loci_of_lower_branch_comparison,Hironaka.HironakaG1NormalCentreComparison.positive_normal_canonical_iff_of_prefix_branch_comparisons}
\leanok
A strictly increasing comparison of finite lower values separately
at every shared level suffices if realized native top points have
a common padded value and every realized finite value is below
that value. With these hypotheses the actual maximum loci are
exact images and preimages. When the comparison holds at each
preceding split, canonicality agrees on the whole lift, including
empty steps.
\end{theorem}
\begin{proof}
Compare different levels by the first coordinate. At each level
split into finite and top branches and use the stated comparisons.
The whole cosupport projects bijectively through the section
inclusion. Transport the maximum and apply the resulting centre
identity inductively along the actual paired steps.
\end{proof}


\begin{theorem}[The normal order of an actual initial positive cut]
\label{hir:g1-original-initial-positive-cut-order}
\uses{hir:g1-positive-normal-clock-equality,hir:g1-finite-example-chosen-normal-contact}
\lean{Hironaka.HironakaG1ProperCarrierNormalContact.iterDerivIdealAlong_normal_power_mem,Hironaka.HironakaG1ProperCarrierNormalContact.coeffMarkedAlong_normal_power_mem,Hironaka.HironakaG1ProperCarrierNormalContact.normal_power_maximalOrder,Hironaka.HironakaG1ProperNormalInitialRegime.nonmonomial_eq_restrictOn_of_boundary_zero,Hironaka.HironakaG1InitialSelectedNormalOrder.initial_selected_normal_data,Hironaka.HironakaG1InitialSelectedNormalOrder.initial_selected_normal_test_order,Hironaka.HironakaG1InitialSelectedNormalOrder.original_initial_positive_lower_test_order}
\leanok
Selecting the old members through a positive point of the initial
padded record leaves a proper cylinder and preserves the normal
power at the new coefficient mark. The unmasked cut has its tight
invariant domain and maximal order, with a normal tangent direction
and zero exponents of all remaining old divisors. On every actual
lower test over its mask the current nonmonomial order on the
cosupport is the true mark of the initial cut.
\end{theorem}
\begin{proof}
Original divisors are pullbacks from the base, so the selected
intersection is a cylinder and the vertical derivative preserves it.
The zero-th coefficient summand supplies the whole normal power.
Its derivatives give nonvanishing and bound the relative order.
The vertical field also excludes divisibility by every retained
old divisor. The initial nonmonomial part is thus the carrier
saturation, and its same-mark companion is restricted-equivalent
to the initial cut. The companion transform estimates retain the
order at its mark along each following actual test.
\end{proof}

\begin{theorem}[The special branch at the first normal contact]
\label{hir:g1-initial-positive-special-recoding}
\uses{hir:g1-original-initial-positive-cut-order,hir:g1-positive-normal-native-special-values}
\lean{Hironaka.HironakaG1ProperCarrierNormalContact.normal_contact_carrier,Hironaka.HironakaG1InitialPositiveSpecialRecoding.codimOnePartAlong_iff_contact_saturated,Hironaka.HironakaG1InitialPositiveSpecialRecoding.initial_selected_special_iff_section_special,Hironaka.HironakaG1InitialPositiveSpecialRecoding.original_initial_masked_positive_special_iff}
\leanok
At a positive cosupport anchor of the initial padded record, the
selected record has level zero at the corresponding point of its own
mask. It is special there exactly when the original section record is
special at the positive anchor.
\end{theorem}
\begin{proof}
The codimension-one part of the cosupport consists of the points where
the cosupport fills the germ of a genuine tangent contact. For the
selected cylinder the normal contact is the section of the selected
carrier, so this condition is exactly the section's positive member
intersection condition. Every old member through the anchor has been
removed. Open restriction transports the level and speciality to the
actual mask.
\end{proof}

\begin{theorem}[The whole order run and companion of the initial cut]
\label{hir:g1-initial-selected-order-phase}
\uses{hir:g1-original-initial-positive-cut-order}
\lean{Hironaka.HironakaG1InitialSelectedOrderPhase.original_initial_positive_lower_order_phase,Hironaka.HironakaG1InitialSelectedOrderPhase.original_initial_positive_lower_companion_eq,Hironaka.HironakaG1InitialSelectedOrderPhase.original_initial_positive_lower_companion_recursion}
\leanok
For the initial positive selected record on its mask, the nonmonomial
order mask at the true mark is the whole base. Every actual test ending
in cosupport has its full length in that order regime. The same-mark
companion retains the selected carrier, ordered boundary and mark. Its
ideal is the carrier restriction of the selected ideal. The companion
has maximal order, equivalent canonicality, and the same full
recursive value on canonical histories at cosupport points.
\end{theorem}
\begin{proof}
The order at the true mark holds along every test, including at points
projected from later cosupport. Induct on the length of the test. The
nonmonomial ideal is the carrier restriction because every divisor
exponent is zero. Apply the invariant comparison under restricted
equivalence to the companion.
\end{proof}

\begin{theorem}[Full original values on the literal supported history]
\label{hir:g1-literal-supported-general-history}
\uses{hir:g1-literal-supported-step-one,hir:g1-supported-native-unequal-depth}
\lean{Hironaka.HironakaG1LiteralSupportedGeneralValue.exists_literal_supported_general_history}
\leanok
A padded test produces a literal test of the original input $Q$ and
identifies its transformed whole record with the pullback of the current
supported normal record. Full canonicality is equivalent for these
two histories. On canonical histories the dimension and depth map
relates their entire original invariant, with exact images and
preimages of the maximum loci.
\end{theorem}
\begin{proof}
Apply the supported comparison to the natural native lift and
transport the whole tower through the zero-section biholomorphism.
The ordered boundary and current carrier are part of the transported
record. Strict monotonicity of the combined dimension and depth map
then identifies the maximum loci.
\end{proof}

\begin{theorem}[The order of the nested normal comparison]
\label{hir:g1-normal-nested-order-map}
\uses{hir:g1-normal-centre-branch-comparison}
\lean{Hironaka.HironakaG1NormalNestedOrderMap.nestedStepEmbedding,Hironaka.HironakaG1NormalNestedOrderMap.lowerBranchComparison_of_nestedLower,Hironaka.HironakaG1NormalNestedOrderMap.max_loci_of_nestedLower,Hironaka.HironakaG1NormalNestedOrderMap.positive_normal_max_loci_of_nestedLower}
\leanok
Fix the new nonmonomial order at each upper level and put the native
lower read inside the finite level-zero companion letter. A native
top read gives an inner top inside that finite letter. The resulting
map is a strict order embedding. When the actual lower reads on the
padded and native histories satisfy this formula, their maximum loci
are exact images and preimages under the section inclusion.
\end{theorem}
\begin{proof}
Different upper levels are compared first. At a fixed level, the
nonmonomial order and companion level are fixed, so comparison
reduces to the strictly increasing lower map extended to top. This
supplies all clauses of the realized-branch comparison.
\end{proof}

\begin{theorem}[Equal-weight compression of two generators]
\label{hir:g1-two-vertex-compression}
\uses{hir:g1-independent-one-vertex-source,hir:actual-weighted-compression-radical}
\lean{Hironaka.HironakaG1NormalMultiVertexSource.compression_two_weight_two_mul_span,Hironaka.HironakaG1NormalMultiVertexSource.compression_two_weight_two_mul_span_eq,Hironaka.HironakaG1NormalMultiVertexSource.exists_same_span_different_two_weight_two_compression}
\leanok
Two generators of a common ideal $I$ with weights $(2,2)$ give a
compressed ideal $J$ with $JI=I^3$. Equality of the original spans
alone can give different compressed ideals, even with these common
weights. In two complex formal variables the lists $(x,y)$ and
$(x+y,y)$ give such an example with the same radical.
\end{theorem}
\begin{proof}
Every cubic monomial in two generators has a square factor. For
the example, the mixed coefficient of $(x+y)^2$ is two, whereas
every element of $(x^2,y^2)$ has mixed coefficient zero.
\end{proof}

\begin{theorem}[The canonical contact in the finite example]
\label{hir:g1-finite-example-canonical-normal-contact}
\uses{hir:g1-finite-example-chosen-normal-contact}
\lean{Hironaka.HironakaG1CanonicalNormalSelectedExample.canonical_normal_selected_next_first_entry,Hironaka.HironakaG1CanonicalNormalSelectedExample.canonical_normal_selected_finite_stepOne}
\leanok
For the actual initial positive selected cut of $(x,y^2)$, the
canonical Step~1 value at the origin has level zero and a finite
lower value. The entire lower value equals the invariant of the
explicit normal contact record, and its first entry is two.
\end{theorem}
\begin{proof}
The padded linear normal generator provides the genuine tangent.
The local tangent-change theorem equates the chosen and prescribed
lower values on their actual masks. The explicit contact retains
the square in the remaining coordinate, giving first entry two.
\end{proof}

\begin{theorem}[Prescribed contact on the actual canonical tail]
\label{hir:g1-prescribed-zero-contact-value}
\uses{hir:g1-initial-selected-order-phase,hir:g1-finite-example-canonical-normal-contact}
\lean{Hironaka.HironakaG1NormalContactChoiceTransfer.zero_step_one_prescribed_contact,Hironaka.HironakaG1NormalContactChoiceTransfer.initial_positive_companion_normal_contact,Hironaka.HironakaG1InitialCompanionNormalTangent.initial_positive_companion_normal_tangent}
\leanok
For a nonspecial level-zero endpoint of a canonical Step~1 history,
the prescribed global tangent uses the same restricted tail,
has the lower canonicality and endpoint cosupport, and gives the exact
finite lower value. For an initial positive selected cut, its literal
order companion has invariant domain, maximal order, and a genuine
global normal tangent at its own mark.
\end{theorem}
\begin{proof}
Transport lower canonicality through the weak equivalence of the
coefficient records and apply tangent-change invariance on their
actual masks. For the selected companion, the normal direction on the
raw cut survives restriction to the carrier and then open pullback.
The identified companion ideal and full order mask supply the domain
and maximal order.
\end{proof}

\begin{theorem}[Canonicality on the bare contact of an initial zero phase]
\label{hir:g1-initial-zero-bare-contact-canonical}
\uses{hir:g1-prescribed-zero-contact-value}
\lean{Hironaka.HironakaG1ProperNormalZeroPhase.prescribed_contact_zero_phase,Hironaka.HironakaG1PrescribedContactBarRecoding.prescribed_contact_bar_data,Hironaka.HironakaG1PrescribedContactBarRecoding.prescribed_contact_bar_canonical_iff,Hironaka.HironakaG1InitialBareContactCanonical.initial_zero_stepOne_bare_contact_canonical}
\leanok
An actual nonspecial zero-phase suffix of an original canonical
maximal-order Step~1 history is canonical on the homogenized
contact of a prescribed global tangent. If the zero phase starts
at the source and the bare coefficient contact has invariant
domain, its original invariant is canonical on the same open
mask and restricted tower. The endpoint remains in its cosupport.
\end{theorem}
\begin{proof}
Use the genuine tangent on the actual canonical tail. The
homogenized and bare coefficient contacts have the same carrier
and local weak class. Their cosupport gaps give equality of
normalized exponents on the mask, so the original lower canonicality
and the endpoint cosupport transfer along the literal tower.
\end{proof}

\begin{theorem}[The original general letter in the finite example]
\label{hir:g1-finite-example-original-general-letter}
\uses{hir:g1-finite-example-canonical-normal-contact,hir:g1-initial-selected-order-phase}
\lean{Hironaka.HironakaG1CanonicalNormalGeneralExample.canonical_normal_selected_general_value,Hironaka.HironakaG1CanonicalNormalGeneralExample.canonical_normal_original_positive_stepOne_value}
\leanok
The original padded selected general value at the origin of the
$(x,y^2)$ example is the depth-two initial embedding of the depth-one
normal contact value of first entry two. The original padded Step~1
value is the finite pair at level one with precisely that embedded
general letter as its lower coordinate.
\end{theorem}
\begin{proof}
The whole order run supplies the actual same-mark companion. Its
Step~1 read compares to the selected read through restricted
equivalence and open restriction. Insert the exact contact value into
the general successor and then into the original positive lower
branch at the empty history.
\end{proof}

\begin{theorem}[Weak test comparison for two quadratic generators]
\label{hir:g1-two-vertex-quadratic-weak-class}
\uses{hir:g1-two-vertex-compression}
\lean{Hironaka.HironakaG1TwoVertexWeakClass.cosupport_eq_of_quadratic_reduction,Hironaka.HironakaG1TwoVertexWeakClass.admissibleCentre_iff_of_common_quadratic_reduction,Hironaka.HironakaG1TwoVertexWeakClass.locWPEquiv_of_ambientQuadraticReduction,Hironaka.HironakaG1TwoVertexWeakClass.quadratic_reduction_of_two_functions,Hironaka.HironakaG1TwoVertexWeakClass.locWPEquiv_two_weight_two_functions}
\leanok
On a common closed submanifold, a quadratic reduction
$J\subseteq I^2$ and $JI=I^3$ gives the same cosupport at marks $2d$
and $d$, and the same admissible centres with the same ordered
boundary. For whole coherent ideals with full ambient carrier,
the reduction gives the local weak test class, including all stage
cosupports. In particular two holomorphic pairs with the same whole
source span and weights $(2,2)$ give that class at compression
mark four.
\end{theorem}
\begin{proof}
If $I$ has finite order $a<d$, then order at least $2d$ for $J$
would make the order of $I^3$ strictly greater than $3a$,
contrary to the sharp power threshold. The other inclusion follows
from $J\subseteq I^2$. For the ambient carriers, controlled products,
powers and restrictions preserve the reduction. Two holomorphic
generators satisfy the required identities at every stalk by the
cubic monomial formula.
\end{proof}

\begin{theorem}[The inner top on the actual initial companion]
\label{hir:g1-initial-positive-companion-special-value}
\uses{hir:g1-initial-positive-special-recoding,hir:g1-initial-selected-order-phase}
\lean{Hironaka.HironakaG1InitialCompanionSpecial.initial_positive_companion_special_iff,Hironaka.HironakaG1InitialCompanionSpecial.initial_positive_companion_special_value}
\leanok
The actual initial order companion on its own mask has level zero at
the positive anchor. The original section's special branch is
equivalent to the companion's special branch there. At such an
anchor the companion reads the pair $(0,\top)$ at every depth.
\end{theorem}
\begin{proof}
Compare the selected record to its carrier-saturated companion at the
same mark using restricted equivalence. The whole ordered boundary
and the common carrier identify the level and speciality. Restrict to
the literal order mask and read the special branch of the original
Step~1 evaluator.
\end{proof}

\begin{theorem}[Chronological comparison through a finite normal flag]
\label{hir:g1-normal-nested-flag}
\uses{hir:g1-normal-nested-order-map}
\lean{Hironaka.HironakaG1NormalNestedFlag.NestedOrderFlag.embedding_append,Hironaka.HironakaG1NormalNestedFlag.NestedOrderFlag.depth_eq,Hironaka.HironakaG1NormalNestedFlag.literalNestedStepComparison,Hironaka.HironakaG1NormalNestedFlag.ActualNestedFlag.compare,Hironaka.HironakaG1NormalNestedFlag.NestedFlagComparison.comparisons,Hironaka.HironakaG1NormalNestedFlag.NestedFlagComparison.max_loci}
\leanok
For a finite flag of literal adjacent histories, suppose the actual
lower read of each adjacent pair satisfies its nested formula on the
endpoint and every preceding split. When the intermediate histories
agree, the endpoint reads compare by the composition of the order
embeddings. This gives exact image and preimage of the maximum
loci and canonicality equivalence at every chronological split.
\end{theorem}
\begin{proof}
Use the adjacent theorem on each literal pair to transport levels,
cosupports and canonicality. Compose the section inclusions and
nested embeddings along the flag. The depth increases by one per
inclusion. Strict monotonicity and the cosupport identities then give
both maximum-locus formulas.
\end{proof}

\begin{theorem}[The initial selected normal contact and the section test class]
\label{hir:g1-initial-selected-normal-section-test-class}
\uses{hir:g1-original-initial-positive-cut-order,hir:g1-prescribed-zero-contact-value}
\lean{Hironaka.HironakaG1InitialNormalSelectedSectionClass.initial_normal_selected_section_wEquiv}
\leanok
On a common initial open containing the positive anchor, the
prescribed normal homogenized coefficient contact and the actual
supported section-selected record have equal carriers and ordered
boundaries, equivalent tests and the same transformed cosupports.
Both coefficient marks are retained. If both records have invariant
domain, their crossed whole powers give this comparison at a common
mark.
\end{theorem}
\begin{proof}
The nested-carrier member restriction compares both selected
records before contact. On the padded cut, the prescribed normal
coordinate is tangent and cuts the supported carrier. The
homogenized coefficient test theorem compares to that actual contact.
The two positive marked powers preserve weak test equivalence and
give the common product mark.
\end{proof}

\begin{theorem}[The level-zero contact on every selected lower history]
\label{hir:g1-initial-selected-zero-history-contact}
\uses{hir:g1-prescribed-zero-contact-value,hir:g1-initial-selected-order-phase}
\lean{Hironaka.HironakaG1InitialCompanionZeroHistory.initial_positive_companion_zero_history,Hironaka.HironakaG1InitialCompanionZeroHistory.initial_positive_companion_zero_contact_value}
\leanok
The initial positive companion has level zero along every finite
history. A tested continuation of the lower record supplies the
actual companion endpoint and its cosupport. If its restricted
companion history is Step~1 canonical and that endpoint is
nonspecial, the prescribed normal contact gives the exact finite
value on the same restricted tail.
\end{theorem}
\begin{proof}
All original boundary slots on the mask are empty and remain empty
under strict transform. Carrier restriction and the companion's
whole order mask identify the tested endpoint cosupport. Apply the
prescribed contact comparison using the proved global normal tangent.
\end{proof}

\begin{theorem}[Quadratic test comparison on a proper closed carrier]
\label{hir:g1-two-vertex-proper-quadratic-weak-class}
\uses{hir:g1-two-vertex-quadratic-weak-class}
\lean{Hironaka.HironakaG1CarrierQuadraticReduction.isLocStable_carrierQuadraticReduction,Hironaka.HironakaG1CarrierQuadraticReduction.locWPEquiv_of_carrierQuadraticReduction,Hironaka.HironakaG1CarrierQuadraticReduction.locWPEquiv_two_weight_two_functions_on_carrier}
\leanok
The quadratic reduction with its inclusions on a common closed
carrier preserves the whole local weak test class at marks $2d$
and $d$. In particular, given the invariant domain of the mark-two
uncompressed record, two whole holomorphic generator pairs with
the same span give equivalent mark-four compressions on that carrier,
with the ordered boundary unchanged.
\end{theorem}
\begin{proof}
Use the sharp power threshold on the carrier to compare cosupports.
Controlled transform preserves both carrier-relative inclusions
because products and powers transform compatibly up to restriction.
Line extension and open restriction give the remaining local tests.
For actual generator pairs the cubic identity supplies the
inclusions.
\end{proof}

\begin{theorem}[Reduction for any finite family of equal weights]
\label{hir:g1-uniform-weight-power-reduction}
\uses{hir:g1-two-vertex-compression}
\lean{Hironaka.HironakaG1UniformVertexReduction.purePowers_mul_span_pow,Hironaka.HironakaG1UniformVertexReduction.uniformCompression_mul_span_pow,Hironaka.HironakaG1UniformVertexReduction.cosupport_eq_of_power_reduction}
\leanok
Let $I$ be the span of a finite family and $J$ the span of its
pure $e$-th powers. If $e+r$ exceeds the number of generators
times $e-1$, then $JI^r=I^{e+r}$. This applies to the
actual compression of a nonempty family of equal positive weights.
On a closed carrier, $J\subseteq I^e$ and this reduction identity
identify the cosupport of $J$ at mark $ed$ with that of $I$ at mark
$d$ for every positive $e$.
\end{theorem}
\begin{proof}
A homogeneous generator of $I^{e+r}$ is a sum of monomials of
degree $e+r$. Each has a pure power factor of exponent at least
$e$, by the degree bound. The product inclusion in the other
direction follows from $J\subseteq I^e$. If the local order of
$I$ is below $d$, the reduction identity and the sharp power
threshold rule out order at least $ed$ for $J$.
\end{proof}

\begin{theorem}[Uniform native weights on an actual carrier]
\label{hir:g1-independent-uniform-native-class}
\uses{hir:g1-uniform-weight-power-reduction,hir:g1-two-vertex-proper-quadratic-weak-class}
\lean{Hironaka.HironakaG1CarrierPowerReduction.locWPEquiv_of_restricted_power_reduction,Hironaka.HironakaG1CarrierPowerReduction.uniformNative_stalk_powerReduction,Hironaka.HironakaG1CarrierPowerReduction.restricted_powerReduction_of_uniformNative,Hironaka.HironakaG1CarrierPowerReduction.locWPEquiv_of_independent_uniformNative_presentations}
\leanok
Two actual families of native point presentations on the whole
carrier are locally weak-test equivalent when they have the same
analytic source, carrier, ordered boundary and common positive
vertex weight, and the marked source vanishing ideal lies in the
invariant domain. The independently chosen vertices, charts and
selectors can give different compression marks.
\end{theorem}
\begin{proof}
Represent each actual compressed stalk and its source in the same
adapted chart. Products and powers commute with that representation.
The uniform-weight reduction supplies restricted inclusions at each
carrier point. The controlled transform, line product and open
restriction preserve these inclusions at their respective marks.
Compare both records to the uncompressed vanishing ideal.
\end{proof}

\begin{theorem}[The actual selected contact class and canonical histories]
\label{hir:g1-initial-selected-normal-section-local-class}
\uses{hir:g1-initial-selected-normal-section-test-class}
\lean{Hironaka.HironakaG1InitialNormalSelectedLocalClass.initial_normal_selected_section_locWP,Hironaka.HironakaG1InitialNormalSelectedLocalClass.initial_normal_selected_mask_exponents_zero,Hironaka.HironakaG1InitialNormalSelectedLocalClass.initial_normal_selected_section_mark_obstruction,Hironaka.HironakaG1InitialNormalSelectedLocalClass.initial_normal_selected_mask_original_can_iff}
\leanok
The actual prescribed initial normal contact and the literal
supported section cut have the same local weak product-test class
on any common open. On the true positive mask their corresponding
ordered divisors are empty and their exponents vanish, including
after marked powers. The contact mark is the iterated factorial of
the section mark, which rules out restricted equivalence when the
section mark is at least three. If both records have the same
invariant-family domain, canonicality agrees along their common
actual tower.
\end{theorem}
\begin{proof}
Separate-carrier member restriction preserves the local test class.
Compare the chosen normal contact using the actual reference contact
on the selected carrier. The empty slots give zero normalized
exponents, and the unequal-mark original canonical comparison
applies directly.
\end{proof}

\begin{theorem}[The full native value in the joint finite example]
\label{hir:g1-finite-example-full-native-value}
\uses{hir:g1-finite-example-original-general-letter,hir:g1-literal-supported-general-history}
\lean{Hironaka.HironakaG1CanonicalNormalJointNativeExample.initial_orderTwoSpecial_value,Hironaka.HironakaG1CanonicalNormalJointNativeExample.joint_native_selected_matches_normal_contact}
\leanok
On the same literal empty lift of the joint $(x,y^2)$ example,
the native selected depth-one read has order two and the finite
level-zero special letter. The prescribed padded contact read is
its same-depth dimension shift, and the actual padded selected
depth-two read is the initial embedding of that image.
\end{theorem}
\begin{proof}
The isolated cosupport on the smooth curve is the codimension-one
special part of the actual mark-two companion. Compute its original
general read on each actual order mask and apply the verified
finite padded recursion.
\end{proof}

\begin{theorem}[The selected normal contact on a later upper chart]
\label{hir:g1-later-positive-selected-chart-contact}
\uses{hir:g1-original-initial-positive-cut-order}
\lean{Hironaka.HironakaG1LaterSelectedNormalTangent.later_normal_bar_contact,Hironaka.HironakaG1LaterSelectedNormalTangent.selected_tangent_of_current_contact,Hironaka.HironakaG1NormalSimultaneousFields.later_positive_member_normal_contact_of_simultaneous_chart}
\leanok
At a tested upper stage the actual transformed bar has a local
normal contact cutting the strict normal section. On a simultaneous
chart straightening the current carrier and the actual selected
members, a free coordinate on which this contact has nonzero
derivative produces a tangent of the selected coefficient ideal.
Its zero trace is the strict normal section on the selected carrier.
\end{theorem}
\begin{proof}
The controlled normal contact produces the bar derivative. Its full
coefficient power belongs to the selected ideal. Divide the free
coordinate field by its nonvanishing derivative of the contact and
differentiate that power along the selected carrier. The resulting
coordinate realizes exactly the prescribed trace.
\end{proof}

\begin{theorem}[A source reduction on the finite normal graph]
\label{hir:g1-normal-graph-power-reduction}
\uses{hir:g1-independent-uniform-native-class,hir:g1-normal-graph-finite-block}
\lean{Hironaka.HironakaG1GraphPowerReduction.locWPEquiv_blockMarked_of_source_power_reduction,Hironaka.HironakaG1GraphPowerReduction.locWPEquiv_blockMarked_of_uniformNative}
\leanok
A carrier-relative source reduction transports to a common actual
finite normal graph whenever the graph records have invariant
domain with their given ordered extra boundary. In particular an
actual uniform-weight native family has the same local weak test
class as the graph of its original uncompressed analytic source.
\end{theorem}
\begin{proof}
The extra normal powers vanish along the graph. Pulling back the
source reduction along the block projection preserves its powers,
products and restricted inclusions. Transfer those inclusions to
the full graph records and apply the local power-reduction theorem.
\end{proof}

\begin{theorem}[The original successor at the actual initial positive cut]
\label{hir:g1-original-initial-positive-successor-read}
\uses{hir:g1-initial-selected-order-phase,hir:g1-initial-selected-zero-history-contact}
\lean{Hironaka.HironakaG1InitialPositiveGeneralRead.original_initial_positive_lower_successor_read,Hironaka.HironakaG1InitialPositiveGeneralRead.original_upper_initial_positive_successor_read}
\leanok
Under padded upper Step~1 canonicality and a full positive run
starting at the initial stage, the original finite lower value is
the exact general successor read on the literal true-mark order
companion. The same restricted companion tail carries its Step~1
canonicality and endpoint cosupport. The additional general
letter raises the recursive depth by one.
\end{theorem}
\begin{proof}
Read the actual lower canonical tail from upper canonicality.
The true order is constant at the initial mark along the tested
lower tail. Its order split therefore starts at the initial record,
and the companion on its own order mask supplies the exact
successor letter of the original evaluator.
\end{proof}

\begin{theorem}[The exact nested top read on the initial literal lift]
\label{hir:g1-original-initial-positive-nested-top}
\uses{hir:g1-original-initial-positive-successor-read,hir:g1-initial-positive-companion-special-value,hir:g1-normal-nested-order-map}
\lean{Hironaka.HironakaG1InitialPositiveNestedRead.original_initial_positive_special_nested_read}
\leanok
At an initial positive section-special, padded-nonspecial anchor
on the actual literal lift, the native read is the top pair at
that positive level. The padded read is exactly its nested image.
The inserted order is the single source-wide coefficient mark, and
the inner lower map is arbitrary on this native top branch.
\end{theorem}
\begin{proof}
The section speciality and level transport to the actual native
record. The initial selected companion reads the level-zero top
pair on its actual order mask. Inserting it in the verified general
successor gives the nested image at the common positive level.
\end{proof}

\begin{theorem}[Finite weight blocks on an actual carrier]
\label{hir:g1-finite-weight-block-reductions}
\uses{hir:g1-uniform-weight-power-reduction,hir:g1-independent-uniform-native-class}
\lean{Hironaka.HironakaG1NonuniformVertexReduction.locWPEquiv_twoBlockFunctions_of_equal_spans,Hironaka.HironakaG1FiniteWeightedBlocks.sum_reductions,Hironaka.HironakaG1FiniteWeightedBlocks.finite_block_compression_reduction,Hironaka.HironakaG1FiniteWeightedBlocks.locWPEquiv_of_independent_finiteBlockAdaptedPresentations}
\leanok
A finite weighted family has a power reduction to the sum of
its normalized block ideals. Each block retains its actual
exponent at the common mark. Independent adapted presentations
with the same realized reference on their common invariant-domain
carrier give the same local weak test class at that mark.
For two blocks of actual holomorphic functions it suffices that
the entire source spans of the two matched blocks agree.
\end{theorem}
\begin{proof}
The uniform reduction applies at each block's own exponent.
Expand a power of the sum of the block ideals. At the combined
offset, some block exponent exceeds its own offset, so apply that
block's reduction equation. Induction adds the remaining blocks.
Read the resulting inclusions in the actual chart stalks and
apply the proper-carrier local weak test comparison.
\end{proof}

\begin{theorem}[Canonical histories on the companion's prescribed section]
\label{hir:g1-initial-companion-prescribed-section}
\uses{hir:g1-initial-selected-zero-history-contact,hir:g1-initial-selected-normal-section-local-class}
\lean{Hironaka.HironakaG1InitialCompanionSectionComparison.initial_companion_normal_section_contact,Hironaka.HironakaG1InitialCompanionSectionComparison.initial_companion_normal_section_can_iff,Hironaka.HironakaG1InitialCompanionSectionComparison.initial_companion_normal_section_finite_refl}
\leanok
On the literal nested masks, the prescribed normal contact of the
initial order companion and the supported section cut have the
same carrier and local weak product-test class. The actual domain
of the contact supplies the same domain for the section, and
original canonicality agrees on common tails. At a nonspecial
surviving point of the empty lower history, the actual finite
contact branch supplies both domains, supported endpoint
cosupport, and the exact original Step~1 value from the contact.
\end{theorem}
\begin{proof}
The whole initial companion uses carrier restriction at the same
mark. Compare its actual reference contact to the section cut
by the separate-carrier member theorem. The remaining ordered
divisors are empty on the true mask, so the unequal-mark
canonical comparison applies. For the empty lower history, read
the exact finite branch from the prescribed normal tangent on
the actual mask.
\end{proof}

\begin{theorem}[The original positive successor at a later split]
\label{hir:g1-original-later-positive-successor-read}
\uses{hir:g1-original-initial-positive-successor-read}
\lean{Hironaka.HironakaG1LaterPositiveSuccessorRead.original_positive_order_successor_read,Hironaka.HironakaG1LaterPositiveSuccessorRead.original_upper_positive_successor_read}
\leanok
A finite positive upper branch of the original Step~1 canonical
history has an exact successor read whenever its actual masked
lower endpoint has positive nonmonomial order. The two splits
are the original positive regime split and then the lower order
regime split. On the latter's actual mask the companion has its
domain, maximal order, canonical tail and endpoint cosupport.
The whole upper finite value is the successor letter with the
companion read as its inner coordinate.
\end{theorem}
\begin{proof}
Read the canonical lower tail from the upper history at its
positive regime split. At the witnessed positive raw order, the
general evaluator supplies the canonical companion at the exact
order-regime head and restricts its literal tail. Combine the two
branch equalities at their respective depths.
\end{proof}

\begin{theorem}[Weak comparison across actual changes of selectors]
\label{hir:g1-normal-graph-chart-reduction}
\uses{hir:g1-normal-graph-power-reduction,hir:g1-normal-graph-finite-block}
\lean{Hironaka.HironakaG1NormalGraphFlags.pullback_blockMarked,Hironaka.HironakaG1NormalGraphFlags.regularAtDim_blockGraph_of,Hironaka.HironakaG1NormalGraphFlags.blockMarked_isInvDomain,Hironaka.HironakaG1GraphPowerReduction.locWPEquiv_pullback_blockMarked_of_source_power_reduction}
\leanok
A carrier-relative power reduction of the pulled-back source
records gives a local weak comparison of the actual graph records
on any finite common normal block. The chart change may use
independent selectors. It transports the whole ordered boundary.
The graph invariant domains follow from the source domains when
that boundary is SNC and transverse to the common graph.
\end{theorem}
\begin{proof}
The chart map carries the whole marked graph record to the graph
of the transported source. Straighten each graph to the zero section
to check the carrier dimension. The exact trace of the whole ideal
on the graph is the source ideal, so its nonvanishing along the
carrier follows by reflecting through the section. Apply the
carrier-relative power comparison on the common graph.
\end{proof}

\begin{theorem}[A later positive survivor in the joint example]
\label{hir:g1-joint-example-later-positive-survivor}
\uses{hir:g1-finite-example-full-native-value,hir:g1-later-positive-selected-chart-contact}
\lean{Hironaka.HironakaG1LaterPositiveNormalExample.joint_first_point_blowup_positive_normal_contact}
\leanok
The first blowup of the common origin of $(x,y^2)$ tests the
actual padded and supported source records. The point in the
$y$-chart over the origin lies in the transformed padded
cosupport, at original level one, on the strict original
divisor and strict normal section. The actual later bar
has a normal contact there cutting out that strict section.
\end{theorem}
\begin{proof}
In the literal chart the blowup divides each pulled-back generator
by the same exceptional coordinate. The resulting generators
are the three local coordinates, which vanish together at the chart
origin. The strict divisor and section contain this surviving
point. Apply the bar contact theorem to the actual controlled test.
\end{proof}

\begin{theorem}[Independent graphs over a common marked reference]
\label{hir:g1-independent-graph-common-reduction}
\uses{hir:g1-normal-graph-chart-reduction}
\lean{Hironaka.HironakaG1GraphPowerReduction.locWPEquiv_pullback_blockMarked_of_common_reduction}
\leanok
Two actual graph records with independent selectors and compression
marks have the same local weak test class after the chart change
when both source records reduce to a common reference. One chart's
SNC and transverse ordered boundary supplies the graph geometry
on the common base.
\end{theorem}
\begin{proof}
Compare each source to the reference on its actual graph.
Pull the right comparison and the full reference domain through
the chart biholomorphism. Its carrier and boundary agree
with those of the left graph, so the comparisons compose.
\end{proof}

\begin{theorem}[Generator changes modulo the normalized reference]
\label{hir:g1-filtered-weighted-presentation-reference}
\uses{hir:g1-finite-weight-block-reductions}
\lean{Hironaka.HironakaG1WeightedPresentationTransport.blockPowerReference_eq_of_mutual_modulo,Hironaka.HironakaG1WeightedPresentationTransport.blockPowerReference_eq_of_crossBlock_shear,Hironaka.HironakaG1WeightedPresentationTransport.exists_crossWeight_shear_with_different_block_span,Hironaka.HironakaG1WeightedPresentationTransport.locWPEquiv_of_filteredBlockChanges}
\leanok
Mutual inclusion of each block modulo the other normalized reference
identifies the two normalized references. This allows adding a
multiple of a normalized power from an unchanged block, even when
the whole block span changes. Actual adapted source data with
these inclusions and the same cross mark give the local weak
test comparison on their common carrier.
\end{theorem}
\begin{proof}
Each transformed generator is a member of the corresponding
block plus the reference ideal. Its normalized power therefore
lies in the reference. Apply the same reasoning in the other
direction. In the actual stalks, the common resulting reference
supplies the finite block comparison.
\end{proof}

\begin{theorem}[Normal tangency on the actual later positive cut]
\label{hir:g1-later-positive-contact-exists}
\uses{hir:g1-later-positive-selected-chart-contact}
\lean{Hironaka.HironakaG1LaterContactFieldExistence.exists_normalized_tangent_field_of_transverse_boundary,Hironaka.HironakaG1LaterContactFieldExistence.exists_later_positive_selected_normal_contact}
\leanok
For an actual padded test and a point on the strict normal section,
there is a genuine open through that point on which the actual
selected member cut has a tangent function. Its zero locus on
the selected carrier is the strict normal section there.
\end{theorem}
\begin{proof}
The transformed padded and section records are well formed and
transverse to the same SNC boundary. Straighten each with that
boundary and compare the tangent differentials. The normal
contact has a nonzero derivative in a direction of the selected
carrier. Normalize this field on a neighborhood and apply the
coefficient power argument.
\end{proof}

\begin{theorem}[The first later finite value in the joint example]
\label{hir:g1-joint-example-later-selected-finite}
\uses{hir:g1-joint-example-later-positive-survivor,hir:g1-later-positive-contact-exists}
\lean{Hironaka.HironakaG1LaterPositiveSelectedExample.joint_first_blowup_masked_selected_normal_tangent,Hironaka.HironakaG1LaterPositiveSelectedExample.joint_first_blowup_selected_nonspecial,Hironaka.HironakaG1LaterPositiveSelectedExample.joint_first_blowup_step_one_finite}
\leanok
At the later positive survivor in the actual $y$-chart,
the selected tangent restricts to the intersection with the true
positive mask. It cuts the strict normal trace of the selected
carrier. The actual upper Step~1 read at this point has first
entry one and a finite lower entry.
\end{theorem}
\begin{proof}
The punctured normal line in the strict original divisor is
outside cosupport, so the point is nonspecial. The exact
positive mask contains it and retains the selected tangent.
The original Step~1 finite-branch theorem gives the numerical
conclusion on the same actual test.
\end{proof}

\begin{theorem}[The literal order ratio on the initial section]
\label{hir:g1-initial-companion-section-order-ratio}
\uses{hir:g1-initial-companion-prescribed-section}
\lean{Hironaka.HironakaG1InitialCompanionSectionNumerics.initial_companion_normal_section_first_entry,Hironaka.HironakaG1InitialCompanionSectionNumerics.initial_companion_normal_section_finite_first_entry}
\leanok
At the actual empty lower history, the first invariant entry of
the prescribed contact is the first supported section entry
multiplied by the ratio of the contact mark to the section mark.
When the section mark is at least three, the actual finite Step~1
read differs from the read obtained by inserting the unscaled
supported section value.
\end{theorem}
\begin{proof}
The proper-contact and supported section records share the
normalized nonmonomial order because all retained indexed slots
are empty. The contact mark is the iterated factorial of the
section mark. The ratio is then an integer and is greater than
one in the indicated range. Apply the original finite read to
the actual prescribed contact on the literal mask.
\end{proof}

\begin{theorem}[Two actual finite positive anchors]
\label{hir:g1-initial-positive-two-finite-centres}
\uses{hir:g1-original-initial-positive-successor-read}
\lean{Hironaka.HironakaG1InitialPositiveFiniteCentres.exists_initialFiniteRead,Hironaka.HironakaG1InitialPositiveFiniteCentres.InitialFiniteRead.upper_order_iff_normalized,Hironaka.HironakaG1InitialPositiveFiniteCentres.InitialFiniteRead.two_anchor_maxLocus_image}
\leanok
For two actual initial positive finite reads from a common
source and level, supplied canonical companion normalizations
at a common mark preserve strict order, weak order, and ties.
If both anchors are on the same stage, the normalized selector
has exactly the original two-point maximum set as its image.
\end{theorem}
\begin{proof}
Read the actual companion with its source-wide selected mark.
At equal levels and that fixed mark, only the recursive
companion coordinate is compared. Apply the original two-pair
order comparison to the normalized companions. The two-point
maximum readback follows in both directions, including a tie.
\end{proof}

\begin{theorem}[An actual later positive subwindow with its full boundary]
\label{hir:g1-later-positive-mask-contact}
\uses{hir:g1-later-positive-contact-exists,hir:g1-original-later-positive-successor-read}
\lean{Hironaka.HironakaG1LaterPositiveMaskContact.positive_member_domain_on_subwindow,Hironaka.HironakaG1LaterPositiveMaskContact.later_positive_mask_contact}
\leanok
At a nonspecial positive cosupport anchor of an actual upper test,
one open containing the anchor lies inside the literal positive
mask and carries the selected record, its invariant domain, and a
normal tangent cutting out the strict normal trace. The whole
record is the literal restriction of the lower mask record. Its
ordered boundary is the transformed upper list with the selected
old slots deleted, with the remaining born slots retained.
\end{theorem}
\begin{proof}
The upper tests put the cosupport anchor on the actual strict
normal section. The positive nonspecial locus gives membership
in the original mask. Intersect that mask with the tangent
window, and restrict the selected coefficient record, its domain
and the entire ordered boundary to the common open.
\end{proof}

\begin{theorem}[The full same-level finite stratum]
\label{hir:g1-initial-positive-finite-stratum}
\uses{hir:g1-initial-positive-two-finite-centres}
\lean{Hironaka.HironakaG1InitialPositiveStratumCentres.initial_finite_stratum_max_loci,Hironaka.HironakaG1InitialPositiveStratumCentres.exists_initial_finite_stratum_max_loci}
\leanok
On any same-level stratum of actual initial finite reads,
the exact maximum loci are images and preimages for the
normalized comparison, with a common normalized mark and the
pointwise canonical companion data. The anchor map can identify
different labels, and empty strata and ties are included.
\end{theorem}
\begin{proof}
Construct the finite reads from the actual padded canonical history
and full initial positive run. Apply pairwise order preservation
at each pair of anchors. Since the realized stratum is exactly the
image of the anchor map, the universal maximum predicates
transport in both directions.
\end{proof}

\begin{theorem}[Independent native presentations in their own charts]
\label{hir:g1-independent-weighted-chart-transport}
\uses{hir:g1-filtered-weighted-presentation-reference,hir:g1-independent-graph-common-reduction}
\lean{Hironaka.HironakaG1WeightedChartTransport.map_adaptedIdeal_transition,Hironaka.HironakaG1WeightedChartTransport.native_point_weighted_data_in_given_chart,Hironaka.HironakaG1WeightedChartTransport.locWPEquiv_of_independentChartFilteredBlockChanges,Hironaka.HironakaG1WeightedCrossMark.locWPEquiv_crossed_of_independentBlockReferences,Hironaka.HironakaG1GraphPowerReduction.locWPEquiv_independent_uniformNative_blockGraphs}
\leanok
The true holomorphic transition of independent centered adapted charts
transports their whole source and compressed ideals. Actual point
presentations supply these data in a separately chosen adapted chart.
Filtered matching at a common cross mark, or equality of the crossed
normalized references, gives equivalence of the whole local test class.
In the uniform-weight case, the native point presentations themselves
supply the reductions to the original common analytic source. Their
actual finite normal graphs then compare across the source
biholomorphism with the full ordered boundary.
\end{theorem}
\begin{proof}
Compose the target chart inverse with the source chart and take its
local ring map. This map commutes with ideal products and powers.
At crossed marks, raise the block reduction identity before applying
the proper-carrier comparison. For uniform weights the source
vanishing ideal pulls back along the chart biholomorphism.
Compare each graph record to that same source reference.
\end{proof}

\begin{theorem}[Later normal contact and the literal section cut]
\label{hir:g1-later-selected-section-class}
\uses{hir:g1-later-positive-mask-contact,hir:g1-initial-selected-normal-section-local-class}
\lean{Hironaka.HironakaG1LaterSelectedSectionClass.selected_normal_contact_section_class_on_window,Hironaka.HironakaG1LaterSelectedSectionClass.later_selected_normal_contact_section_class}
\leanok
The genuine normal-contact restriction on the later positive
subwindow and the supported section member cut of the same tested
upper prefix have the same carrier, full ordered boundary, and
local weak product-test class.
\end{theorem}
\begin{proof}
The actual member restrictions of the padded and supported records
on the common window share the local class. The verified normal
tangent cuts out exactly the strict section trace, and the
reference-contact theorem gives the local comparison with the
normal-contact coefficient record.
\end{proof}

\begin{theorem}[True order from later slot-preserving fields]
\label{hir:g1-later-selected-true-order}
\uses{hir:g1-later-positive-mask-contact,hir:g1-original-later-positive-successor-read}
\lean{Hironaka.HironakaG1LaterSelectedNormalOrder.divExpAlong_eq_zero_of_normal_power_field,Hironaka.HironakaG1LaterSelectedNormalOrder.later_selected_normal_order_data}
\leanok
If the genuine later tangent admits the stated normalized tangent
fields for the touching retained slots on its mask, each retained
divisor exponent is zero, and the literal true-mark companion
has the carrier-restricted ideal. The order remains exactly the
true mark at every surviving cosupport endpoint of every actual
tested lower continuation.
\end{theorem}
\begin{proof}
The true power of the normal function lies in the actual selected
coefficient ideal. A positive divisor exponent would make that
function vanish along the selected trace of the touching divisor.
The normalized tangent field preserving the trace contradicts this.
The true-mark companion law then preserves the order along tests.
\end{proof}

\begin{theorem}[The born exceptional slot in the joint example]
\label{hir:g1-joint-example-later-born-level}
\uses{hir:g1-joint-example-later-selected-finite}
\lean{Hironaka.HironakaG1LaterPositiveBornLevelExample.joint_first_blowup_selected_boundary,Hironaka.HironakaG1LaterPositiveBornLevelExample.joint_positive_inner_first_component_changes,Hironaka.HironakaG1LaterPositiveBornLevelExample.joint_first_blowup_upper_regime_run}
\leanok
In the actual $y$-chart the later selected record retains exactly
the born exceptional slot. Its next Step~1 first component is one,
while at the initial selected anchor the same component is zero.
Both upper levels are one. Along this particular upper trajectory
the positive regime starts at the initial stage.
\end{theorem}
\begin{proof}
The ordered selection removes the strict original divisor and
leaves the actual exceptional set. This set contains the chart
origin and hence contributes one to the next level. At the initial
selection there are no remaining incident slots. The regime
calculation uses the initial and later nonspecial anchors.
\end{proof}

\begin{theorem}[Exact initial finite reads on the special companion branch]
\label{hir:g1-initial-companion-section-special-transfer}
\uses{hir:g1-initial-companion-section-order-ratio}
\lean{Hironaka.HironakaG1InitialCompanionSectionSpecialRead.initial_companion_normal_section_finite_special_reads,Hironaka.HironakaG1InitialCompanionSectionSpecialTransfer.initial_order_special_iff_of_empty_slots,Hironaka.HironakaG1InitialCompanionSectionSpecialTransfer.initial_companion_normal_section_finite_special_read}
\leanok
The actual contact and supported section order masks coincide
at the initial finite read, and their companions are special
together. When the contact companion is special the two
recursive values have the same inner level-zero top payload,
with first entries related by the exact mark ratio.
\end{theorem}
\begin{proof}
Empty retained slots and normalized order identify both literal
masks. The scaled companion class comparison gives the original
Step~1 skeleton on the common mask. Its special branch reads
the same level-zero top, and the nonmonomial first letters are
related by the previously established ratio.
\end{proof}

\begin{theorem}[The initial mixed stratum maximum]
\label{hir:g1-initial-positive-mixed-stratum}
\uses{hir:g1-original-initial-positive-nested-top,hir:g1-initial-positive-finite-stratum}
\lean{Hironaka.HironakaG1InitialPositiveMixedStratum.initial_finite_le_section_special_and_lt_iff,Hironaka.HironakaG1InitialPositiveMixedStratum.initial_mixed_stratum_max_locus}
\leanok
On an initial same-positive-level stratum of padded-nonspecial
cosupport, a realized section-special anchor bounds every
padded read. Equality holds exactly at the section-special
anchors. If such an anchor exists, those anchors are the exact
padded maximum locus.
\end{theorem}
\begin{proof}
The actual companion has the same source-wide true mark and
level zero at every anchor. Section speciality gives its inner
top branch. Nonspeciality gives a finite inner read, so the
lexicographic comparison reduces to a strict comparison with top.
\end{proof}

\begin{theorem}[Actual later slot fields and true order]
\label{hir:g1-later-selected-actual-true-order}
\uses{hir:g1-later-selected-true-order,hir:g1-later-positive-mask-contact}
\lean{Hironaka.HironakaG1LaterSelectedActualOrder.exists_later_selected_normal_slot_fields,Hironaka.HironakaG1LaterSelectedActualOrder.exists_later_selected_actual_normal_order}
\leanok
On a subwindow of the actual later positive mask containing a
nonspecial upper cosupport anchor, the true normal contact has
normalized fields preserving the touching retained slots. On this
same window, all retained divisor exponents vanish. The true-mark
companion is the whole carrier restriction, and the raw order along
each actual tested lower continuation stays at that mark on cosupport.
\end{theorem}
\begin{proof}
At a touching retained slot, include that slot along with the
selected old slots in the tangent calculation. The transverse
upper and strict section records supply the normalized fields
near the anchor. Finiteness of the ordered slots gives a
single subwindow on which the true-order criterion applies.
\end{proof}

\begin{theorem}[The later supported domain and canonical comparison]
\label{hir:g1-later-selected-section-canonical}
\uses{hir:g1-later-selected-section-class}
\lean{Hironaka.HironakaG1LaterSelectedSectionCanonical.exists_born_member_of_mem_retained_on_mask,Hironaka.HironakaG1LaterSelectedSectionCanonical.later_selected_section_domain_and_canonical_iff}
\leanok
The actual prescribed later contact transfers its invariant
domain to the literal supported section cut on the common window.
For canonicality comparison the normalized divisor exponents
need only agree on the retained born members meeting the
contact cosupport.
\end{theorem}
\begin{proof}
The two records share a carrier, boundary and cosupport, and the
selected section ideal is coherent at its positive mark. These
data transfer the domain. A retained old member misses the
positive mask. Apply the original class comparison with the
remaining normalized exponents.
\end{proof}

\begin{theorem}[The finite stratum from the actual supported lower reads]
\label{hir:g1-initial-finite-section-stratum}
\uses{hir:g1-initial-companion-prescribed-section,hir:g1-initial-positive-finite-stratum}
\lean{Hironaka.HironakaG1InitialCompanionSectionPairOrder.initial_finite_section_records,Hironaka.HironakaG1InitialCompanionSectionPairOrder.FiniteSectionPair.lt_iff,Hironaka.HironakaG1InitialCompanionSectionPairOrder.initial_finite_section_stratum_max_loci}
\leanok
When the actual initial finite companions are nonspecial, the
prescribed contacts and supported section cuts on their literal
masks preserve and reflect the order of the two-anchor reads.
On such an initial same-level stratum this gives exact upper
maximum-locus images and preimages from the supported reads,
including ties.
\end{theorem}
\begin{proof}
The finite contact branch produces the actual supported lower
records and both domains. The marks are constant across anchors
on each side. Apply the original two-pair comparison at their
literal unequal marks and transport the maximum predicates.
\end{proof}

\begin{theorem}[A positive split with a retained born slot]
\label{hir:g1-joint-example-advanced-head}
\uses{hir:g1-joint-example-later-born-level}
\lean{Hironaka.HironakaG1LaterAdvancedHeadExample.joint_added_boundary_upper_regime_advances,Hironaka.HironakaG1LaterAdvancedHeadExample.joint_advanced_head_step_one}
\leanok
Adding the second coordinate divisor to the actual joint example
gives a tested nonempty upper blowup with level decreasing from
two to one and a nonspecial endpoint. The actual positive run
at that endpoint has length zero. Its split head is the
blowup, and the actual selected mask has next level one from
the retained born exceptional slot.
\end{theorem}
\begin{proof}
The literal chart misses the strict second divisor while meeting
the strict first divisor and the exceptional set. Both upper
records are tested and the controlled generators give the surviving
cosupport point. The preceding level differs from the endpoint
level, so the actual split takes the empty tail at the later
head. The ordered lower record counts the single nonempty slot.
\end{proof}

\begin{theorem}[Exact order of realized later positive reads]
\label{hir:g1-later-positive-branch-order}
\uses{hir:g1-original-later-positive-successor-read,hir:g1-joint-example-advanced-head}
\lean{Hironaka.HironakaG1LaterPositiveBranchOrder.exists_laterPositiveRead,Hironaka.HironakaG1LaterPositiveBranchOrder.LaterPositiveRead.upper_lt_iff,Hironaka.HironakaG1LaterPositiveBranchOrder.LaterPositiveRead.upper_leaf_order_iff,Hironaka.HironakaG1LaterPositiveBranchOrder.LaterPositiveRead.upper_lt_of_companion_level_lt}
\leanok
At a tied original upper level, two realized later positive
reads compare by the actual lower order, then the actual
companion level, then the companion's inner value. If the
order and companion level agree, the inner value decides
strict order, weak order and ties, at any shared companion level.
\end{theorem}
\begin{proof}
Read the two actual positive successors separately on their own
chronological splits. Expand the resulting lexicographic order
at the common original level. The remaining coordinates are
exactly the measured order, companion level and inner value.
\end{proof}

\begin{theorem}[Native weighted comparisons from owned coefficients]
\label{hir:g1-native-prefix-weight-comparison}
\uses{hir:g1-independent-weighted-chart-transport}
\lean{Hironaka.HironakaG1NativePrefixWeightComparison.native_owned_coefficients_separate,Hironaka.HironakaG1NativePrefixWeightComparison.native_prefix_match_identifies_germ,Hironaka.HironakaG1NativePrefixWeightComparison.locWPEquiv_of_native_prefix_comparisons}
\leanok
The actual native prefix-support and supported-division laws
reduce normalized reference comparison to agreement on owned
coefficients. For each generator use a representative in the
matched-degree span plus the normalized reference, with vanishing
lower owned coefficients and the given higher owned coefficients.
Mutual matches give the whole weak test class on the actual
common carrier and ordered boundary.
\end{theorem}
\begin{proof}
Supported division and the whole-jet laws identify two native
germs with the same owned coefficients. Native prefix support
supplies the lower-side vanishing for the generator. The
resulting filtered generator change compares the normalized
block references at the literal cross mark.
\end{proof}

\begin{theorem}[Full later true order on the actual lower mask]
\label{hir:g1-later-actual-full-order-run}
\uses{hir:g1-later-selected-actual-true-order,hir:g1-original-later-positive-successor-read}
\lean{Hironaka.HironakaG1ActualLaterSuccessorRead.original_upper_positive_lower_raw_mark,Hironaka.HironakaG1ActualLaterSuccessorRead.original_upper_positive_successor_at_true_mark,Hironaka.HironakaG1ActualLaterOrderRun.original_upper_positive_lower_full_order_run,Hironaka.HironakaG1ActualLaterOrderRun.original_upper_positive_successor_at_initial_order_companion}
\leanok
At a positive nonspecial upper cosupport endpoint of an
actual Step~1 canonical history, the full lower tested tail
has order equal to its true mark at each surviving projected
prefix. The actual upper finite read uses the same-mark
companion at the initial stage of that lower tail and retains
its full ordered boundary on the literal order mask.
\end{theorem}
\begin{proof}
The upper canonical read supplies the actual lower canonical
tail. Choose one normal-contact subwindow through its starting
anchor. Nested open restriction compares its exact lower
transforms with the full-mask transforms. The true order on
the window transfers to every surviving projected prefix.
Thus the whole lower tail lies in the same order regime and
the exact successor formula reads its initial order companion.
\end{proof}

\begin{theorem}[An actual canonical later split needs its born level]
\label{hir:g1-joint-example-canonical-born-level}
\uses{hir:g1-joint-example-advanced-head,hir:g1-later-actual-full-order-run}
\lean{Hironaka.HironakaG1LaterAdvancedHeadCanonicalExample.joint_advanced_head_isCanonical,Hironaka.HironakaG1LaterAdvancedHeadCanonicalExample.joint_advanced_canonical_positive_successor,Hironaka.HironakaG1LaterAdvancedHeadCanonicalExample.joint_advanced_canonical_not_level_zero_inner}
\leanok
The actual single point blowup in the enlarged-boundary joint
example is canonical for the original upper invariant. Its later
positive value has true lower order one and reads the actual
order companion at inner level one. Inserting inner level zero
changes the whole upper value for every inner leaf.
\end{theorem}
\begin{proof}
The initial cosupport is exactly the blown-up point. Its
successor advances the actual positive split. The true-order
read identifies the ordinary companion on its own order mask.
The born exceptional divisor gives inner level one, which
cannot equal zero in this lexicographic coordinate.
\end{proof}

\begin{theorem}[Actual recursive domains and canonical class on later cuts]
\label{hir:g1-later-contact-recursive-domain}
\uses{hir:g1-later-selected-section-canonical,hir:g1-later-selected-section-class}
\lean{Hironaka.HironakaG1LaterContactRecursiveDomain.exists_later_contact_recursive_domain_of_section_nonspecial,Hironaka.HironakaG1LaterContactRecursiveDomain.same_mark_companion_tangent_on_order_mask,Hironaka.HironakaG1LaterSelectedSectionGapCanonical.later_selected_section_canonical_iff_of_contact_domain,Hironaka.HironakaG1LaterSectionPairOrder.later_section_comparison,Hironaka.HironakaG1LaterSectionPairOrder.LaterSectionComparison.lt_iff}
\leanok
At a positive upper cosupport anchor that is nonspecial on the
supported section, the genuine later contact window in both
literal masks supplies the contact and section recursive
domains, with the existing ordered born members. Whenever
such a prescribed contact has its recursive domain, its
normalized exponents and original canonical histories agree
with the section cut. Two such later canonical endpoints
compare by their associated supported section values, even
with independent cuts and lower histories.
\end{theorem}
\begin{proof}
The actual section mask supplies its selected domain. The
contact has the same carrier and cosupport, so this domain
transfers. The two coefficient gaps identify normalized
exponents for the entire ordered boundary. The original
unequal-mark comparison gives canonicality. For two endpoints,
the literal coefficient marks on each side depend only on the
original source, so the original two-pair cut compares them.
\end{proof}

\begin{theorem}[The initial positive stratum with both branches]
\label{hir:g1-initial-positive-all-stratum}
\uses{hir:g1-initial-positive-mixed-stratum,hir:g1-initial-finite-section-stratum}
\lean{Hironaka.HironakaG1InitialPositiveAllStratumCentres.initial_positive_all_stratum_max_loci}
\leanok
The realized actual initial same-level stratum has a common
supported comparison of maxima. At section-special anchors the
literal section upper value is top at that level. At nonspecial
anchors the actual supported lower read decides the finite
comparison. This includes ties and the empty stratum.
\end{theorem}
\begin{proof}
Read the branches using the literal section speciality. The
special anchor dominates the finite branch at their shared
initial order. Two finite anchors compare by their actual
supported lower reads, so the pairwise order criterion
transports the maximum set in each case.
\end{proof}

\begin{theorem}[Later highest-level maxima and actual canonical steps]
\label{hir:g1-later-highest-stratum-canonical-step}
\uses{hir:g1-later-positive-branch-order}
\lean{Hironaka.HironakaG1LaterPositiveStratumCentres.later_positive_stratum_max_loci,Hironaka.HironakaG1LaterPositiveStratumCentres.later_positive_top_stratum_canonical_step_iff}
\leanok
Pairwise order agreement on realized later positive tuples
transports maximum loci on the corresponding strata. When
those strata are the full nonempty highest-level cosupport
strata and the actual centres and prefix canonicalities
correspond, the next actual blowups have equivalent canonicality.
\end{theorem}
\begin{proof}
Translate the realized tuple comparisons into upper value
comparisons. Every point at a lower original level is
strictly below every point in the nonempty highest stratum.
The transported highest-stratum maxima thus identify both
full maxima. Apply the original canonical successor criterion.
\end{proof}

\begin{theorem}[Source weak equivalence on independent normal graphs]
\label{hir:g1-source-weak-graph-lift}
\uses{hir:g1-independent-uniform-native-class}
\lean{Hironaka.HironakaG1SourceWeakClassGraphLift.locWPEquiv_sectionMarked_of_source,Hironaka.HironakaG1SourceWeakClassGraphLift.locWPEquiv_properNormalMarked_of_source,Hironaka.HironakaG1SourceWeakClassGraphLift.locWPEquiv_selectorMarked_of_source,Hironaka.HironakaG1SourceWeakClassGraphLift.locWPEquiv_selectorMarked_of_native}
\leanok
Source weak test equivalence lifts to the support-preserving
zero sections and to independent proper normal graphs. The
source marks and carriers may differ. The graph chart maps
pull back every added boundary member in order. Independent
native weighted data on the same intrinsic Hilbert--Samuel
cut supply the source weak equivalence.
\end{theorem}
\begin{proof}
Apply the source test equivalence to the line projection,
adjoin the zero divisor and take its selected member. Its
relative coefficient trace is the section class. For the
proper graph, include the normal generator at the original
mark and transport through the graph shift.
\end{proof}

\begin{theorem}[The actual later companion and its born-aware selected cut]
\label{hir:g1-actual-later-selected-companion-tangent}
\uses{hir:g1-later-actual-full-order-run,hir:g1-later-contact-recursive-domain}
\lean{Hironaka.HironakaG1ActualLaterCompanionTangent.exists_actual_later_initial_order_companion_tangent,Hironaka.HironakaG1BornSelectedCompanionPhase.exists_born_selected_contact_of_actual_compAt}
\leanok
The literal first order companion on the actual later lower
tail has its domain and a tangent whose zero trace is the
strict section. On a positive nonspecial member of an actual
order companion, a prescribed tangent with the supported
section as its trace gives a selected cut of one lower
recursion depth and a coefficient cut weakly equivalent to
the corresponding supported cut. The selection retains the
actual ordered incidence of the born members.
\end{theorem}
\begin{proof}
Transfer the upper strict section direction through the
same-mark companion on the exact order mask. Shrink the
tangent window inside the next positive mask and pass to
the selected member there. The coefficient trace compares
with the supported section through the local class theorem.
\end{proof}

\begin{theorem}[Initial comparison with the literal section upper read]
\label{hir:g1-initial-literal-section-stratum}
\uses{hir:g1-initial-positive-all-stratum}
\lean{Hironaka.HironakaG1InitialFiniteSectionUpperRead.initial_finite_section_upper_read,Hironaka.HironakaG1InitialFiniteSectionUpperRead.initial_positive_literal_section_stratum}
\leanok
At the initial positive level the padded nonspecial stratum
has the same strict order and maxima as the literal supported
section upper read. This applies to both branches of the
section invariant, including ties and the empty stratum.
\end{theorem}
\begin{proof}
The section lower read on its own nested masks equals the
value on the common supported cut. In the section-special
case the upper value is the top branch. Compare the two
resulting upper reads at the shared level.
\end{proof}

\begin{theorem}[Canonical paired histories from their actual predecessor reads]
\label{hir:g1-later-positive-prefix-history}
\uses{hir:g1-later-highest-stratum-canonical-step}
\lean{Hironaka.HironakaG1LaterPositiveCANHistory.PrefixRead.max_image,Hironaka.HironakaG1LaterPositiveCANHistory.canonical_iff_of_later_positive_prefix_reads}
\leanok
For each proper predecessor of a genuine padded and native
lift, it suffices to give a realized top-stratum comparison,
empty cosupports on both sides, or exact maximum-locus
readback. Under these conditions the two entire histories
have equivalent canonicality.
\end{theorem}
\begin{proof}
The realized order comparison transports the highest-stratum
maxima, which are the full maxima. In the empty and exact
readback cases the full maximum equality is immediate. Then
induct on the actual paired histories using the centre maps.
\end{proof}

\begin{theorem}[A finite inner leaf at the actual born level]
\label{hir:g1-joint-example-finite-born-leaf}
\uses{hir:g1-joint-example-canonical-born-level}
\lean{Hironaka.HironakaG1LaterAdvancedHeadLeafExample.joint_advanced_head_finite_inner_leaf}
\leanok
The level-one companion in the canonical advanced-head example
takes its finite branch. The final recursive read is on the
next literal selected mask, with empty history, mark one and
the ordered boundary consisting of one empty member.
\end{theorem}
\begin{proof}
On the actual strict chart a companion generator cuts out a
smooth normal direction while the born member cuts out the
other coordinate. The companion is nonspecial, and its
positive Step~1 formula gives the recursive read on the
selected mask.
\end{proof}

\begin{theorem}[Independent weighted extractions along native test prefixes]
\label{hir:g1-native-weighted-prefix-state}
\uses{hir:g1-native-prefix-weight-comparison,hir:g1-source-weak-graph-lift}
\lean{Hironaka.HironakaG1NativeWeightedPrefixState.independent_nativeHSStates_at_test_prefix,Hironaka.HironakaG1NativeWeightedPrefixState.independent_native_extractions_at_test_prefix,Hironaka.HironakaG1NativeWeightedPrefixState.locWPEquiv_selectorMarked_after_native_test_prefix}
\leanok
Two actual native records over the same intrinsic cut retain
the same intrinsic analytic source and Hilbert--Samuel value
through every common weak test prefix. Their independent
weighted extractions in an adapted chart present the whole
transported source ideal. On an open coordinate stage these
states give the comparison of the actual selector graphs.
\end{theorem}
\begin{proof}
Track strict transforms, projections and open restrictions of
the intrinsic source. At each step the native test-state
theorem gives the common transported value. Extract each native
presentation separately in the chosen target chart and apply
the graph class comparison using the transported domains.
\end{proof}

\begin{theorem}[The Hilbert--Samuel cuts on the actual common graph]
\label{hir:g1-common-graph-hs-cuts}
\uses{hir:g1-native-weighted-prefix-state}
\lean{Hironaka.HironakaG1CommonGraphHSCut.commonGraphSubset_analytic,Hironaka.HironakaG1CommonGraphHSCut.ambient_cut_and_bound_of_intrinsic_iso,Hironaka.HironakaG1CommonGraphHSCut.exists_commonGraph_ambient_cuts}
\leanok
An intrinsic isomorphism of two embedded reduced analytic
sets gives a local common analytic graph. The actual two
projections identify every positive ambient Hilbert--Samuel
cut with its image in the original embeddings. The ambient
upper bounds transport as well, including points off the
analytic carriers.
\end{theorem}
\begin{proof}
The graph is cut out by the base analytic ideal and the
holomorphic graph equations. Each projection induces the
isomorphism of analytic stalks on the carriers. Off these
carriers the stalk of the vanishing ideal is the unit ideal,
so the quotient has Hilbert--Samuel value zero.
\end{proof}

\begin{theorem}[Whole supported division after one closed normal extension]
\label{hir:g1-closed-normal-supported-division}
\uses{hir:g1-common-graph-hs-cuts}
\lean{Hironaka.HironakaG1ClosedNormalDivision.normalOwner_of_one_le,Hironaka.HironakaG1ClosedNormalDivision.normalOwner_liftExp_iff,Hironaka.HironakaG1ClosedNormalDivision.historicalWholeSupportedDivision_closedNormal}
\leanok
Adjoining the actual degree-one normal equation to an
extended weighted source family preserves whole supported
division. Positive normal exponents belong to the new
generator, and exponents on the zero-normal slice preserve
their original owners in the extended ordered family.
\end{theorem}
\begin{proof}
Divide the positive-normal part by the normal coordinate.
The remaining zero-normal slice is in the original ideal,
where the source division law supplies the owned
multipliers. Extend them through the zero-normal slice.
\end{proof}

\begin{theorem}[Actual later born cuts and original canonicality]
\label{hir:g1-actual-born-cut-canonicality}
\uses{hir:g1-actual-later-selected-companion-tangent}
\lean{Hironaka.HironakaG1ActualLaterBornCut.exists_actual_later_born_cut,Hironaka.HironakaG1ActualLaterBornCut.supported_member_carrier_on_initial_order_mask,Hironaka.HironakaG1BornSelectedCompanionCanonical.born_selected_contact_canonical_of_domain,Hironaka.HironakaG1BornSelectedCompanionCanonical.exists_born_selected_contact_canonical_of_supported_nonspecial,Hironaka.HironakaG1ActualLaterBornCanonical.exists_actual_later_born_canonical}
\leanok
The actual later companion tangent gives a next member cut
with the inherited born slots when the supported reference
supplies its domain and class on the order mask. The literal
supported member has the required strict-section carrier.
For the realized contact and supported cuts, either recursive
domain gives both domains and equivalent original canonicality.
When the supported anchor is nonspecial, its literal mask
supplies the domains on a common open neighborhood.
For the literal supported reference with its domain and
class, its positive nonspecial anchor supplies companion
nonspeciality in the actual canonical history.
\end{theorem}
\begin{proof}
The actual tangent's zero trace yields the supported carrier.
Apply positive selection with every ordered incident member.
The coefficient cosupport gaps determine the normalized
boundary exponents. Transport the recursive domain and apply
the original class comparison on the common window.
\end{proof}

\begin{theorem}[The literal native maximum locus at the initial prefix]
\label{hir:g1-initial-literal-native-read}
\uses{hir:g1-later-positive-prefix-history,hir:g1-initial-literal-section-stratum}
\lean{Hironaka.HironakaG1InitialLiteralNativeCentres.initial_literal_native_max_image,Hironaka.HironakaG1InitialLiteralNativeCentres.initial_literal_native_prefix_read}
\leanok
When the initial padded cosupport is nonspecial at one shared
positive level and the section has maximal order, the entire
initial padded maximum locus is the image of the literal native
maximum locus. This supplies the first actual predecessor read
for the paired canonical histories.
\end{theorem}
\begin{proof}
Identify the initial padded and section cosupports and compare
the actual strict order on them. The supported maximal-order
lift maps the native maxima to the section maxima. Compose
the identifications at the empty prefix.
\end{proof}

\begin{theorem}[Highest positive initial native maxima]
\label{hir:g1-initial-literal-native-top-read}
\uses{hir:g1-initial-literal-native-read}
\lean{Hironaka.HironakaG1InitialLiteralNativeTopCentres.initial_padded_positive_not_special,Hironaka.HironakaG1InitialLiteralNativeTopCentres.initial_literal_native_top_max_image,Hironaka.HironakaG1InitialLiteralNativeTopCentres.initial_literal_native_top_prefix_read}
\leanok
An initial positive padded cosupport point is nonspecial.
With an attained positive highest level and section maximal
order, the maximum-locus image and the initial predecessor
read hold for the full cosupport, including points at
arbitrary lower levels.
\end{theorem}
\begin{proof}
The initial normal line preserves incidence in every original
member while its nearby nonzero points leave the zero-section
cosupport. Thus all highest-level points are nonspecial on the
padded side. At their common level apply the literal section
order comparison. Lower first entries give no additional
maxima on either side.
\end{proof}

\begin{theorem}[Domain of the supported reference on the exact later order window]
\label{hir:g1-later-supported-exact-order-window}
\uses{hir:g1-actual-born-cut-canonicality}
\lean{Hironaka.HironakaG1LaterOrderMaskSupportedLocal.supported_domain_on_nested_order_window,Hironaka.HironakaG1LaterOrderMaskSupportedLocal.later_supported_reference_on_exact_order_window}
\leanok
Supported nonspeciality supplies an open neighborhood of the
actual anchor inside the exact later companion order mask. The
literal supported reference on this window has its recursive
domain. Its local weak test class agrees with that of the
companion when the localized padded lower record and companion
are restricted equivalent.
\end{theorem}
\begin{proof}
Intersect the two open masks through their pullback diagram
and restrict the ordinary supported recursive domain. Compose
the original local padded--section class with the given
restricted comparison on the same window.
\end{proof}

\begin{theorem}[Local companion equivalence from the actual nonfactor window]
\label{hir:g1-later-order-nonfactor-window}
\uses{hir:g1-later-supported-exact-order-window}
\lean{Hironaka.HironakaG1LaterOrderMaskNonfactor.initial_companion_restrictEquiv_on_nonfactor_window,Hironaka.HironakaG1LaterOrderMaskNonfactor.initial_companion_restrictEquiv_of_ambient_nonfactor}
\leanok
On the subwindow supplied by the actual retained zero
exponents, the exact same-mark companion and the padded
lower record are restricted equivalent on the actual
order mask.
\end{theorem}
\begin{proof}
The nonfactor identity gives the same-mark companion on
the ambient window. Restrict through the nested opens
of the actual order mask and carry the comparison under
their exact transition maps.
\end{proof}

\begin{theorem}[An actual common supported reference on the later order mask]
\label{hir:g1-actual-later-common-supported-reference}
\uses{hir:g1-later-order-nonfactor-window,hir:g1-later-supported-exact-order-window}
\lean{Hironaka.HironakaG1ActualLaterCommonReference.exists_actual_later_common_supported_reference}
\leanok
At the actual positive nonspecial split, suppose the literal
supported section is nonspecial at the split anchor. There is
one open neighborhood on the exact order mask where the
actual companion and literal supported reference each have
their recursive domain and are locally weakly equivalent.
The companion keeps its actual mark and ordered members.
\end{theorem}
\begin{proof}
Intersect the supported domain window and the nonfactor
window inside the order mask. Construct the domains on this
intersection by restriction. Combine the literal padded and
section comparison with the actual local companion class.
\end{proof}

\begin{theorem}[Literal later section read at a positive restart]
\label{hir:g1-later-literal-section-restart}
\uses{hir:g1-later-contact-recursive-domain,hir:g1-initial-literal-section-stratum}
\lean{Hironaka.HironakaG1LaterLiteralSectionUpperRead.later_literal_section_restart_read,Hironaka.HironakaG1LaterLiteralSectionUpperRead.later_literal_section_restart_stratum}
\leanok
When the literal supported section starts its positive
nonspecial run at the current actual stage, its upper finite
read is given by the supported section cut on the actual
inherited masks. The contact values and literal section upper
values give the same strict order and maxima on a shared
positive restart stratum.
\end{theorem}
\begin{proof}
Use the true contact and supported recursive domains on the
common open. The literal restart uses the empty lower tail
and retains the existing ordered born members. Transport
the supported cut value across its open restriction, and
apply the two-anchor comparison on that stratum.
\end{proof}

\begin{theorem}[Finite section reads along the actual canonical suffix]
\label{hir:g1-later-literal-section-canonical-tail}
\uses{hir:g1-later-literal-section-restart}
\lean{Hironaka.HironakaG1LaterLiteralSectionCanonicalTail.later_literal_section_finite_read,Hironaka.HironakaG1LaterLiteralSectionCanonicalTail.later_literal_section_finite_stratum}
\leanok
For an original canonical section history of a maximal-order
source, the finite read uses the actual positive nonspecial
run and its whole suffix. At a shared level, the associated
contact values therefore give the same strict order and
maximum locus as the literal upper section values.
\end{theorem}
\begin{proof}
Split the canonical history at the actual run head and
construct the recursive comparison domains there. Restrict
the lower canonical tail to the common window and transport
it to the actual nested masks. Apply the exact read and the
same-level comparison for the two contacts.
\end{proof}

\begin{theorem}[Weak test equivalence through a finite normal block]
\label{hir:g1-finite-source-weak-graph-lift}
\uses{hir:g1-source-weak-graph-lift,hir:g1-native-weighted-prefix-state}
\lean{Hironaka.HironakaG1FiniteSourceWeakClassGraphLift.locWPEquiv_blockMarked_of_source,Hironaka.HironakaG1FiniteSourceWeakClassGraphLift.locWPEquiv_blockMarked_after_native_test_prefix}
\leanok
The source weak class lifts through any finite normal block
and through a base biholomorphism between independent graph
charts. It retains both original marks and the complete
ordered extra boundary. A common native test prefix supplies
the source weak class and domains at its coordinate-open
endpoint.
\end{theorem}
\begin{proof}
Iterate the one-coordinate zero-section comparison. Projection
composition identifies the whole trace with the actual block
trace. The normal powers vanish along that graph in restricted
equivalence. Transport the result through the chart change
and pull back the additional members in order.
\end{proof}

\begin{theorem}[The original record on the native manifold of a finite graph]
\label{hir:g1-finite-native-graph-record}
\uses{hir:g1-finite-source-weak-graph-lift}
\lean{Hironaka.HironakaG1FiniteNativeGraphSection.nativeBlockGraph_regular,Hironaka.HironakaG1FiniteNativeGraphSection.nativeBlockGraphEquiv,Hironaka.HironakaG1FiniteNativeGraphSection.pullback_native_blockMarked}
\leanok
The actual native manifold of a full holomorphic normal graph
is biholomorphic to its source open. The marked pullback to
that manifold recovers the whole original record in its source
chart, with each additional ambient member restricted to the
graph in its original order.
\end{theorem}
\begin{proof}
The section and projection are inverse, and holomorphic for
the flattening atlas. The composite of the native inclusion
and the graph section pulls the native record back to the
literal graph trace. All normal generators vanish there.
\end{proof}

\begin{theorem}[The original canonical history on the literal source of a finite graph]
\label{hir:g1-finite-native-graph-canonical-history}
\uses{hir:g1-finite-native-graph-record,hir:g1-original-native-chart-comparison}
\lean{Hironaka.HironakaG1FiniteNativeGraphSection.exists_original_native_block_history}
\leanok
An actual original canonical history of a graph-supported
marked record gives the literal source a canonical history
of the same length, with the ordered restricted boundary.
The original recursive values have the shift by the normal
codimension. The original maximum loci are exact images and
preimages, and a resolved endpoint remains resolved.
\end{theorem}
\begin{proof}
Use the canonical native lift on the actual full graph.
Transport that lift through the graph's source
biholomorphism and its literal marked trace. Combine the
original dimension-shifted value map with the two transport
maps for the entire endpoint cosupport and its maxima.
\end{proof}

\begin{theorem}[Exact numerical leaf of the advanced-head example]
\label{hir:g1-joint-example-exact-born-leaf}
\uses{hir:g1-joint-example-finite-born-leaf}
\lean{Hironaka.HironakaG1LaterAdvancedHeadNumericalExample.joint_advanced_head_exact_upper_value}
\leanok
The original upper value of the actual canonical example has
outer level one, positive order one, born inner level one,
and positive order one. Its last branch is special at level
zero and has the top leaf.
\end{theorem}
\begin{proof}
The lifted normal line on the last selected mask detects
order one. The resulting smooth curve carrier has a singleton
cosupport germ, so its order companion is special. The ordered
empty boundary member contributes zero to the level. Substitute
the exact last read into the two positive successors.
\end{proof}

\begin{theorem}[Finite and prefix jet frames after adjoining the normal equation]
\label{hir:g1-closed-normal-jet-frames}
\uses{hir:g1-closed-normal-supported-division}
\lean{Hironaka.HironakaG1ClosedNormalJet.historicalWholeJetMap_closedNormal_bijective,Hironaka.HironakaG1ClosedNormalJet.historicalDegreePrefixFrames_closedNormal}
\leanok
Adjoining the degree-one normal equation preserves the whole
finite jet frame and the degree-sorted prefix frames. The
prefix statement includes the homogeneous frames and every
degree cutoff, including zero.
\end{theorem}
\begin{proof}
In the zero-normal part the columns are the lifted original
columns. In the positive-normal part the new equation gives
identity columns. The decomposition restricts to every
degree prefix and its homogeneous piece.
\end{proof}

\begin{theorem}[The actual tangent and next born cut on the common window]
\label{hir:g1-localized-next-born-canonicality}
\uses{hir:g1-actual-later-common-supported-reference,hir:g1-actual-born-cut-canonicality}
\lean{Hironaka.HironakaG1ActualLaterCommonTangent.exists_actual_later_common_supported_tangent,Hironaka.HironakaG1LocalizedLaterBornCanonical.initial_positive_special_pullback_open_iff,Hironaka.HironakaG1LocalizedLaterBornCanonical.exists_actual_localized_later_born_canonical}
\leanok
The two actual localized records carry the genuine tangent
whose zero trace equals the literal supported carrier. For
the positive born level and the stated supported nonspecial
anchors, the next selected cuts have their literal coefficient
marks and recursive domains. The selected contact and literal
supported cut have equivalent original canonical histories.
\end{theorem}
\begin{proof}
Intersect the common domain and class window with the
actual companion tangent window. The strict section trace
identifies the supported carrier on this intersection.
Positive speciality reflects through the open pullback.
Select the born members and apply the coefficient gap
comparison to the two records on the same window.
\end{proof}

\begin{theorem}[Initial native read with empty or singleton cosupport]
\label{hir:g1-initial-small-support-native-read}
\uses{hir:g1-initial-literal-native-read}
\lean{Hironaka.HironakaG1InitialLiteralNativeSmallSupport.initial_literal_native_subsingleton_max_image,Hironaka.HironakaG1InitialLiteralNativeSmallSupport.initial_literal_native_subsingleton_prefix_read}
\leanok
When the initial padded cosupport has at most one point,
its whole maximum locus is the image of the literal native
maximum locus. This supplies the actual initial predecessor
read, including the empty and level-zero singleton cases.
\end{theorem}
\begin{proof}
The native inclusion identifies the actual cosupports. A
set with at most one point is its own maximum locus for
each of the two value functions.
\end{proof}

\begin{theorem}[Actual mixed and highest strata of the later section invariant]
\label{hir:g1-later-literal-mixed-section-stratum}
\uses{hir:g1-later-literal-section-canonical-tail}
\lean{Hironaka.HironakaG1LaterLiteralSectionMixedStratum.later_literal_section_special_read,Hironaka.HironakaG1LaterLiteralSectionMixedStratum.later_literal_section_mixed_stratum,Hironaka.HironakaG1LaterLiteralSectionMixedStratum.later_literal_section_max_locus_highest_level}
\leanok
At a common positive level on an actual canonical section
history, the finite contact values and special top values
recover the literal section order and maximum loci. If a
special point occurs at that level, exactly the special
points maximize. An attained highest level suffices for
the whole comparison on a set spanning several levels.
\end{theorem}
\begin{proof}
Use the actual chronological suffix at each finite anchor
and the original special branch at each special anchor.
At the shared level a special value is top and ties every
other special value. For several levels, the first entry
gives strict dominance over every lower level.
\end{proof}

\begin{theorem}[Canonical comparison in a chart of a genuine native lift]
\label{hir:g1-original-native-chart-comparison}
\uses{hir:g1-same-bound-native-maxima-canonicality}
\lean{Hironaka.HironakaG1NativeChartCanonicalComparison.original_native_chart_comparison}
\leanok
An independent chart of the actual native tower preserves
and reflects original canonicality of a supported ambient
lift. At the ambient offset bound the source chart has the
same cosupport and maxima under exact images and preimages.
Its original recursive values transport with the dimension
shift.
\end{theorem}
\begin{proof}
Pull back both canonical histories along the stagewise
tower isomorphism and its inverse. Use the original
supported bound comparison. The endpoint inclusion and
biholomorphism cover the ambient cosupport. Strict
monotonicity of the recursive shift gives both maxima
formulas on it.
\end{proof}

\begin{theorem}[Two ambient realizations of the same native chart history]
\label{hir:g1-cross-ambient-native-comparison}
\uses{hir:g1-original-native-chart-comparison}
\lean{Hironaka.HironakaG1CrossAmbientNativeComparison.NativeChartRealization.compare,Hironaka.HironakaG1CrossAmbientNativeComparison.cross_ambient_native_comparison}
\leanok
Two supported ambient histories with the same literal native
chart history have equivalent original canonicality. The
paired points of their cosupports give the same strict
order, although the two ambient dimensions and recursive
offset bounds may differ. The maximum locus of the common
chart maps onto both original maxima, with exact preimages.
\end{theorem}
\begin{proof}
Compare each actual ambient tower with the given native
chart tower at its own offset. The two strictly monotone
value maps preserve and reflect order. Apply their exact
cosupport and maximum-locus transport formulas.
\end{proof}

\begin{theorem}[The actual extended selector and its Jacobian]
\label{hir:g1-closed-normal-actual-selector}
\uses{hir:g1-closed-normal-jet-frames}
\lean{Hironaka.HironakaG1ClosedNormalSelector.historicalWholePrefixSupport_normalGerms,Hironaka.HironakaG1ClosedNormalSelector.normalSelector_contact,Hironaka.HironakaG1ClosedNormalSelector.normalSelector_jacobian_det}
\leanok
The actual linear normal germ and the pulled-back original
germs retain whole prefix support. The extended selector
uses the original choices and the new normal equation.
It has genuine analytic contact and preserves the selected
Jacobian determinant at every cutoff.
\end{theorem}
\begin{proof}
Compare the literal Taylor series of the extended germs.
Lift the original selected derivatives and insert the
derivative of the normal coordinate. At positive cutoffs
the selected matrix is block diagonal with a unit normal
block and the original matrix. At cutoff zero it is empty.
\end{proof}

\begin{theorem}[The native value on the actual advanced-head lift]
\label{hir:g1-joint-example-actual-native-value-obstruction}
\uses{hir:g1-joint-example-exact-born-leaf,hir:g1-original-native-chart-comparison}
\lean{Hironaka.HironakaG1LaterAdvancedHeadNativeExample.joint_advanced_head_literal_native_discrepancy,Hironaka.HironakaG1LaterAdvancedHeadNativeExample.joint_advanced_head_literal_native_no_shift}
\leanok
The actual advanced-head blowup has a literal canonical
native lift. At the matched endpoint, the native upper
read is special top or has finite raw order zero. Thus
any lower value map preserving the first numerical entry
fails to give the cylindrical value of this example.
\end{theorem}
\begin{proof}
The native initial cosupport is the singleton centre,
so the test lift is canonical. At the surviving native
endpoint, the retained born member on the selected curve
absorbs its order. The finite native read has first entry
zero, whereas the actual cylindrical finite read has first
entry one by the exact numerical computation.
\end{proof}

\begin{theorem}[The whole weighted presentation after a closed normal extension]
\label{hir:g1-closed-normal-weighted-presentation}
\uses{hir:g1-closed-normal-actual-selector,hir:g1-closed-normal-supported-division}
\lean{Hironaka.HironakaG1ClosedNormalLaws.actualHistoricalFivePointLaws_closedNormal,Hironaka.HironakaG1ClosedNormalLaws.weightedCommonMark_closedNormal,Hironaka.HironakaG1ClosedNormalLaws.weightedCompression_closedNormal,Hironaka.HironakaG1ClosedNormalNativeState.span_normalGerms,Hironaka.HironakaG1ClosedNormalNativeState.nativePoint_closedNormal,Hironaka.HironakaG1ClosedNormalNativeState.nativeHSState_closedNormal}
\leanok
The family of actual extended analytic germs satisfies all
five historical laws. Its common mark stays fixed, and its
compressed ideal is the extension of the source compressed
ideal plus the normal coordinate to that mark. Given the
true ambient Hilbert--Samuel cut and upper bound and the
two adapted whole-ideal identities, the extended family
supplies the native point presentations and the whole
target native state.
\end{theorem}
\begin{proof}
Combine the actual selector, finite and prefix frames, and
whole supported division. The added generator has degree
one, so it leaves the least common multiple fixed. Its
compressed power generates the new principal summand.
Apply the two adapted ideal identities to the whole
uncompressed and compressed spans.
\end{proof}

\begin{theorem}[Literal maximum loci at an initial zero-level stabilization]
\label{hir:g1-initial-zero-literal-maxima}
\uses{hir:g1-initial-small-support-native-read}
\lean{Hironaka.HironakaG1InitialZeroRecodedCentres.initial_zero_stepOne_max_loci,Hironaka.HironakaG1InitialZeroRecodedCentres.native_mark_one_initial_general_read,Hironaka.HironakaG1InitialZeroRecodedCentres.initial_zero_stepOne_literal_read,Hironaka.HironakaG1InitialZeroRecodedCentres.initial_zero_stepOne_literal_max_loci}
\leanok
For a coherent mark-one ideal on the full coordinate
carrier with empty boundary, the initial stabilized
Step~1 comparison at level zero reads the native general
value. When the ideal has maximal order, this is a
strict recoding of the literal lower-depth native
Step~1 value. The corresponding maximum loci are
exact images and preimages for the entire cosupport.
\end{theorem}
\begin{proof}
Read the actual prescribed zero contact of the
stabilization. On the native mark-one input the initial
order companion is the original ideal on its order mask.
Read its positive general branch and compose the two
strict constructors. Apply the zero section to both
cosupports and their maximum loci.
\end{proof}

\begin{theorem}[The supported section reader on an actual cylindrical history]
\label{hir:g1-native-section-reader-canonical}
\uses{hir:g1-later-literal-mixed-section-stratum,hir:g1-initial-zero-literal-maxima}
\lean{Hironaka.HironakaG1SectionReaderCanonical.sectionMarked_maximalOrder,Hironaka.HironakaG1SectionReaderCanonical.canonical_iff_of_section_reader,Hironaka.HironakaG1SectionReaderCanonical.Lift.original_native_section_reader}
\leanok
On the actual cylindrical tower, evaluate each prefix
using the original supported section reader. For maximal
order on the source, canonicality for this reader agrees
with original native canonicality on the actual lift.
The canonical endpoint values use the strictly monotone
dimension shift, with exact maximum images and preimages.
\end{theorem}
\begin{proof}
The section trace retains the source order bound.
At each canonical prefix the cylindrical and section
records have the same tests and cosupport. For the shared
section value function their maximum centres therefore
agree. Transport the supported section histories and
values along the actual native lift.
\end{proof}

\begin{theorem}[The actual native cut under one normal padding]
\label{hir:g1-proper-normal-native-cut}
\uses{hir:g1-common-graph-hs-cuts,hir:g1-closed-normal-weighted-presentation}
\lean{Hironaka.HironakaG1ProperNormalNativeCut.liftedCarrier_isAnalytic,Hironaka.HironakaG1ProperNormalNativeCut.liftedCarrier_cut_and_bound,Hironaka.HironakaG1ProperNormalNativeCut.properNormalMarked_native_cut_and_bound}
\leanok
For an actual native source on a coordinate open, its
zero-section carrier is analytic in the padded open.
The same intrinsic Hilbert--Samuel profile supplies its
ambient upper bound and top cut. That cut is exactly the
cosupport of the actual proper normal padding.
\end{theorem}
\begin{proof}
The section and projection are inverse intrinsic maps of
the analytic carriers. Their stalk isomorphisms determine
the Hilbert--Samuel profiles and the ambient cut. The
literal padded cosupport is the same lift of the original
native cosupport.
\end{proof}

\begin{theorem}[Different native marks in two actual ambient realizations]
\label{hir:g1-cross-ambient-native-class-comparison}
\uses{hir:g1-cross-ambient-native-comparison}
\lean{Hironaka.HironakaG1CrossAmbientNativeClassComparison.cross_ambient_native_comparison_of_class}
\leanok
The actual native representatives of the two ambient
histories may have different marks. A common local weak
test class, equal carrier dimensions and normalized
boundary exponents give equivalent original canonicality,
the same strict order on matched cosupport points, and
the same maximum loci in the common chart.
\end{theorem}
\begin{proof}
Use unequal-mark canonical comparison on the original
native history. The two-pair cut compares its actual
values at both points. Carry this strict order through
the two ambient dimension maps and their cosupport
and maximum-locus formulas.
\end{proof}

\begin{theorem}[Actual padded native states in finite normal codimension]
\label{hir:g1-actual-finite-normal-native-state}
\uses{hir:g1-proper-normal-native-cut,hir:g1-closed-normal-weighted-presentation}
\lean{Hironaka.HironakaG1ProperNormalNativeChart.adaptedIdeal_productFlat_normalPadding,Hironaka.HironakaG1ProperNormalNativeState.liftedCarrier_vanishing_mem_iff,Hironaka.HironakaG1ProperNormalNativeState.properNormalMarked_nativeHSState,Hironaka.HironakaG1IteratedProperNormalNativeState.iteratedNormalMarked_data,Hironaka.HironakaG1IteratedProperNormalNativeState.iteratedNormalMarked_nativeHSState,Hironaka.HironakaG1IteratedProperNormalNativeState.locWPEquiv_iteratedNormalMarked_of_native}
\leanok
An actual proper normal padding of a full native source
is a native state for its closed zero-section carrier and
original Hilbert--Samuel profile. This applies at every
step of an arbitrary finite normal padding. The actual
mark stays fixed and the entire ordered boundary pulls
back by the composite projection. Independent states
over the same intrinsic source cut therefore have the
same native weak test class on the iterated ambient.
\end{theorem}
\begin{proof}
Extend an actual adapted chart by the normal coordinate.
The uncompressed vanishing ideal and compressed marked
ideal both have exact whole-ideal identities in that
chart. Apply the historical-law extension with the
actual Hilbert--Samuel cut, and repeat at each normal
coordinate. Apply native weak test equivalence at the
matching terminal cut.
\end{proof}

\begin{theorem}[The literal original input for the section reader]
\label{hir:g1-literal-native-section-reader}
\uses{hir:g1-native-section-reader-canonical}
\lean{Hironaka.HironakaG1LiteralNativeSectionReader.section_native_input_record,Hironaka.HironakaG1LiteralNativeSectionReader.exists_literal_native_section_reader}
\leanok
An actual padded test tower produces the test tower of
the literal original source. Its whole cosupport is
identified by the actual endpoint maps, and its original
Step~1 canonicality is equivalent to the supported
section reader on the cylindrical history. The canonical
endpoint gives the dimension-shifted values and both
maximum-locus formulas.
\end{theorem}
\begin{proof}
The actual native test construction gives the geometric
lift. The section record in its source chart has the
literal original carrier, ideal, mark and ordered
boundary. Transport the original Step~1 invariant
through the tower isomorphism in both directions and
use the section reader comparison on the lift.
\end{proof}

\begin{theorem}[Original finite read at the localized next born cut]
\label{hir:g1-localized-next-born-finite-read}
\uses{hir:g1-localized-next-born-canonicality}
\lean{Hironaka.HironakaG1LocalizedLaterBornRead.localized_positive_stepOne_finite_read,Hironaka.HironakaG1LocalizedLaterBornRead.contact_recursive_strict_order_on_canonical_suffixes,Hironaka.HironakaG1LocalizedLaterBornRead.localized_born_witness_read}
\leanok
When the actual canonical positive companion run starts at
the localized cut, its finite value is the recursive read
on the actual nested selected record and restricted suffix.
This selected history is canonical and retains the
endpoint cosupport. On that suffix, canonicality for
the prescribed contact transfers to the supported cut,
and normalized exponents give the same strict order.
\end{theorem}
\begin{proof}
Use the actual full-run split and lower canonical data.
The nested restriction maps identify the literal nested
selected record. For the contact comparison, use the
normalized exponents and the original unequal-mark cut
theorem on the same suffix.
\end{proof}

\begin{theorem}[The precise native branch and maxima of the advanced-head blowup]
\label{hir:g1-advanced-head-native-monomial-maxima}
\uses{hir:g1-joint-example-actual-native-value-obstruction}
\lean{Hironaka.HironakaG1LaterAdvancedHeadNativeValueExample.joint_advanced_head_section_nonspecial,Hironaka.HironakaG1LaterAdvancedHeadNativeValueExample.joint_advanced_head_exact_native_monomial_read,Hironaka.HironakaG1LaterAdvancedHeadNativeValueExample.joint_advanced_head_endpoint_cosupport,Hironaka.HironakaG1LaterAdvancedHeadNativeValueExample.joint_advanced_head_full_max_loci}
\leanok
The literal native lift of the actual canonical blowup
takes the finite branch at level one. Its selected
record has nonmonomial order zero and the original
monomial value on its actual mask. Both full endpoint
maximum loci are singletons at the paired survivor.
\end{theorem}
\begin{proof}
A punctured strict-section curve proves nonspeciality
and the actual native chronology transfers it.
The born divisor gives the monomial branch at the
restart. The three genuine projective charts cover the
whole endpoint. On the two linear pivots the ideal is
a unit, and the remaining pivot has the unique survivor.
The native lift identifies that whole cosupport.
\end{proof}

\begin{theorem}[Independent actual nonlinear normal selectors]
\label{hir:g1-actual-nonlinear-native-selector-comparison}
\uses{hir:g1-actual-finite-normal-native-state,hir:g1-closed-normal-actual-selector}
\lean{Hironaka.HironakaG1NonlinearGraphNative.selectorAmbientMarked_native,Hironaka.HironakaG1NonlinearGraphNative.locWPEquiv_independent_selectorAmbient}
\leanok
An actual nonlinear selector graph carries the native
state with its literal normal power, original mark and
complete ordered boundary. Two independently chosen
original records with one added direction have the
same weak test class after their actual ambient chart
changes when the graph subsets and ordered boundary
lists agree there.
\end{theorem}
\begin{proof}
Transport the verified zero-section state through the
actual holomorphic graph shift. The whole ideal pulls
back with its original normal power. Replace the ordered
boundary by the complete prescribed list and apply the
native test comparison on the common chart cut.
\end{proof}

\begin{theorem}[The full zero graph and the iterated normal opens]
\label{hir:g1-finite-normal-block-geometry}
\uses{hir:g1-actual-finite-normal-native-state}
\lean{Hironaka.HironakaG1FiniteSourceWeakClassGraphLift.iteratedNormalOpen_eq_blockOpen,Hironaka.HironakaG1FiniteSourceWeakClassGraphLift.iteratedNormalProjection_blockEquiv,Hironaka.HironakaG1FiniteSourceWeakClassGraphLift.iteratedNormalCarrier_preimage_blockEquiv}
\leanok
The iterated normal opens are the actual full-block open,
with its literal base projection. Their closed iterated
carrier is the full zero graph of the original carrier.
\end{theorem}
\begin{proof}
Iterate the first-coordinate projections. The normal
coordinates split into the inherited block and the new
last coordinate. Their simultaneous vanishing cuts out
the literal zero graph on the full-block open.
\end{proof}

\begin{theorem}[Independent finite normal paddings on a common actual chart]
\label{hir:g1-independent-finite-normal-native-class}
\uses{hir:g1-actual-finite-normal-native-state,hir:g1-finite-normal-block-geometry}
\lean{Hironaka.HironakaG1IndependentFiniteNormalNative.normalMarked_boundary_on_chart,Hironaka.HironakaG1IndependentFiniteNormalNative.locWPEquiv_independent_finiteNormal,Hironaka.HironakaG1IndependentFiniteNormalNative.locWPEquiv_independent_finiteNormal_open}
\leanok
Two independent native source records with the same
intrinsic profile may have different ambient dimensions,
marks and padding codimensions. On an actual common
open chart whose maps pull back the same closed carrier
and full ordered boundary, their finite normal padded
records have the same local weak test class.
\end{theorem}
\begin{proof}
The finite padding constructs both actual native states.
Transport along the two open embeddings after dimension
identification. The pulled carriers and ordered
boundaries give the hypotheses of native weak test
equivalence on the common chart.
\end{proof}

\begin{theorem}[The literal weighted record of a finite nonlinear graph]
\label{hir:g1-finite-nonlinear-graph-native-state}
\uses{hir:g1-finite-normal-block-geometry,hir:g1-actual-nonlinear-native-selector-comparison}
\lean{Hironaka.HironakaG1IteratedNormalBlockIdeal.iteratedNormalMarked_ideal,Hironaka.HironakaG1IteratedNormalBlockIdeal.iteratedNormalMarked_pullback_blockEquiv,Hironaka.HironakaG1IteratedNormalBlockIdeal.blockAmbientMarked_pullback_shift,Hironaka.HironakaG1IteratedNormalBlockIdeal.blockAmbientMarked_nativeHSState}
\leanok
The actual iterated ideal has the base pullback and each
normal coordinate to the original mark. It gives the
whole finite-block record in the full coordinate chart.
An arbitrary holomorphic selector transports this record
to the actual nonlinear graph, with its original native
state, ordinary domain and full ordered boundary.
\end{theorem}
\begin{proof}
At each closed normal step, pull back the preceding
coherent generators and adjoin the power of the actual
new coordinate. The block dictionary identifies all
these coordinates on the full normal open. Then apply
the finite block chart shift to the whole marked record
and its native state.
\end{proof}

\begin{theorem}[Independent native originals on their actual intrinsic common graph]
\label{hir:g1-actual-common-graph-native-class}
\uses{hir:g1-finite-nonlinear-graph-native-state,hir:g1-independent-finite-normal-native-class,hir:g1-common-graph-hs-cuts}
\lean{Hironaka.HironakaG1NativeCoordinateWindow.nativeHSState_coordinateWindow,Hironaka.HironakaG1CommonGraphRestrictionMaps.commonGraphFirstRestriction_isOpenHoloEmbedding,Hironaka.HironakaG1CommonGraphRestrictionMaps.commonGraphSecondRestriction_isOpenHoloEmbedding,Hironaka.HironakaG1CommonGraphRestrictionMaps.exists_actual_commonBlockGraph_restrictions,Hironaka.HironakaG1ActualCommonGraphNativeComparison.commonGraphLeftMarked_E,Hironaka.HironakaG1ActualCommonGraphNativeComparison.commonGraphRightMarked_E,Hironaka.HironakaG1ActualCommonGraphNativeComparison.exists_actual_commonGraph_native_comparison}
\leanok
Two independent native originals with the same intrinsic
Hilbert--Samuel profile and a specified intrinsic
isomorphism admit an actual common weighted graph chart
at every original carrier point. Both whole records on
it have the native state of the same closed carrier.
Each retains its original mark, and its ordered
boundary pulls back along its own actual ambient map.
When those ordered lists agree, the two weighted records
have the same whole local weak test class.
\end{theorem}
\begin{proof}
The intrinsic isomorphism gives actual paired original
coordinate windows and the two holomorphic graph
extensions. Restrict the original native records to
these windows and form both finite nonlinear padded
records. The inclusion and the exchanged-coordinate
inclusion are actual open holomorphic embeddings.
They pull back the two nonlinear carriers to exactly
the same closed graph. The native weak class theorem
applies to the two transported states.
\end{proof}

\begin{theorem}[The actual advanced positive cut and its ordered born suffix]
\label{hir:g1-advanced-born-finite-read}
\uses{hir:g1-localized-next-born-finite-read}
\lean{Hironaka.HironakaG1AdvancedBornFiniteRead.retained_boundary_eq_old_filter_append_new,Hironaka.HironakaG1AdvancedBornFiniteRead.advanced_positive_stepOne_finite_read}
\leanok
The actual head of a canonical positive run supplies
the selected record on every smaller mask window through
its anchor. The restricted suffix is canonical there
and reads the whole finite value. Selection filters the
incident old slots and retains the complete born suffix
in its chronological order.
\end{theorem}
\begin{proof}
Use the original lower data at the actual chronological
split and restrict through the exact open diagram. In
the ordered SNC list a touching born slot cannot equal
a touching old slot, so the old-slot deletion leaves
the entire born suffix.
\end{proof}

\begin{theorem}[True order and strict trace of the localized next born record]
\label{hir:g1-localized-born-order-and-trace}
\uses{hir:g1-localized-next-born-canonicality}
\lean{Hironaka.HironakaG1LocalizedBornNormalOrder.exists_localized_born_normal_order,Hironaka.HironakaG1LocalizedBornNormalHistory.same_mark_test_order_and_strict_contact,Hironaka.HironakaG1LocalizedBornNormalHistory.exists_localized_born_normal_history}
\leanok
On a smaller actual window, the literal selected born
record contains the incoming normal to its original
mark. Every retained divisor exponent vanishes there.
On every actual tested continuation the nonmonomial
order at each survivor equals that mark. The contact
is the strict transform of the supported carrier, and
the full order clock and strict contact hold at each
surviving projected prefix.
\end{theorem}
\begin{proof}
Normalize the fields on the retained-slot intersections
using the actual supported transversality. The normal
power and the vanishing retained exponents identify the
same-mark companion with the carrier restriction.
Propagate the true order and the actual strict trace
along the same tested tower.
\end{proof}

\begin{theorem}[The supported special branch on the actual next born mask]
\label{hir:g1-localized-next-born-supported-special}
\uses{hir:g1-localized-next-born-canonicality}
\lean{Hironaka.HironakaG1LocalizedLaterBornSpecial.positive_special_member_no_domain,Hironaka.HironakaG1LocalizedLaterBornSpecial.exists_actual_localized_later_born_special_read,Hironaka.HironakaG1LocalizedLaterBornSpecial.born_selected_contact_no_domain_of_supported_special}
\leanok
At a supported special anchor of positive born level,
the original upper read has the exact special top at
that level and retains it under the actual localization.
The corresponding finite selected record on every open
through the anchor has no recursive domain. A padded
nonspecial contact on this supported branch has the
same domain obstruction.
\end{theorem}
\begin{proof}
Positive speciality reflects through literal open
pullback. It gives local saturation of the selected
carrier in the cosupport. The nonvanishing condition
of the ordinary domain excludes this saturation.
Transfer the saturation to any constructed contact
through the same-carrier local comparison.
\end{proof}

\begin{theorem}[Mixed highest positive strata and the actual native maxima]
\label{hir:g1-later-positive-mixed-centre-criterion}
\uses{hir:g1-later-literal-mixed-section-stratum,hir:g1-later-positive-prefix-history}
\lean{Hironaka.HironakaG1LaterPositiveMixedCentreCriterion.later_positive_mixed_max_eq_iff,Hironaka.HironakaG1LaterPositiveMixedCentreCriterion.later_positive_mixed_prefix_read_iff,Hironaka.HironakaG1LaterPositiveMixedCentreCriterion.later_positive_mixed_no_prefix_read_of_special_order_lt,Hironaka.HironakaG1LaterPositiveMixedCentreCriterion.later_positive_mixed_no_prefix_read_of_reversed_order,Hironaka.HironakaG1LaterPositiveMixedCentreCriterion.later_positive_mixed_prefix_read_of_unique_special_order}
\leanok
On the attained highest positive stratum of the mixed
case, the literal section maximizers are its special
anchors. With the actual native lift, exact raw maximum
readback holds precisely when the padded tuples tie on
those anchors and strictly dominate every finite
section anchor. A strict difference in true order
between special anchors or a strict reversal across
the branches obstructs readback.
\end{theorem}
\begin{proof}
Use the actual finite and top section reads and the
literal padded tuple order. A single top stratum
dominates every lower level. Compare both maximum sets
on the top stratum and carry the section set through
the supported native inclusion.
\end{proof}

\begin{theorem}[Actual boundary traces of independent native weighted records]
\label{hir:g1-native-raw-boundary-trace}
\uses{hir:g1-actual-common-graph-native-class,hir:g1-cross-ambient-native-class-comparison}
\lean{Hironaka.HironakaG1NativeBoundaryTrace.native_boundary_divides_once_iff,Hironaka.HironakaG1NativeBoundaryTrace.native_boundary_submark_of_trace_gap,Hironaka.HironakaG1NativeBoundaryTrace.normalized_native_boundary_exponents_of_trace,Hironaka.HironakaG1NativeBoundaryTrace.native_canonical_iff_of_boundary_trace}
\leanok
The raw weighted record is divisible by an original
boundary slot exactly when its local trace is inside the
intrinsic reduced carrier. A contained trace that escapes
the Hilbert--Samuel cut gives an exponent strictly
between zero and the mark. When every contained trace
stays in the cut, independent native marks with the same
ordered boundary give equal normalized exponents and
equivalent original canonicality.
\end{theorem}
\begin{proof}
Use the whole reduced-stalk dictionary for the actual
native presentations. Containment in the cut tests
the exponent at the full mark. The coefficient and
trace estimates then compare the two literal exponents
after cross-multiplication by their marks.
\end{proof}

\begin{theorem}[The exact numerical native read throughout the advanced-head example]
\label{hir:g1-advanced-head-exact-native-payload}
\uses{hir:g1-advanced-head-native-monomial-maxima}
\lean{Hironaka.HironakaG1LaterAdvancedHeadNativePayloadExample.joint_advanced_head_native_numerical_payload,Hironaka.HironakaG1LaterAdvancedHeadNativePayloadExample.joint_advanced_head_initial_exact_comparison,Hironaka.HironakaG1LaterAdvancedHeadNativePayloadExample.joint_advanced_head_nested_recoding_changes_on_history}
\leanok
The actual initial native read is special at level two.
The exact final native read has level one, order zero
and monomial key \((1,1,\{1\})\). The initial level-zero
insertion recodes the initial pair on this actual
history. All such insertions fail at the matched final
survivor.
\end{theorem}
\begin{proof}
The literal selected native record contains the actual
exceptional chart coordinate. Testing it along the native
curve determines exponent one in the born slot and
computes the monomial key. Compare this exact finite
letter with the already verified cylindrical value
at the same canonical endpoint.
\end{proof}

\begin{theorem}[The true-mark companion on the actual initial order mask]
\label{hir:g1-same-mark-initial-stepOne-read}
\uses{hir:g1-initial-zero-bare-contact-canonical}
\lean{Hironaka.HironakaG1ActualLaterOrderRun.original_positive_initial_order_companion,Hironaka.HironakaG1SameMarkInitialStepOne.same_mark_initial_stepOne_read,Hironaka.HironakaG1NestedBareContactTransport.nested_initial_bare_contact_data,Hironaka.HironakaG1NestedBareContactTransport.initial_order_zero_canonical_for_earlier_contact,Hironaka.HironakaG1NestedBareContactTransport.initial_order_prescribed_contact_at_tested_head}
\leanok
When a positive order run spans the entire original canonical
tower, its recursive letter reads the companion at the initial
head. When the true-mark companion is the carrier restriction,
the same letter reads original Step~1 on the literal marked
restriction. On a zero phase starting there, canonicality
transfers to the earlier bare coefficient record on the exact
nested mask. The prescribed contacts are also
restricted-equivalent at any tested advancing head.
\end{theorem}
\begin{proof}
Identify the actual order split with the initial head and
transport restricted equivalence on its order mask. The
coefficient construction respects restricted equivalence
and open pullback. Apply the initial zero-phase theorem
on the actual companion and transfer along the literal
restricted tower with the same coefficient mark.
\end{proof}

\begin{theorem}[The carrier test for disagreement of the next-born special letters]
\label{hir:g1-born-special-carrier-cutoff}
\uses{hir:g1-localized-next-born-supported-special}
\lean{Hironaka.HironakaG1BornSpecialCarrierCriterion.positive_special_carrier_criterion,Hironaka.HironakaG1BornSpecialCarrierCriterion.exists_actual_born_special_carrier_criterion,Hironaka.HironakaG1BornSpecialCarrierCriterion.positive_special_letters_restrict,Hironaka.HironakaG1BornSpecialCarrierCriterion.positive_special_run_cutoff_iff}
\leanok
At a positive level on a common literal test, the padded
anchor is special exactly when the supported anchor is
special and the points outside the common cosupport on the
larger selected carrier do not accumulate at the anchor.
Both tests respect open restriction and the whole ordered
boundary. The supported nonspecial run is strictly shorter
exactly when the two tests disagree at a projected predecessor
inside the padded run.
\end{theorem}
\begin{proof}
The exact selected carriers are nested and the transformed bar
cosupports agree. Saturation of the larger carrier is the
conjunction of saturation of the smaller carrier and absence
of accumulation of the exterior points. Apply this
decomposition at each projected predecessor in the run.
\end{proof}

\begin{theorem}[Retained boundary counts for the realized mixed maximum comparison]
\label{hir:g1-later-mixed-centre-boundary-trace}
\uses{hir:g1-later-positive-mixed-centre-criterion}
\lean{Hironaka.HironakaG1LaterMixedCentreBoundaryTrace.laterPositiveRead_order_eq_fixed_mark,Hironaka.HironakaG1LaterMixedCentreBoundaryTrace.laterPositiveRead_companion_boundary_trace,Hironaka.HironakaG1LaterMixedCentreBoundaryTrace.later_positive_mixed_max_eq_iff_companions,Hironaka.HironakaG1LaterMixedCentreBoundaryTrace.later_positive_mixed_max_eq_of_boundary_trace,Hironaka.HironakaG1LaterMixedCentreBoundaryTrace.later_positive_mixed_no_max_eq_of_boundary_mismatch}
\leanok
On a canonical padded prefix, all realized padded finite
letters in the highest positive section stratum read the
same selected order. The comparison of the two maximum
loci then uses the actual companion letters. The companion
level counts the retained ordered members through the
survivor and its top letter tests actual companion
specialness. Unequal retained counts among section special
anchors obstruct equality of the full maximum loci.
\end{theorem}
\begin{proof}
The selected raw mark is invariant under the actual
canonical splitting and localization. Identify the
realized companion's original members with the strict
transforms of its retained list. Substitute their counts
into the mixed maximum criterion and use the lexicographic
order of the companion letters.
\end{proof}

\begin{theorem}[Native boundary exponents in local codimension at least two]
\label{hir:g1-native-codimension-two-traces}
\uses{hir:g1-native-raw-boundary-trace}
\lean{Hironaka.HironakaG1NativeBoundaryCodimension.locDimGE_of_regularAtDim_eventually_subset,Hironaka.HironakaG1NativeBoundaryCodimension.native_boundary_exponent_zero_of_local_codimension_two,Hironaka.HironakaG1NativeBoundaryCodimension.independent_native_boundary_exponents_zero_of_local_codimension_two,Hironaka.HironakaG1NativeBoundaryCodimension.independent_native_canonical_iff_of_local_codimension_two}
\leanok
Intrinsic local codimension at least two along the common
Hilbert--Samuel cut forces zero raw exponent in every
ordered SNC slot of either independent native presentation.
Original canonicality agrees at their own marks along the
same tested history under the common graph comparison.
\end{theorem}
\begin{proof}
Containment of the germ of a regular boundary hypersurface
would force the native analytic set to have local dimension
at least that of the hypersurface. The codimension bound
rules out containment. The raw boundary trace theorem
then gives equal normalized exponents and canonicality.
\end{proof}

\begin{theorem}[The original prescribed contact at the actual advancing zero split]
\label{hir:g1-advancing-order-zero-contact}
\uses{hir:g1-same-mark-initial-stepOne-read}
\lean{Hironaka.HironakaG1AdvancingOrderZeroContact.order_companion_actual_zero_head_contact}
\leanok
The nonspecial zero letter of the actual initial same-mark
companion reads the earlier prescribed contact at the actual
zero split. The suffix is canonical for that contact and its
endpoint survives. The contact retains the strict zero trace,
the recursive mark, and the whole ordered boundary at the head.
\end{theorem}
\begin{proof}
Extract the canonical tail and its exact finite value using the
prescribed tangent on the companion. The genuine tested head
permits comparison with the original prescribed contact.
Restricted equivalence transports both canonicality and the value
on the literal suffix.
\end{proof}

\begin{theorem}[Comparing independent native histories on a common graph refinement]
\label{hir:g1-native-graph-overlap-read}
\uses{hir:g1-actual-common-graph-native-class,hir:g1-native-raw-boundary-trace}
\lean{Hironaka.HironakaG1NativeGraphOverlapRead.native_graph_overlap_canonical_and_max,Hironaka.HironakaG1NativeGraphOverlapRead.native_graph_trace_on_refinement,Hironaka.HironakaG1NativeGraphOverlapRead.native_graph_overlap_descend_to_charts,Hironaka.HironakaG1NativeGraphOverlapRead.exists_refinement_of_ordered_incident_germs,Hironaka.HironakaG1NativeGraphOverlapRead.exists_native_graph_refinement_canonical_and_max}
\leanok
Matching ordered incident boundary germs give a common open
graph refinement. When contained boundary traces stay inside
the intrinsic cut, original canonicality, endpoint cosupports
and maximum loci correspond under the actual chart tower
isomorphisms. The original charts may carry different slots that
restrict to the empty set.
\end{theorem}
\begin{proof}
Shrink until the corresponding incident germs agree as sets
and the other listed members are empty. Compare normalized
exponents after removing the empty slots on the common graph.
Transport canonicality and maxima back to the two raw records
and descend along the supplied tower isomorphisms.
\end{proof}

\begin{theorem}[The complete joint advanced-head history]
\label{hir:g1-advanced-head-terminal-example}
\uses{hir:g1-advanced-head-exact-native-payload}
\lean{Hironaka.HironakaG1LaterAdvancedHeadTerminalExample.joint_advanced_head_padded_complete_history,Hironaka.HironakaG1LaterAdvancedHeadTerminalExample.joint_advanced_head_common_complete_history}
\leanok
In the actual joint example the second point blowups occur at
the unique padded and native survivors of the first step.
They extend the same first lift and both histories are canonical.
The two stages with nonempty cosupport have their exact four raw
letters, and both terminal cosupports are empty.
\end{theorem}
\begin{proof}
The three middle-chart generators give the whole padded point
ideal at its surviving anchor. Blow up each singleton maximum and
use the actual lifted native point blowup. In each terminal chart
a controlled generator is a unit.
\end{proof}

\begin{theorem}[The full positive letter through an advancing zero contact]
\label{hir:g1-true-mark-zero-contact-read}
\uses{hir:g1-advancing-order-zero-contact,hir:g1-same-mark-initial-stepOne-read}
\lean{Hironaka.HironakaG1ActualLaterOrderRun.order_run_of_all_test_orders,Hironaka.HironakaG1TrueMarkZeroContactRead.true_mark_original_advancing_zero_read,Hironaka.HironakaG1TrueMarkZeroContactRead.original_positive_advancing_zero_read_on}
\leanok
On an actual positive selected subwindow with the same-mark
normal geometry, the canonical lower history runs at its
true mark from the source. When the actual companion ends
in a finite zero letter, the original upper Step~1 letter
uses the entire recursive value of the prescribed zero contact
on its actual suffix. The contact is canonical there with
its surviving endpoint and current ordered boundary.
\end{theorem}
\begin{proof}
The same-mark tangent forces the true order on every tested
prefix, so the order companion lies at the initial cut.
Recover its canonical history and apply the advancing zero-contact
comparison. Substitute the exact resulting recursive value into
the finite read at the original positive head.
\end{proof}

\begin{theorem}[A native crossing with a strictly submark boundary exponent]
\label{hir:g1-native-crossing-submark}
\uses{hir:g1-native-raw-boundary-trace,hir:g1-native-codimension-two-traces}
\lean{Hironaka.HironakaG1NativeCrossingGeometry.crossing_origin_singular,Hironaka.HironakaG1NativeCrossingGeometry.crossing_hs_drop_on_axis,Hironaka.HironakaG1AnchoredBoundaryTraceWitness.exists_anchored_submark_boundary,Hironaka.HironakaG1NativeCrossingSubmark.exists_actual_native_crossing_submark}
\leanok
The two coordinate axes in the complex plane, with the
first axis as the single ordered SNC member, have an actual
anchored native presentation in the invariant domain.
The boundary trace lies in the analytic crossing and leaves
the Hilbert--Samuel cut arbitrarily close to the origin.
The raw native boundary exponent at the anchor is strictly
positive and strictly below its original mark.
\end{theorem}
\begin{proof}
Both branches are regular away from their intersection, and
their distinct tangent directions give singularity at the
origin. The Hilbert--Samuel value on the punctured boundary
axis differs from the value at the origin. Use the actual
anchored native phase and its original-function point laws
to apply the boundary trace criterion on its open chart.
\end{proof}

\begin{theorem}[The born selected comparison on smaller common windows]
\label{hir:g1-localized-born-window-restriction}
\uses{hir:g1-localized-next-born-canonicality,hir:g1-true-mark-zero-contact-read}
\lean{Hironaka.HironakaG1LocalizedBornNormalHistory.positive_members_zero_trace,Hironaka.HironakaG1LocalizedBornWindowRestriction.selected_bare_contact_pullback_val,Hironaka.HironakaG1LocalizedBornWindowRestriction.localized_born_window_data_restrict,Hironaka.HironakaG1LocalizedBornWindowRestriction.localized_born_witness_inside}
\leanok
The actual next-born selected pair restricts to any smaller
common window, retaining its exact zero trace and recursive
coefficient record. The restricted pair has the same ordered
boundary, both domains, normalized exponent comparison, and
equivalent canonicality. In particular, the window can be
chosen inside any specified neighborhood of the split anchor.
\end{theorem}
\begin{proof}
Member selection commutes with restriction of its fixed
ambient data, and the actual coefficient ideal commutes with
open pullback. Transport the local test class to the smaller
window and compare the normalized exponents on the literal
slots. Intersect with the given neighborhood to choose a window.
\end{proof}

\begin{theorem}[The true order and the supported contact on one born window]
\label{hir:g1-common-born-zero-contact-read}
\uses{hir:g1-localized-born-window-restriction,hir:g1-true-mark-zero-contact-read}
\lean{Hironaka.HironakaG1LocalizedBornNormalOrder.exists_localized_born_normal_order_subwindow,Hironaka.HironakaG1LocalizedBornZeroContactRead.exists_localized_born_zero_contact_read,Hironaka.HironakaG1AdvancingSupportedBornContact.supported_contact_local_class,Hironaka.HironakaG1AdvancingSupportedBornContact.advancing_kept_contact_supported_stage,Hironaka.HironakaG1AdvancingSupportedBornContact.advancing_kept_contact_supported_canonical,Hironaka.HironakaG1AdvancingSupportedBornContact.actual_zero_head_advances_supported_born}
\leanok
The prescribed-to-supported comparison and the true-order normal
geometry hold on one actual subwindow of the original positive phase.
At a finite zero endpoint the original positive value reads a contact
on the strict transform of the supported selected carrier. The
supported record advances along the same actual head and suffix,
with the entire stage boundary and its original mark. Its suffix is
canonical when the advanced supported record has the cosupport gap.
\end{theorem}
\begin{proof}
Choose the true-order window inside the comparison window and restrict
the whole contact pair. Transfer the selected zero trace to the order
mask and its actual head. The supported local test class and both
domains follow along the tested prefix. The gap then provides the
normalized exponents for canonicality on the actual suffix.
\end{proof}

\begin{theorem}[The supported gap on source opens avoided by the centres]
\label{hir:g1-supported-gap-empty-epoch}
\uses{hir:g1-common-born-zero-contact-read}
\lean{Hironaka.HironakaG1SupportedGapEmptyEpoch.cosupport_gap_pullback_val,Hironaka.HironakaG1SupportedGapEmptyEpoch.cosupport_gap_geoTransformAlong_empty,Hironaka.HironakaG1SupportedGapEmptyEpoch.cosupport_gap_geoTransformAlong_centresAvoid,Hironaka.HironakaG1SupportedGapEmptyEpoch.cosupport_gap_on_stage_open_of_centresAvoid,Hironaka.HironakaG1SupportedGapEmptyEpoch.advancing_supported_canonical_of_centresAvoid}
\leanok
The supported cosupport gap persists through an empty
epoch and on any actual stage open lying over a source window
avoided by all centres of the prefix. On that stage open, the
prescribed contact and advanced supported record have equivalent
canonicality along the actual restricted suffix.
\end{theorem}
\begin{proof}
On the empty restriction the controlled and total pullbacks agree
at both the original mark and mark one. Transport the source gap
to the stage mask and use it for the normalized exponent comparison.
\end{proof}

\begin{theorem}[The count of born slots in a restarted inner companion]
\label{hir:g1-restarted-inner-companion-born-count}
\uses{hir:g1-true-mark-zero-contact-read}
\lean{Hironaka.HironakaG1NestedPositiveBirthObstruction.restarted_positive_companion_level_eq_born_count,Hironaka.HironakaG1NestedPositiveBirthObstruction.restarted_positive_inner_level_requires_birth,Hironaka.HironakaG1NestedPositiveBirthObstruction.restarted_positive_inner_level_mem_exceptional_support,Hironaka.HironakaG1NestedPositiveBirthObstruction.restarted_positive_companion_level_zero_of_isEmptyTower,Hironaka.HironakaG1NestedPositiveBirthObstruction.actual_positive_successor_read_inner_zero_of_isEmptyTower,Hironaka.HironakaG1NestedPositiveBirthObstruction.actual_positive_successor_read_ne_positive_inner_of_isEmptyTower}
\leanok
On the actual new companion mask the inner level counts the
incident born members retained at the positive selection.
A recursive history with only empty centres forces inner level
zero and excludes a positive inner letter at the true order.
\end{theorem}
\begin{proof}
The selected list drops every incident original member.
The exact companion pullback preserves the remaining list
and the incidence count on the actual mask. A born member is
in the exceptional support, which is empty for an empty history.
\end{proof}

\begin{theorem}[Attained recursive orders across independent native charts]
\label{hir:g1-native-chart-attained-value-order}
\uses{hir:g1-native-graph-overlap-read}
\lean{Hironaka.HironakaG1NativeChartValueOrder.attainedValueOrderIso_apply,Hironaka.HironakaG1NativeChartValueOrder.nativeDimensionImageOrderIso_apply,Hironaka.HironakaG1NativeChartValueOrder.exists_original_empty_boundary_value_embedding,Hironaka.HironakaG1NativeChartValueOrder.cross_ambient_native_chart_value_read,Hironaka.HironakaG1NativeChartValueOrder.native_graph_overlap_attained_order}
\leanok
The recursive values on each actual ambient lift are the exact
dimension shifts of its native chart values. Native normalized
comparison gives an isomorphism between the attained orders.
Literal empty boundary slots contribute their own insertion maps,
which also reconstruct the original orders on common graph refinements.
\end{theorem}
\begin{proof}
Transport both chart histories through their actual native
comparison. Compose the order embeddings for their dimensions
and their empty slots. Common strict order cuts identify the
attained ranges and give the pointwise isomorphism.
\end{proof}

\begin{theorem}[The normalized exponent of a principal native stalk]
\label{hir:g1-native-principal-boundary-ratio}
\uses{hir:g1-native-crossing-submark}
\lean{Hironaka.HironakaG1NativePrincipalBoundaryRatio.native_uniformPrincipal_boundary_ratio,Hironaka.HironakaG1NativePrincipalBoundaryRatio.analyticPrincipal_initial_vertices,Hironaka.HironakaG1NativePrincipalBoundaryRatio.native_uniformDegree_original_cut,Hironaka.HironakaG1NativePrincipalBoundaryRatio.independent_native_uniformPrincipal_boundary_ratio}
\leanok
For a nonzero principal analytic stalk with uniform inherited
native degree $\mu$, a literal native compression at mark $d$
has boundary exponent $e_{\rm native}$ satisfying
$\mu e_{\rm native}=d e_{\rm analytic}$. Two independent uniform
presentations of the same principal stalk have the same normalized
exponent. The canonical initial Taylor diagram has one vertex.
\end{theorem}
\begin{proof}
The whole uniform weighted compression at the principal stalk
is the corresponding power of the original ideal. Taking its
divisor exponent gives the scaled formula. Both inherited degrees
are the same by the next-cut exclusion at the original stalk.
The initial exponents of a principal Taylor ideal form the
translate of the positive orthant by the generator's initial exponent.
\end{proof}

\begin{theorem}[Recognizing the finite zero branch from the strict supported trace]
\label{hir:g1-companion-zero-branch-recognition}
\uses{hir:g1-common-born-zero-contact-read}
\lean{Hironaka.HironakaG1CompanionZeroBranchRecognition.codim_one_absent_of_strict_supported_trace,Hironaka.HironakaG1CompanionZeroBranchRecognition.zero_step_regime_of_strict_supported_trace,Hironaka.HironakaG1CompanionZeroBranchRecognition.step_regime_zero_iff_level_of_strict_supported_trace,Hironaka.HironakaG1CompanionZeroBranchRecognition.full_zero_clock_iff_no_source_incidence,Hironaka.HironakaG1CompanionZeroBranchRecognition.exists_actual_localized_companion_zero_recognition}
\leanok
The actual strict supported trace rules out codimension-one
specialness on a common tested continuation. At a zero endpoint
this supplies nonspecialness. On a nonempty history the zero
run reaches the source exactly when the source projection
misses all retained original members. In the actual localized
companion it suffices to miss the inherited born members.
\end{theorem}
\begin{proof}
The common cosupport has smaller local dimension than the
strict supported carrier, since its supported ideal is nonzero
along that carrier. This excludes saturation in codimension
one of the larger carrier. The original member count can only
drop on the history. Apply this to all projected predecessors
and note that positive selection deletes the incident old slots.
\end{proof}

\begin{theorem}[The mixed maximum count at the actual upper head]
\label{hir:g1-later-mixed-upper-born-trace}
\uses{hir:g1-later-mixed-centre-boundary-trace}
\lean{Hironaka.HironakaG1LaterMixedUpperBornTrace.laterPositiveRead_companion_level_eq_upperHeadNonOriginal,Hironaka.HironakaG1LaterMixedUpperBornTrace.upperHeadNonOriginalBoundaryCount_eq_zero_of_full_run,Hironaka.HironakaG1LaterMixedUpperBornTrace.later_positive_mixed_max_eq_of_shared_initial_head}
\leanok
The realized companion level counts the strict transforms
of the non-original upper members that are already present
at the actual head and meet the endpoint. When all the upper
runs start at the original source, this count is zero and
agreement of the companion-special and section-special loci
gives equality of the mixed maximum loci.
\end{theorem}
\begin{proof}
Read the actual initial companion of the positive lower
history. Pull its ordered members through both actual masks.
All original members still in the list missed the head
anchor, so their strict transforms miss the endpoint. At a
shared initial head there are no non-original upper members.
\end{proof}

\begin{theorem}[The restarted second companion at level zero]
\label{hir:g1-restarted-second-companion-zero}
\lean{Hironaka.HironakaG1RestartedSecondCompanionZero.restarted_second_true_mark_companion_zero}
\leanok
Under the canonical history and strict supported trace hypotheses, the restarted second true-mark companion is in the finite zero branch.
\end{theorem}
\begin{proof}
Restrict to the actual born chart and apply the common tested zero-trace criterion.
\end{proof}

\begin{theorem}[Supported exponents on an advancing contact]
\label{hir:g1-advancing-supported-contact-exponents}
\lean{Hironaka.HironakaG1AdvancingSupportedContactExponents.inherited_supported_exponent_at_least_mark,Hironaka.HironakaG1AdvancingSupportedContactExponents.advancing_supported_canonical_of_new_trace}
\leanok
The inherited boundary exponent is at least the mark under the stated hypotheses. A new trace supplying the exponent comparison transports canonicality across the advancing supported contact.
\end{theorem}
\begin{proof}
Use the original lower old-exponent bound and compare the normalized exponents on the selected history.
\end{proof}

\begin{theorem}[Canonicality from the controlled born trace]
\label{hir:g1-advancing-born-controlled-trace}
\lean{Hironaka.HironakaG1AdvancingBornControlledTrace.advancing_supported_canonical_of_born_controlled_trace}
\leanok
The controlled born trace supplies the exponent comparison for canonicality of the advancing supported contact.
\end{theorem}
\begin{proof}
Transport the scaled exponent along the trace and apply the advancing-contact criterion.
\end{proof}

\begin{theorem}[A native two-axis crossing at mark two]
\label{hir:g1-native-crossing-ideal}
\lean{Hironaka.HironakaG1NativeCrossingIdeal.axisProductGerm_ne_zero,Hironaka.HironakaG1NativeCrossingIdeal.adapted_origin_axis,Hironaka.HironakaG1NativeCrossingIdeal.adapted_origin_crossing,Hironaka.HironakaG1NativeCrossingIdeal.axisProductGerm_exact_order,Hironaka.HironakaG1NativeCrossingIdeal.crossing_axis_exponent_one,Hironaka.HironakaG1NativeCrossingIdeal.crossing_adapted_canonical_mark_two}
\leanok
The product of the two coordinate germs has order two. Both adapted boundary axes have exponent one. The actual canonical marked input realizes this crossing.
\end{theorem}
\begin{proof}
Compute the orders in adapted coordinates and identify the initial canonical marked input.
\end{proof}

\begin{theorem}[The exact initial crossing ratio]
\label{hir:g1-native-crossing-exact-ratio}
\lean{Hironaka.HironakaG1NativeCrossingExactRatio.exists_actual_crossing_initial_ratio}
\leanok
The actual initial native crossing has boundary exponent one at mark two, giving ratio one half.
\end{theorem}
\begin{proof}
Compute the exponent from the reduced axis-product ideal in the adapted initial record.
\end{proof}

\begin{theorem}[Paired saturation and the mixed maximum]
\label{hir:g1-later-mixed-paired-saturation}
\lean{Hironaka.HironakaG1LaterMixedPairedSaturation.PairedSaturation.codimOne_iff_member,Hironaka.HironakaG1LaterMixedPairedSaturation.paired_saturation_special_iff,Hironaka.HironakaG1LaterMixedPairedSaturation.laterPositiveRead_leaf_top_iff_of_paired_saturation,Hironaka.HironakaG1LaterMixedPairedSaturation.later_positive_mixed_max_eq_of_paired_saturation}
\leanok
Matching contact and selected member germs and local cosupport thresholds identifies the special branch of the two readers. On a shared initial upper run, this gives equality of their mixed maximum loci.
\end{theorem}
\begin{proof}
Compare local saturation through the two actual open restrictions. The shared run gives zero companion incidence.
\end{proof}

\begin{theorem}[Actual codimension-two chart transport]
\label{hir:g1-codimension-two-chart-transport}
\lean{Hironaka.HironakaG1CodimensionTwoChartTransport.native_good_trace_of_codimension_two,Hironaka.HironakaG1CodimensionTwoChartTransport.native_codimension_two_zero_normalized,Hironaka.HironakaG1CodimensionTwoChartTransport.exists_codimension_two_native_refinement,Hironaka.HironakaG1CodimensionTwoChartTransport.actual_common_graph_domains,Hironaka.HironakaG1CodimensionTwoChartTransport.actual_common_graph_dimLE,Hironaka.HironakaG1CodimensionTwoChartTransport.exists_actual_common_graph_codimension_two,Hironaka.HironakaG1CodimensionTwoChartTransport.codimension_two_native_chart_transport}
\leanok
A codimension-two trace supplies the required native refinement. The common graph construction gives the domain and dimension bounds used for native chart transport.
\end{theorem}
\begin{proof}
Use the local zero-normalized trace and transport the chart data on the refined native opens.
\end{proof}

\begin{theorem}[An actual first born chart]
\label{hir:g1-recursive-born-finite-example}
\lean{Hironaka.HironakaG1RecursiveBornFiniteExample.initial_cosupport,Hironaka.HironakaG1RecursiveBornFiniteExample.canonical,Hironaka.HironakaG1RecursiveBornFiniteExample.quotientGen_mem,Hironaka.HironakaG1RecursiveBornFiniteExample.controlled_span,Hironaka.HironakaG1RecursiveBornFiniteExample.domain,Hironaka.HironakaG1RecursiveBornFiniteExample.test,Hironaka.HironakaG1RecursiveBornFiniteExample.stage_carrier,Hironaka.HironakaG1RecursiveBornFiniteExample.stage_span,Hironaka.HironakaG1RecursiveBornFiniteExample.stage_cosupport_iff,Hironaka.HironakaG1RecursiveBornFiniteExample.survivor}
\leanok
The example has a canonical first blowup and explicit controlled generators on its chart. The computation identifies the transformed carrier and cosupport and shows that the distinguished point survives.
\end{theorem}
\begin{proof}
Divide the selected generators on the actual chart and compute the controlled ideal and cosupport from their span.
\end{proof}

\begin{theorem}[The localized born companion on its full zero run]
\label{hir:g1-localized-born-full-zero-read}
\uses{hir:g1-restarted-second-companion-zero}
\lean{Hironaka.HironakaG1LocalizedBornFullZeroRead.exists_localized_born_full_zero_read}
\leanok
The localized born construction supplies its finite zero branch
and the canonical supported suffix on the actual companion run.
\end{theorem}
\begin{proof}
Use the restarted true-mark companion to obtain the zero regime.
Transport canonicality across the preserved supported gap at
the empty zero head and read the actual contact suffix.
\end{proof}

\begin{theorem}[An actual finite positive born companion]
\label{hir:g1-recursive-born-order-example}
\uses{hir:g1-recursive-born-finite-example,hir:g1-restarted-inner-companion-born-count}
\lean{Hironaka.HironakaG1RecursiveBornOrderExample.stage_ordinary_order,Hironaka.HironakaG1RecursiveBornOrderExample.stage_exceptional_exponent,Hironaka.HironakaG1RecursiveBornOrderExample.stage_nonmonomial_order,Hironaka.HironakaG1RecursiveBornOrderExample.order_run_restarts,Hironaka.HironakaG1RecursiveBornOrderExample.actual_finite_positive_companion,Hironaka.HironakaG1RecursiveBornSelectedExample.actual_selected_born_read}
\leanok
For the canonical point blowup of the mark-one ideal
$(x^2,y^3,t^2)$, the middle-chart survivor has ordinary order two,
born divisor exponent one, and nonmonomial order one. Its actual
order run restarts there, and the companion has a finite level-one
letter. The lower value is the original invariant of the actual
selected record on the born trace, with mark one and empty boundary.
\end{theorem}
\begin{proof}
Use the exact middle-chart controlled generators. Their order
and their quotient along the born divisor determine the two orders
and the exponent. The order drops at the step. At the new cut the
born divisor contributes one to the companion count, and punctured
points of its trace rule out saturation. Apply the original finite
read on its literal mask.
\end{proof}

\begin{theorem}[The native crossing on independent initial charts and powers]
\label{hir:g1-native-crossing-independent-histories}
\uses{hir:g1-native-crossing-exact-ratio}
\lean{Hironaka.HironakaG1CrossingUnequalPowerHistory.powered_history_boundary,Hironaka.HironakaG1CrossingUnequalPowerHistory.normalized_powered_history_boundary,Hironaka.HironakaG1CrossingUnequalPowerHistory.powered_history_canonical_order,Hironaka.HironakaG1CrossingUnequalPowerHistory.exists_crossing_unequal_power_histories,Hironaka.HironakaG1IndependentInitialCrossing.exists_crossing_initial_ratio_in_centered_chart,Hironaka.HironakaG1IndependentInitialCrossing.crossing_initial_common_window_eq,Hironaka.HironakaG1IndependentInitialCrossing.exists_independent_crossing_initial_history_comparison}
\leanok
Two independent initial presentations of the two-axis crossing
have mark two and axis exponent one. Their entire pulled marked
records agree on a supplied common chart. Independently chosen
positive coherent powers have unequal marks and the same normalized
exponents on their actual common controlled histories. They retain
canonicality, maxima, and the order of their attained recursive values.
\end{theorem}
\begin{proof}
The actual initial adapted charts produce the same whole
crossing ideal and the same literal ordered axis on the overlap.
Run the invariant on the resulting equal records. For the
positive powers, use exact scaling of the controlled ideals
and divisor exponents and the original marked-power comparison.
\end{proof}

\begin{theorem}[Actual comparison with one positive normal block]
\label{hir:g1-positive-graph-boundary-comparison}
\uses{hir:g1-codimension-two-chart-transport,hir:g1-native-chart-attained-value-order}
\lean{Hironaka.HironakaG1SingularGraphCodimension.native_cut_singular,Hironaka.HironakaG1SingularGraphCodimension.native_cut_locDimLE_of_smooth_envelope,Hironaka.HironakaG1SingularGraphCodimension.native_cut_empty_of_ambient_dim_le_one,Hironaka.HironakaG1SingularGraphCodimension.actual_common_graph_codim_two_of_positive_right,Hironaka.HironakaG1PositiveGraphBoundaryComparison.exists_positive_actual_graph_boundary_comparison}
\leanok
Inside a smooth graph envelope the actual singular native cut
has local dimension strictly smaller than the envelope. A single
positive complementary normal block therefore suffices to rule
out contained boundary traces. For independently padded native
records, matching ordered incident boundary germs then give a
common refinement with equivalent canonicality, matching
cosupport and maxima, and isomorphic attained recursive orders.
\end{theorem}
\begin{proof}
The singularity of a cut point excludes local saturation in the
smooth graph envelope. Apply the resulting dimension bound to
each regular boundary hypersurface. Intersect the windows where
the incident germs coincide, normalize the boundary exponents
on the actual graph, and use original class comparison.
\end{proof}

\begin{theorem}[Holomorphic unit representatives on the actual native graph]
\label{hir:g1-relative-boundary-units-on-graph}
\uses{hir:g1-snc-reader-full-label-normalization,hir:g1-positive-graph-boundary-comparison}
\lean{Hironaka.HironakaG1RelativeBoundaryUnitsOnGraph.ambient_unit_representative,Hironaka.HironakaG1RelativeBoundaryUnitsOnGraph.original_germ_relation_on_graph,Hironaka.HironakaG1RelativeBoundaryUnitsOnGraph.simultaneous_unit_representatives,Hironaka.HironakaG1RelativeBoundaryUnitsOnGraph.labelledFunction_quotient,Hironaka.HironakaG1RelativeBoundaryUnitsOnGraph.exists_relative_labelled_units_on_graph,Hironaka.HironakaG1FixedReaderBoundaryUnits.labelled_span_of_member,Hironaka.HironakaG1FixedReaderBoundaryUnits.source_labelled_trace_span,Hironaka.HironakaG1FixedReaderBoundaryUnits.target_labelled_trace_span,Hironaka.HironakaG1FixedReaderBoundaryUnits.fixed_reader_units_on_graph,Hironaka.HironakaG1FixedReaderBoundaryUnits.source_labelledFunction_eq_normalized,Hironaka.HironakaG1FixedReaderBoundaryUnits.target_labelledFunction_eq_normalized,Hironaka.HironakaG1FixedReaderBoundaryUnits.active_reader_units_on_graph}
\leanok
In fixed independent normalized SNC readers, each original
labelled native trace has its coordinate-generated whole
quotient ideal. If the actual intrinsic isomorphism respects
the structure-point map and identifies these labelled ideals,
there are simultaneous nonvanishing ambient holomorphic unit
representatives on both opens. On the actual paired native
graph their values agree, and each active target coordinate is
the corresponding source coordinate times that common unit.
\end{theorem}
\begin{proof}
The fixed reader membership dictionary identifies the entire
divisor ideal with the coordinate ideal, including the unit
ideal for an absent slot. Apply local-ring unit equivalence
of generators and lift the units through the ambient stalk
maps. Represent the finite family on simultaneous opens and
shrink so every representative stays nonzero. Equality of
the intrinsic germs spreads to the paired native graph.
The structure-point map gives the same pointwise values there.
\end{proof}

\begin{theorem}[A common ratio window with the original ordered boundaries]
\label{hir:g1-actual-relative-reader-window}
\uses{hir:g1-relative-boundary-actual-projections,hir:g1-relative-boundary-units-on-graph}
\lean{Hironaka.HironakaG1RelativeReaderWindow.ratio_anchor_mem_readerUnitWindow,Hironaka.HironakaG1RelativeReaderWindow.reader_maps_at_ratio_anchor,Hironaka.HironakaG1RelativeReaderWindow.reader_maps_slot_preimage_eq,Hironaka.HironakaG1RelativeReaderWindow.reader_maps_ordered_boundary_eq,Hironaka.HironakaG1RelativeReaderWindow.reader_maps_contMDiff,Hironaka.HironakaG1RelativeReaderWindow.exists_readerUnitWindow_of_intrinsic_iso,Hironaka.HironakaG1RelativeReaderWindow.exists_reader_lift_of_coordinate_ratio}
\leanok
The common nonzero-ratio open of the two actual independent
SNC readers maps holomorphically to both original ambient opens
and contains a lift of their actual anchors. Its two inverse
images of each corresponding original boundary slot coincide.
Position-preserving labels give equality of the whole ordered
lists. Any actual pair of reader points with the compatible
active coordinate ratios lifts with precisely those ratios.
\end{theorem}
\begin{proof}
Intersect the nonzero-ratio locus with both inverse images of
the original reader coordinate opens. Join each active block
to its complementary block and apply the inverse actual reader.
For a paired point place the first reader's active block, the
prescribed ratios, and the two complementary blocks in the
product. The ratio identity gives the second reader's active
block. For every slot the common selector is either absent or
gives one active index. Multiplication by a nonzero ratio
preserves its zero set, giving equality slot by slot.
\end{proof}

\begin{theorem}[The actual paired native graph in the common ratio window]
\label{hir:g1-actual-paired-ratio-graph}
\uses{hir:g1-actual-relative-reader-window,hir:g1-relative-boundary-units-on-graph}
\lean{Hironaka.HironakaG1FixedReaderComplementExtensions.exists_holomorphic_ambient_graph_extension,Hironaka.HironakaG1FixedReaderComplementExtensions.simultaneous_holomorphic_graph_extensions,Hironaka.HironakaG1FixedReaderComplementExtensions.paired_ambient_subwindows,Hironaka.HironakaG1FixedReaderComplementExtensions.paired_reader_complement_extensions,Hironaka.HironakaG1ActualPairedRatioGraph.exists_actual_paired_ratio_graph}
\leanok
The paired native traces lift to the common ratio window of
the two original readers. Both projections are bijective on
the lifted graph and induce the actual given intrinsic pairing.
There the two holomorphic ratio blocks agree, and the complementary
coordinate functions extend holomorphically to the two ambient
opens in their respective fixed readers.
\end{theorem}
\begin{proof}
Lift each paired complementary coordinate through the actual
quotient stalk and choose simultaneous holomorphic representatives.
Shrink the two ambient opens so their native traces correspond
under the given isomorphism and its inverse. Lift every paired
point to the ratio window using its active coordinate identity.
The first projection recovers the source point and the second
recovers its image. The two inverse identities for the intrinsic
isomorphism give bijectivity onto the paired traces.
\end{proof}

\begin{theorem}[Both original reader maps are actual Euclidean submersions]
\label{hir:g1-actual-relative-euclidean-submersions}
\uses{hir:g1-actual-relative-reader-window,hir:g1-actual-paired-ratio-graph}
\lean{Hironaka.HironakaG1ReaderProductSubmersions.window_product_isSubmersion,Hironaka.HironakaG1ReaderProductSubmersions.leftReader_isSubmersion,Hironaka.HironakaG1ReaderProductSubmersions.rightReader_isSubmersion,Hironaka.HironakaG1ReaderProductSubmersions.product_submersion_reparameterize,Hironaka.HironakaG1ReaderProductSubmersions.firstReaderMap_isSubmersionOfComplement,Hironaka.HironakaG1ReaderProductSubmersions.secondReaderMap_isSubmersionOfComplement,Hironaka.HironakaG1ReaderProductSubmersions.euclidean_reader_maps_submersions,Hironaka.HironakaG1ReaderProductSubmersions.euclidean_reader_maps_slot_preimage_eq}
\leanok
On a single Euclidean parameterization of the actual ratio
window, the original reader maps are holomorphic submersions
with their actual independent complementary dimensions.
Their inverse images of corresponding boundary slots agree
on that same Euclidean window.
\end{theorem}
\begin{proof}
Reorder the left active and complementary coordinates as
the first factor of the source product. For the right reader,
scale by the nonzero ratio variables and then reorder the
right active and complementary coordinates. Both actual
local charts therefore make the reader maps product
projections. Carry the window and the slot equality through
the common Euclidean linear change of coordinates.
\end{proof}

\begin{theorem}[The two exact native states on the actual common analytic graph]
\label{hir:g1-intrinsic-ratio-native-states}
\uses{hir:g1-actual-relative-euclidean-submersions,hir:g1-actual-paired-ratio-graph}
\lean{Hironaka.HironakaG1RatioGraphProductEquations.paired_ratio_graph_iff_product,Hironaka.HironakaG1RatioGraphProductEquations.leftGraphValue_hol,Hironaka.HironakaG1RatioGraphProductEquations.analytic_ratio_graph_of_equation,Hironaka.HironakaG1RatioGraphProductEquations.exists_actual_paired_ratio_graph_equations,Hironaka.HironakaG1RatioBlockGraphAdapters.leftBlockMap_base,Hironaka.HironakaG1RatioBlockGraphAdapters.rightBlockMap_base,Hironaka.HironakaG1RatioBlockGraphAdapters.blockGraph_preimages,Hironaka.HironakaG1RatioBlockGraphAdapters.leftRatioMarked_mark,Hironaka.HironakaG1RatioBlockGraphAdapters.rightRatioMarked_mark,Hironaka.HironakaG1RatioBlockGraphAdapters.leftRatioMarked_boundary,Hironaka.HironakaG1RatioBlockGraphAdapters.rightRatioMarked_boundary,Hironaka.HironakaG1RatioBlockGraphAdapters.ratioMarked_ordered_boundary_eq,Hironaka.HironakaG1RatioBlockGraphAdapters.leftBlockMap_hol,Hironaka.HironakaG1RatioBlockGraphAdapters.rightBlockMap_hol,Hironaka.HironakaG1RatioBlockGraphAdapters.leftBlockMap_openEmbedding,Hironaka.HironakaG1RatioBlockGraphAdapters.rightBlockMap_openEmbedding,Hironaka.HironakaG1RatioBlockGraphAdapters.leftBlockGraphValue_hol,Hironaka.HironakaG1RatioBlockGraphAdapters.rightBlockGraphValue_hol,Hironaka.HironakaG1RatioBlockNativeStates.nativeState_markedPullback_openImage,Hironaka.HironakaG1RatioBlockNativeStates.blockGraph_preimage_codim_two,Hironaka.HironakaG1RatioBlockNativeStates.normalizedBaseMap_openImage,Hironaka.HironakaG1RatioBlockNativeStates.leftBlockMap_openImage,Hironaka.HironakaG1RatioBlockNativeStates.rightBlockMap_openImage,Hironaka.HironakaG1RatioBlockNativeStates.ratioMarked_nativeStates,Hironaka.HironakaG1IntrinsicRatioNativeComparison.exists_intrinsic_ratio_native_comparison}
\leanok
An actual intrinsic isomorphism respecting the original ordered
boundary trace ideals gives two literal padded native records
on a common analytic ratio graph. Each record preserves its
own original mark, ordinary invariant domain, and full ordered
boundary. The positive normal complement gives local codimension
at least two at the points of the common cut.
\end{theorem}
\begin{proof}
The exact paired graph is cut out in each product by its native
carrier and the holomorphic complementary equation. In the second
product use the actual ratio scaling. Both product maps restrict
to open embeddings. Construct the normal block marked records and
pull back along these embeddings. Their native states, domains, and
common support then follow from their literal block graphs. The
normal envelope and its singular native cut give the local dimension
bound. The ratio identity identifies every original boundary slot.
\end{proof}

\begin{theorem}[Canonical comparison on the intrinsic ratio graph]
\label{hir:g1-intrinsic-ratio-canonical-comparison}
\uses{hir:g1-intrinsic-ratio-native-states,hir:g1-codimension-two-chart-transport}
\lean{Hironaka.HironakaG1IntrinsicRatioCanonical.exists_intrinsic_ratio_canonical_comparison,Hironaka.HironakaG1RatioBlockUnpadding.markedPullback_blockAmbientMarked_section,Hironaka.HironakaG1RatioBlockUnpadding.exists_openImage_section_unpadding}
\leanok
The two produced native records on the actual common ratio graph
have equivalent original canonicality along every common tower.
For a canonical tower their endpoint cosupports and maximum loci
coincide and their attained recursive values are order-isomorphic.
The holomorphic section of each actual block graph recovers the
whole original normalized source record on the literal section open.
\end{theorem}
\begin{proof}
The common native state and ordered boundary give local weak test
equivalence. Local codimension two forces zero initial exponents,
so the normalized boundary condition follows at the original
independent marks. Apply the original canonical two-pair
comparison to the actual common tower and to pairs of its
active points. Their strict comparisons define the attained
order isomorphism and also identify both maximum loci. On the
section every added normal graph equation is zero, so its whole
marked pullback is the original normalized record.
\end{proof}

\begin{theorem}[An actual blowup of the native graph section]
\label{hir:g1-native-ratio-section-blowup}
\uses{hir:g1-intrinsic-ratio-canonical-comparison,hir:g1-intrinsic-ratio-native-states}
\lean{Hironaka.HironakaG1RatioGraphCentreGeometry.blockGraph_isTransverse,Hironaka.HironakaG1RatioGraphCentreGeometry.blockAmbientMarked_cosupp,Hironaka.HironakaG1RatioGraphCentreGeometry.liftedCarrier_hasNormalCrossings,Hironaka.HironakaG1RatioGraphCentreGeometry.blockGraph_hasNormalCrossings,Hironaka.HironakaG1RatioGraphCentreGeometry.blockGraph_admissibleCentre,Hironaka.HironakaG1RatioGraphCentreGeometry.ratioSection_centreGeometry,Hironaka.HironakaG1RatioSectionBlowup.exists_native_and_padded_section_blowup,Hironaka.HironakaG1RatioSectionBlowup.nativePullback_withN_inter,Hironaka.HironakaG1RatioSectionBlowup.sectionSupport_closed,Hironaka.HironakaG1RatioSectionBlowup.sectionSupport_regular,Hironaka.HironakaG1RatioSectionBlowup.blockGraph_supported_isWellFormed,Hironaka.HironakaG1RatioSectionBlowup.exists_blockGraph_section_blowup,Hironaka.HironakaG1RatioGraphNativeSource.sectionBaseEquiv_centre_image,Hironaka.HironakaG1RatioGraphNativeSource.sectionBaseEquiv_record,Hironaka.HironakaG1RatioGraphNativeSource.exists_ratioSection_source_blowup}
\leanok
The padded cosupport is the graph over the native cosupport.
The whole graph section and the graph of any native admissible
centre have the required compatibility with the pulled ordered
boundary. On a literal block chart the paired genuine native
and padded blowup projections commute. The whole transformed
section record is the native transform, and the actual padded
transform is restricted-equivalent on the strict supported
carrier with all inherited and born slots.
\end{theorem}
\begin{proof}
The added normal graph powers cut out the actual graph at
the original mark. Simultaneous coordinates for the native
centre and boundary pull back along the block projection and
the graph shift. Form the induced native blowup of the graph
section. The exact controlled restriction formula identifies
the section record with the native transform. Its compatibility
with the original ambient control keeps the mark and every
ordered boundary slot on the resulting supported carrier.
\end{proof}

\begin{theorem}[Mixed maxima with advancing head traces]
\label{hir:g1-mixed-maxima-advancing-head-traces}
\uses{hir:g1-later-mixed-paired-saturation,hir:g1-later-mixed-upper-born-trace}
\lean{Hironaka.HironakaG1NonemptyHeadMixedTrace.pairedSaturation_of_tested_source_trace,Hironaka.HironakaG1NonemptyHeadMixedTrace.HasTestedPairedSourceTrace.to_paired,Hironaka.HironakaG1NonemptyHeadMixedTrace.upperHeadNonOriginalBoundaryCount_eq_zero_iff_off_endpoint,Hironaka.HironakaG1NonemptyHeadMixedTrace.upperHeadNonOriginalBoundaryCount_eq_zero_of_off_head,Hironaka.HironakaG1NonemptyHeadMixedTrace.later_positive_mixed_max_eq_of_zero_head_counts,Hironaka.HironakaG1NonemptyHeadMixedTrace.later_positive_mixed_max_eq_of_tested_zero_birth_traces,Hironaka.HironakaG1MixedSelectedThresholdTrace.selected_threshold_of_source_window_class,Hironaka.HironakaG1MixedSelectedThresholdTrace.codimOne_iff_member_of_selected_threshold,Hironaka.HironakaG1MixedSelectedThresholdTrace.codimOne_iff_member_of_tested_selected_trace,Hironaka.HironakaG1MixedSelectedThresholdTrace.zero_companion_special_iff_of_selected_source_class}
\leanok
The mixed comparison allows distinct advancing upper heads
when all strict transforms of their non-original members miss
the relevant endpoints. The actual tested source tangent gives
the endpoint tangent. Under the source-window class comparison
the selected cosupport threshold follows from the actual test.
Matching strict selected carrier germs then identifies specialness.
\end{theorem}
\begin{proof}
Propagate the tangent along the actual ordinary test and read
the selected cosupport using the retained source incidences.
Compare saturation on the matched carrier germs. The literal
endpoints have zero non-original count, so the mixed maximum
criterion reduces to the resulting specialness comparison.
\end{proof}

\begin{theorem}[Specialness and actual retained counts on mixed normal strata]
\label{hir:g1-mixed-tested-one-chart-special}
\uses{hir:g1-mixed-maxima-advancing-head-traces,hir:g1-prescribed-companion-selected-trace}
\lean{Hironaka.HironakaG1MixedOneChartSpecial.selected_zero_special_iff_of_tested_local_source,Hironaka.HironakaG1MixedOneChartSpecial.actual_initial_order_zero_special_iff_of_source_class,Hironaka.HironakaG1MixedOneChartSpecial.later_mixed_padded_max_subset_max_head_count,Hironaka.HironakaG1MixedOneChartSpecial.later_mixed_max_ne_of_upper_head_birth_mismatch}
\leanok
Along the actual tested initial companion run, the local
source-window class identifies specialness in the selected
zero-level companion and in the section. On a realized later
highest stratum with fixed true order and nonspecial padded
points, a padded maximum must have maximal retained upper-head
birth count. The two-anchor count discrepancy in the mixed
criterion forces different maximum loci at a section-maximal
anchor.
\end{theorem}
\begin{proof}
Transport the prescribed local source tangent along the twice
restricted ordinary test. The source class identifies the
selected cosupport threshold, and the transported tangent
identifies the carrier germ needed in the specialness criterion.
On the later stratum, substitute the actual retained count into
the companion value. Strict increase in that count forces strict
increase in the full padded value when the true order is fixed.
\end{proof}

\begin{theorem}[Retained boundary count and live leaf determine the mixed maxima]
\label{hir:g1-later-mixed-live-centres}
\uses{hir:g1-mixed-tested-one-chart-special}
\lean{Hironaka.HironakaG1LaterMixedLiveCentres.later_padded_max_eq_max_leaf_on_max_head_count,Hironaka.HironakaG1LaterMixedLiveCentres.later_padded_max_eq_zero_special_of_occupied,Hironaka.HironakaG1LaterMixedLiveCentres.later_mixed_max_eq_iff_max_count_and_live_leaf}
\leanok
On the realized nonspecial padded highest stratum of fixed
true order, the padded maxima first maximize the actual
retained head count and then the live recursive leaf among
equal counts. With an occupied section-special maximal anchor,
the entire mixed maxima agree precisely when all section-special
anchors attain the maximum count and tie its leaf, and every
section-finite anchor at that count has a strictly smaller leaf.
At maximal count zero an occupied, recognized special branch
identifies the padded maxima with the section-special locus.
\end{theorem}
\begin{proof}
Read each original value from its fixed true order, actual
retained count and recursive leaf. Compare counts first.
Within the maximal count compare the live leaves literally.
The section-special anchor already maximizes the section value,
so equality of entire maxima gives exactly the stated equalities
and strict inequalities. In the zero case the recognized special
branch is the top recursive leaf.
\end{proof}

\begin{theorem}[The finite live leaf on an actual common chronological chart]
\label{hir:g1-chronological-mixed-leaf-charts}
\uses{hir:g1-later-mixed-live-centres,hir:g1-chronological-companion-suffix}
\lean{Hironaka.HironakaG1ChronologicalLeafTies.ChronologicalChart.read,Hironaka.HironakaG1ChronologicalLeafTies.chronological_chart_recursive_tie,Hironaka.HironakaG1ChronologicalLeafTies.positive_companion_leaf_at_actual_selected_suffix,Hironaka.HironakaG1ChronologicalLeafTies.realized_companion_of_positive_order_witness,Hironaka.HironakaG1ChronologicalLeafTies.finite_companion_leaves_tie_on_chronological_chart,Hironaka.HironakaG1ChronologicalLeafTies.compAt_d_eq_of_true_mark,Hironaka.HironakaG1ChronologicalLeafTies.laterPositiveRead_companion_on_actual_order_mask,Hironaka.HironakaG1ChronologicalLeafTies.laterPositiveRead_actual_companion_mark,Hironaka.HironakaG1ChronologicalMixedMaxima.common_chronological_chart_of_overlapping_support,Hironaka.HironakaG1ChronologicalMixedMaxima.SelectedLeafRead.sameChart_of_matching_records,Hironaka.HironakaG1ChronologicalMixedMaxima.SelectedLeafRead.eq_of_sameChart,Hironaka.HironakaG1ChronologicalMixedMaxima.selectedLeafRead_of_finite_positive_letter,Hironaka.HironakaG1ChronologicalMixedMaxima.later_positive_selectedLeafRead_of_finite,Hironaka.HironakaG1ChronologicalMixedMaxima.later_mixed_max_eq_of_selected_chronological_charts,Hironaka.HironakaG1ChronologicalMixedMaxima.later_mixed_max_ne_of_tied_count_finite_section_top_companion,Hironaka.HironakaG1ChronologicalMixedMaxima.later_mixed_max_ne_of_tied_count_cross_branch_sameChart}
\leanok
Finite live leaves at a positive actual retained count tie on
common selected chronological refinements with compatible
local class, carrier dimension, marks and normalized exponents.
Actual companion marks agree for reads of the same proper-normal
prefix. Overlapping masks of the same supported suffix have
a common refinement with literally identical whole records.
The resulting leaf equality applies to the exact mixed maximum
criterion. A tied-count top leaf or a tied finite leaf at a
section-finite anchor obstructs maximum agreement.
\end{theorem}
\begin{proof}
Read each finite positive leaf on its literal next selected
suffix. Transport the records and values by the actual nested
open restrictions and tower isomorphisms. When the two records
come from overlapping masks, both pulls are the same restricted
whole record. Otherwise the stated local class and normalized
boundary compatibility compare their attained values. First
compare retained counts, then apply the exact leaf criterion on
the maximal-count stratum.
\end{proof}

\begin{theorem}[Actual survivor symmetries and the mixed maximum]
\label{hir:g1-actual-mixed-survivor-symmetry}
\uses{hir:g1-chronological-mixed-leaf-charts}
\lean{Hironaka.HironakaG1ChronologicalSurvivorSymmetry.common_chronological_chart_of_tower_symmetry,Hironaka.HironakaG1ChronologicalSurvivorSymmetry.recursive_tie_of_canonical_survivor_symmetry,Hironaka.HironakaG1ChronologicalSurvivorSymmetry.canonical_distinct_max_anchors_of_survivor_symmetry,Hironaka.HironakaG1MixedMaskSymmetry.positive_mask_preimage_of_towerIso,Hironaka.HironakaG1MixedMaskSymmetry.MixedTowerSymmetry.preserves_section_and_upper_run,Hironaka.HironakaG1MixedMaskSymmetry.MixedTowerSymmetry.actual_positive_masks,Hironaka.HironakaG1MixedMaskSymmetry.MixedTowerSymmetry.companion_eq,Hironaka.HironakaG1MixedMaskSymmetry.later_mixed_max_eq_of_section_special_orbit}
\leanok
A tower symmetry preserving the two full original marked records
preserves the literal positive masks, section-special branch,
and actual chronological values at corresponding survivors.
The mixed maximum loci agree if the highest section-special
anchors form one such orbit and every section-finite anchor at
that level has strictly smaller retained count.
\end{theorem}
\begin{proof}
Pull back the actual ordered boundaries and cosupports along
the tower isomorphism. This identifies both actual run cuts and
the open masks, so the selected histories and recursive values
are transported literally. The orbit ties the section-special
values. The retained-count inequality already excludes every
section-finite competitor in the exact mixed maximum criterion.
\end{proof}

\begin{theorem}[Unequal retained counts on an actual common canonical prefix]
\label{hir:g1-common-canonical-unequal-born-counts}
\uses{hir:g1-actual-mixed-survivor-symmetry,hir:g1-later-mixed-live-centres}
\lean{Hironaka.HironakaG1AdvancedHeadCountExhaustion.advanced_first_stage_every_survivor_retains_one,Hironaka.HironakaG1AdvancedHeadCountExhaustion.advanced_two_occupied_stages_exhaust_retained_competitors,Hironaka.HironakaG1TwoZeroBornCounts.common_canonical_prefix_has_unequal_retained_counts}
\leanok
The source ideal $(x_0,x_1^2(x_1-1))$ has two actual initial
cosupport points of levels two and one. Its padded and section
records have the same first canonical point blowup. Two padded
survivors at level one have retained upper-head counts one and
zero, respectively.
\end{theorem}
\begin{proof}
The unique initial point at level two gives the common centre.
On the born middle chart the controlled ideal inclusion gives
the surviving point and its exceptional incidence. The second
point survives off the centre by the actual projection isomorphism.
The literal level runs then retain the one born slot at the first
survivor and no born slot at the second.
\end{proof}

\begin{theorem}[The two actual section windows of the unequal-count example]
\label{hir:g1-two-zero-actual-source-charts}
\uses{hir:g1-common-canonical-unequal-born-counts}
\lean{Hironaka.HironakaG1TwoZeroSourceCharts.born_source_window_ideal_eq,Hironaka.HironakaG1TwoZeroSourceCharts.off_source_window_ideal_eq,Hironaka.HironakaG1TwoZeroSourceCharts.positive_initial_selected_raw_order_of_cosupport,Hironaka.HironakaG1TwoZeroSourceCharts.section_window_record_eq,Hironaka.HironakaG1TwoZeroSourceCharts.off_source_section_nonspecial,Hironaka.HironakaG1TwoZeroSourceCharts.born_anchor_mem_section_window,Hironaka.HironakaG1TwoZeroSourceCharts.off_anchor_mem_section_window,Hironaka.HironakaG1TwoZeroSourceCharts.born_section_stage_chart_eq}
\leanok
Near the born source point the two-zero ideal has the literal
double-zero record. Near the off-centre source point it has the
literal simple-zero record, whose section is nonspecial there.
On the actual born endpoint window the transformed section
record agrees with the double-zero model. The initial positive
selection of a cosupport point has positive raw selected order.
\end{theorem}
\begin{proof}
Invert the factor nonzero on each specified source open.
The resulting literal ideal equalities preserve the marks and
the full ordered boundary on the restricted section records.
Apply actual controlled open restriction on the born window.
For the positive selection the actual incident boundary is
removed, and the selected cosupport then gives positive raw
ordinary order.
\end{proof}

\begin{theorem}[Actual unequal maximum loci for one Cartier source]
\label{hir:g1-cartier-two-anchors-maximum-discrepancy}
\uses{hir:g1-two-zero-actual-source-charts,hir:g1-common-canonical-unequal-born-counts,hir:g1-later-mixed-live-centres}
\lean{Hironaka.HironakaG1CartierTwoAnchors.record_domain,Hironaka.HironakaG1CartierTwoAnchors.record_cosupport,Hironaka.HironakaG1CartierTwoAnchors.padded_cosupport,Hironaka.HironakaG1CartierTwoAnchors.section_cosupport,Hironaka.HironakaG1CartierTwoAnchors.level_two_iff_origin,Hironaka.HironakaG1CartierTwoAnchors.cartier_common_first_canonical,Hironaka.HironakaG1CartierTwoAnchors.cartier_born_and_off_survive,Hironaka.HironakaG1CartierTwoAnchors.cartier_endpoint_cosupport,Hironaka.HironakaG1CartierTwoAnchors.cartier_section_endpoint_all_special,Hironaka.HironakaG1CartierTwoAnchors.cartier_section_maximum_eq_cosupport,Hironaka.HironakaG1CartierTwoAnchors.cartier_both_anchors_maximize_section,Hironaka.HironakaG1CartierTwoAnchors.cartier_born_count,Hironaka.HironakaG1CartierTwoAnchors.cartier_off_count,Hironaka.HironakaG1CartierTwoAnchors.cartier_common_prefix_maximum_discrepancy}
\leanok
For the same Cartier source $(x_0)$ with ordered boundary
$x_0=0,x_1=0$, the literal first point blowup is canonical
for both the padded and section records. The born and off-centre
survivors both attain the section maximum. Their respective
retained upper-head counts are one and zero, and the padded
value strictly ranks the born survivor higher. The actual
two maximum loci are different.
\end{theorem}
\begin{proof}
The origin is the unique initial point at level two for either
record. Its blowup leaves the actual strict source line as the
cosupport. The section values there are the special top of
level one. The born survivor starts its upper head at the
new chart and retains the exceptional slot. The off-centre
survivor retains its source run and has no new slot. In the
padded comparison the ordinary order is common, so strict
comparison of these counts strictly compares its full values.
\end{proof}

\begin{theorem}[The prescribed section germ at the actual selected endpoint]
\label{hir:g1-prescribed-companion-selected-trace}
\uses{hir:g1-mixed-maxima-advancing-head-traces,hir:g1-companion-zero-branch-recognition}
\lean{Hironaka.HironakaG1MixedPrescribedSelectedTrace.selectedCarrier_strictAlong,Hironaka.HironakaG1MixedPrescribedSelectedTrace.strictAlong_twice_restrict_preimage,Hironaka.HironakaG1MixedPrescribedSelectedTrace.bar_carrier_subset_of_source_subset,Hironaka.HironakaG1MixedPrescribedSelectedTrace.source_zero_eq_selected_of_carrier_and_zero,Hironaka.HironakaG1MixedPrescribedSelectedTrace.initial_order_companion_carrier_selected,Hironaka.HironakaG1MixedPrescribedSelectedTrace.initial_order_zero_companion_special_iff_of_prescribed_source,Hironaka.HironakaG1ActualPrescribedCompanionTrace.strict_prescribed_zero_eq_section_bar_carrier,Hironaka.HironakaG1ActualPrescribedCompanionTrace.exists_actual_initial_companion_prescribed_zero,Hironaka.HironakaG1ActualPrescribedCompanionTrace.strictAlong_eventually_eq_of_source_window,Hironaka.HironakaG1ActualPrescribedCompanionTrace.local_prescribed_zero_transports_selected_trace,Hironaka.HironakaG1ActualPrescribedCompanionTrace.no_strict_supported_domain_at_saturated_zero}
\leanok
The actual initial order companion has a neighborhood where
the prescribed zero-set cuts out the selected section carrier.
On the actual tested suffix over that neighborhood, retained
original level identifies the strict selected germ at the endpoint.
The source-window class then supplies its specialness comparison.
At saturation the finite supported construction has no invariant
domain on the strict contact.
\end{theorem}
\begin{proof}
Take the genuine local tangent at the initial order mask.
Identify its zero-set with the literal section using the
transformed bar carriers. The strict transform commutes
with both literal restrictions. The equality of the retained
level identifies the selected germ. Saturation on the resulting
carrier contradicts the nonzero condition of the finite domain.
\end{proof}

\begin{theorem}[The next selected ideal on the actual born trace]
\label{hir:g1-recursive-born-residual-example}
\uses{hir:g1-recursive-born-order-example}
\lean{Hironaka.HironakaG1RecursiveBornResidualExample.residual_span,Hironaka.HironakaG1RecursiveBornResidualExample.selected_born_trace_gens,Hironaka.HironakaG1RecursiveBornResidualExample.selected_born_payload_first}
\leanok
The actual selected ideal along the born trace in the finite
positive example has generators $(a^2,c^2)$. Its first recursive
entry is two in the actual finite letter.
\end{theorem}
\begin{proof}
Divide by the literal born power to compute the colon
ideal with generators $(a^2,y,c^2)$. Restrict it through the
actual selection to $y=0$. Restriction of a squared coordinate
to its axis rules out order three and computes the exact first
recursive entry.
\end{proof}

\begin{theorem}[The original upper letter on the actual common companion window]
\label{hir:g1-original-upper-born-full-zero}
\uses{hir:g1-localized-born-full-zero-read,hir:g1-prescribed-companion-selected-trace}
\lean{Hironaka.HironakaG1ActualCompanionWindowRead.step_one_open_window_history,Hironaka.HironakaG1ActualCompanionWindowRead.original_upper_born_full_zero_read}
\leanok
The actual initial order companion of an original canonical
positive upper history restricts to the genuine common tangent
window with its canonicality, value and phase clock intact.
If its own positive nonspecial run spans that window's history
and the supported anchors are nonspecial, the actual contact
and supported records have canonical tails and surviving
endpoints. The original outer letter reads that contact at the
exact true marks of both positive phases.
\end{theorem}
\begin{proof}
Read the initial order companion from the original upper letter.
Construct the actual supported tangent window at the same anchor.
Use open restriction of Step~1 and its clock. The restarted
positive phase then gives the full zero branch and both canonical
contact tails on their exact masks. Substitute their finite value
back into the two original positive letters.
\end{proof}

\begin{theorem}[The order-two companion and its prescribed trace in the example]
\label{hir:g1-recursive-born-order-two-companion}
\uses{hir:g1-recursive-born-residual-example}
\lean{Hironaka.HironakaG1RecursiveBornOrderTwoCompanionExample.selected_born_order_two_companion,Hironaka.HironakaG1RecursiveBornOrderTwoCompanionExample.selected_order_two_coordinate_trace}
\leanok
The actual selected order-two record in the finite positive
example has its original maximal-order companion at mark two
and level zero. Along the prescribed $a=0$ section of the
literal order mask its raw ideal is generated by $c^2$ at the
anchor.
\end{theorem}
\begin{proof}
Read the positive order-two successor on the actual mask.
The empty boundary and true-order mark identify the companion
as the carrier restriction. Restrict the squared coordinate
generators along $a=0$ to obtain the principal raw trace.
\end{proof}

\begin{theorem}[Exact order after the next coefficient operation]
\label{hir:g1-recursive-born-order-two-coefficient}
\uses{hir:g1-recursive-born-order-two-companion}
\lean{Hironaka.HironakaG1RecursiveBornOrderTwoBarExample.homogenized_mark_two,Hironaka.HironakaG1RecursiveBornOrderTwoBarExample.coefficient_mark_two,Hironaka.HironakaG1RecursiveBornOrderTwoBarExample.original_bar_and_contact_mark_two,Hironaka.HironakaG1RecursiveBornOrderTwoBarExample.actual_recursive_born_contact_marks,Hironaka.HironakaG1RecursiveBornOrderTwoBarExample.actual_recursive_born_coordinate_coefficient_bounds,Hironaka.HironakaG1RecursiveBornOrderTwoBarExample.actual_recursive_born_coordinate_coefficient_exact_order}
\leanok
In the finite positive born example, the actual bar and the next
coefficient operation keep mark two. On the coordinate section
$a=0$ inside the three literal masks, the next coefficient ideal
has order two at the distinguished point.
\end{theorem}
\begin{proof}
At mark two, homogenization and each coefficient operation use
the square of the actual first derivative ideal. The restricted
ideal therefore still has order at least two and contains the
square of the third coordinate in its span. The third-coordinate
axis in the prescribed section pulls that square back to $t^2$.
The axis lifts through all three open masks at the distinguished
point, so order at least three would force $t^2$ to vanish to
order three there, a contradiction.
\end{proof}

\begin{theorem}[The prescribed coordinate as the actual canonical contact]
\label{hir:g1-recursive-born-order-two-canonical-tangent}
\uses{hir:g1-recursive-born-order-two-coefficient}
\lean{Hironaka.HironakaG1RecursiveBornOrderTwoTangentExample.ambient_coordinate_tangent,Hironaka.HironakaG1RecursiveBornOrderTwoTangentExample.actual_order_two_coordinate_tangent,Hironaka.HironakaG1RecursiveBornOrderTwoTangentExample.actual_order_two_canonical_coordinate_read}
\leanok
The prescribed first-coordinate section in the order-two
companion is an admissible contact. The survivor is nonspecial
and the ordered companion boundary is empty. The original
canonical value reads the whole actual contact record on a
literal open in the third mask, giving the original composed
successor value with the last curve entry read on that section.
\end{theorem}
\begin{proof}
The first-coordinate field preserves the born fibre and
differentiates the square of the first coordinate to supply the
tangent. Transport this field through the actual open masks.
The cosupport is locally a point on the two-dimensional carrier,
which makes the survivor nonspecial. Remove the closed bad locus
and use the canonical zero-level contact formula on the resulting
open. Substitute that contact value into both earlier letters.
\end{proof}

\begin{theorem}[The full numerical value on the actual positive born chart]
\label{hir:g1-recursive-born-numerical-payload}
\uses{hir:g1-recursive-born-order-two-canonical-tangent,hir:g1-recursive-born-order-two-coefficient}
\lean{Hironaka.HironakaG1RecursiveBornNumericalPayloadExample.order_two_special_at_true_mark_two,Hironaka.HironakaG1RecursiveBornNumericalPayloadExample.actual_recursive_born_numerical_payload}
\leanok
The canonical survivor of the actual finite positive born example
has its full original value. Its successive order and level pairs
are $(1,1)$, $(2,0)$ and $(2,0)$. The last pair is the special
curve letter at true mark two.
\end{theorem}
\begin{proof}
The coefficient ideal keeps its true mark and has exact order
two on the prescribed actual contact. The square of the final
coordinate is in the restricted span, so its vanishing makes
the cosupport locally the point on the one-dimensional contact
carrier. This is the special zero-level branch. Substitute its
top leaf into the exact canonical contact formula and then into
the two preceding original successor letters.
\end{proof}

\begin{theorem}[Canonical restriction at an advancing positive companion head]
\label{hir:g1-chronological-companion-suffix}
\uses{hir:g1-original-upper-born-full-zero,hir:g1-mixed-maxima-advancing-head-traces}
\lean{Hironaka.HironakaG1ChronologicalCompanionSuffix.positive_companion_chronological_suffix,Hironaka.HironakaG1ChronologicalCompanionSuffix.positive_companion_chronological_tangent,Hironaka.HironakaG1ChronologicalCompanionSuffix.exists_companion_source_bar_tangent_window}
\leanok
An arbitrary actual positive nonspecial companion run
has a tested head and a canonical literal selected suffix.
Its boundary preserves the chronological distinction between
original and born slots. A genuine inherited source tangent
with the supported zero trace gives the actual tangent and
contact class at this head. Restriction to that tangent
window retains the exact finite value.
\end{theorem}
\begin{proof}
Read the actual positive split and restrict its selected
record on the given subwindow. Propagate the tangent through
the tested head and use the current supported strict carrier
to normalize the tangent fields. Restrict on the selected
original members while carrying the whole born suffix.
Use the original local canonicality and values on the
literal restricted tower.
\end{proof}

\begin{theorem}[The next positive read from the original upper history]
\label{hir:g1-actual-chronological-next-positive}
\uses{hir:g1-chronological-companion-suffix,hir:g1-original-upper-born-full-zero}
\lean{Hironaka.HironakaG1ActualChronologicalNext.next_positive_of_ordinary_supported_tangent,Hironaka.HironakaG1ActualChronologicalNext.original_upper_chronological_next_positive}
\leanok
At an actual positive nonspecial endpoint of the upper
order companion, its next positive letter reads the literal
selected suffix from its own run head. The suffix is canonical
and carries the supported contact with the whole retained
ordered boundary. Both original positive letters have their
exact values at the matched survivor.
\end{theorem}
\begin{proof}
Extract the genuine order companion from the original upper
canonical history. Use the common supported tangent window
and transfer its tangent to the actual bar. Read the positive
split at its run head, retain the chronological boundary
and restrict the canonical suffix there. Compose its exact
letter with the original upper read.
\end{proof}

\begin{theorem}[The true mark and next finite zero contact on a chronological suffix]
\label{hir:g1-chronological-true-mark-zero}
\uses{hir:g1-chronological-companion-suffix,hir:g1-actual-chronological-next-positive,hir:g1-restarted-second-companion-zero}
\lean{Hironaka.HironakaG1ChronologicalTrueMarkZero.true_mark_of_retained_boundary_zero,Hironaka.HironakaG1ChronologicalTrueMarkZero.selected_suffix_true_mark_zero_contact,Hironaka.HironakaG1ChronologicalTrueMarkZero.exists_retained_zero_window_of_normal_power,Hironaka.HironakaG1ChronologicalTrueMarkZero.historical_selected_full_normal_power,Hironaka.HironakaG1ChronologicalTrueMarkZero.exists_selected_true_mark_canonical_window,Hironaka.HironakaG1ChronologicalTrueMarkZero.chronological_companion_true_mark_window}
\leanok
The genuine advancing positive head supplies a canonical selected
suffix on a smaller window once the actual selected supported record
has the required domain on its positive mask. There the companion has
the true mark, all retained exponents vanish, and the exact positive
letter is preserved. A nonspecial level-zero endpoint gives the
finite zero read on that suffix with all ordered slots retained. The
advanced cosupport gap transfers canonicality to the supported cut.
\end{theorem}
\begin{proof}
The inherited bar tangent contributes its full normal power to the
historical selected ideal. Along the supported transverse carrier
this power supplies nonfactor witnesses for every retained slot.
Shrink to their common open window and use the maximal order to
identify the true-mark companion with the restricted ideal. Restrict
the actual canonical suffix without changing its read. At the
nonspecial zero endpoint use the literal finite-contact formula and
transport the chronological value along the supported test.
\end{proof}

\begin{theorem}[Supported domain and true-mark gap along chronological contacts]
\label{hir:g1-advancing-support-power-gap}
\uses{hir:g1-chronological-true-mark-zero}
\lean{Hironaka.HironakaG1AdvancingSupportPowerGap.cosupport_gap_of_restrictEquiv,Hironaka.HironakaG1AdvancingSupportPowerGap.power_mark_one_has_gap,Hironaka.HironakaG1AdvancingSupportPowerGap.cosupport_gap_geoTransformAlong_of_mark_one_power,Hironaka.HironakaG1AdvancingSupportPowerGap.advancing_supported_canonical_of_mark_one_power,Hironaka.HironakaG1AdvancingSupportPowerGap.chronological_supported_domain_of_nonspecial_anchor,Hironaka.HironakaG1AdvancingSupportPowerGap.finite_zero_iff_off_retained_strict_boundary}
\leanok
A supported record whose restricted ideal is the whole positive
power of a mark-one record preserves the order gap through actual
controlled transforms, including occupied steps. This gives the
supported canonical transfer at its advancing masks. If the
supported record is nonspecial at the actual positive head, the
intersection of its mask with the upper mask gives both selected
domains locally. In the tested strict-contact setting the finite
zero branch occurs precisely off every retained strict boundary
slot, including slots born before the companion formed.
\end{theorem}
\begin{proof}
The power law multiplies every allowed order by the positive
power and commutes with controlled transforms. The larger
threshold therefore has the same cosupport as the true one.
Use the nonspeciality hypothesis to put the actual anchor in
both masks and apply their selected-domain reductions. The
strict supported trace gives nonspeciality at level zero, and
the literal original-slot formula reads that level as the
number of incident retained strict transforms.
\end{proof}

\begin{theorem}[The geometric domain split at an actual positive head]
\label{hir:g1-chronological-support-domain-split}
\uses{hir:g1-advancing-support-power-gap}
\lean{Hironaka.HironakaG1ChronologicalSupportDimensionSplit.analytic_trace_fills_iff_local_full_dimension,Hironaka.HironakaG1ChronologicalSupportDimensionSplit.supported_positive_anchor_dimension_split,Hironaka.HironakaG1ChronologicalSupportDimensionSplit.chronological_support_dimension_domain_split,Hironaka.HironakaG1ChronologicalSupportDomainObstruction.positive_special_head_no_selected_domain,Hironaka.HironakaG1ChronologicalSupportDomainObstruction.chronological_support_domain_iff_dim_drop}
\leanok
At the actual tested positive run head, the supported point is
special precisely when the analytic cosupport trace has full local
dimension in the selected carrier. Strict local dimension drop
gives a common positive submask on which both selected domain
bounds hold on smaller windows. Full dimension excludes a selected
recursive domain on any open through the anchor.
\end{theorem}
\begin{proof}
An analytic subset of a smooth carrier with full local dimension
contains that carrier as a germ. In the full-dimension case the
literal selected cosupport therefore fills the selected carrier
near the point, contradicting nonvanishing in any invariant
domain. In the strict-drop case both positive-head mask reductions
apply on the intersection of their actual opens.
\end{proof}

\begin{theorem}[A canonical point blowup can lose the source order gap]
\label{hir:g1-canonical-controlled-gap-counterexample}
\uses{hir:g1-advancing-support-power-gap,hir:g1-recursive-born-finite-example}
\lean{Hironaka.HironakaG1ControlledGapCounterexample.canonical_point_blowup_loses_gap}
\leanok
The mark-two coordinate ideal generated by $(x^2,y^3,z^2)$
has the source cosupport gap and a genuine canonical point
blowup. Its actual controlled chart has an ordinary order-one
survivor, so the controlled gap fails at that stage.
\end{theorem}
\begin{proof}
Compute the source ordinary order at the origin and away from
it. The unique maximal centre is the origin. On the middle
chart, division by the exceptional factor at mark two gives
the span of $(a^2,b,c^2)$. At the survivor the linear $b$
has order one, which witnesses failure of the next threshold.
\end{proof}

\begin{theorem}[The two branches at the actual positive head]
\label{hir:g1-exhaustive-actual-positive-head}
\uses{hir:g1-chronological-support-domain-split,hir:g1-chronological-true-mark-zero}
\lean{Hironaka.HironakaG1ExhaustivePositiveHead.regular_chronological_head_true_mark,Hironaka.HironakaG1ExhaustivePositiveHead.selected_suffix_finite_zero_of_strict_avoidance_and_power,Hironaka.HironakaG1ExhaustivePositiveHead.exhaustive_actual_positive_head}
\leanok
Full local dimension at the actual positive head gives the
supported top letter and rules out a finite supported domain
through the anchor. Strict dimension drop gives a smaller
actual window carrying the canonical true-mark suffix and its
original positive letter. On that suffix the finite zero
continuation applies when the endpoint avoids every retained
strict slot and the supported order-mask record is a whole
power of a mark-one record on its carrier.
\end{theorem}
\begin{proof}
Use the exact dimension criterion at the actual head. In the
strict-drop case choose the common domain and tangent window,
then use the full normal power to remove all retained factors
on a smaller true-mark chart. The actual slot-avoidance formula
identifies the zero branch. The whole-power law supplies the
controlled gap for the supported canonical transfer.
\end{proof}

\begin{theorem}[Local born orders and continuation at a true mark]
\label{hir:g1-local-born-order-continuation}
\uses{hir:g1-exhaustive-actual-positive-head,hir:g1-canonical-controlled-gap-counterexample}
\lean{Hironaka.HironakaG1LocalBornOrderContinuation.selected_suffix_finite_zero_of_local_born_order,Hironaka.HironakaG1LocalBornOrderContinuation.chronological_supported_source_gap,Hironaka.HironakaG1LocalBornOrderContinuation.cosupport_gap_of_current_coefficient_presentation,Hironaka.HironakaG1LocalBornOrderContinuation.true_mark_positive_continuation_of_surviving_slot}
\leanok
For the supported selected coefficient member the source gap
is automatic. At the actual advancing zero head it suffices
to bound the supported exponent on a born incident slot by
the supported mark when the contact exponent is zero and the
supported exponent is positive. Alternatively an exact
current coefficient presentation supplies the gap. If a
retained strict slot survives, the actual inner level is
positive and yields its special top or finite canonical suffix
read with the original true mark and chronological slots.
\end{theorem}
\begin{proof}
Use the ordinary gap of the selected coefficient presentation.
Transport all old-slot exponents through the actual controlled
history. Only the stated born divergence requires a lower
bound for the current gap. Apply the advancing zero theorem.
If a retained slot survives, the original-slot count is
positive and the chronological positive read gives the two
actual successor cases.
\end{proof}

\begin{theorem}[Advanced supported comparison from actual coefficient witnesses]
\label{hir:g1-advancing-supported-coefficient-witnesses}
\uses{hir:g1-local-born-order-continuation,hir:g1-advancing-support-power-gap}
\lean{Hironaka.HironakaG1BornCoefficientTraceWitness.advancing_supported_canonical_of_finite_born_coefficient_witnesses,Hironaka.HironakaG1ChronologicalPositiveSupported.chronological_positive_supported_recursive_comparison}
\leanok
For the advancing zero branch, exact inherited coefficient
equivalence and finite refreshed generator membership on each
offending born trace give the normalized exponents and supported
canonicality. At an actual next positive nonspecial supported
anchor, the genuine local mask gives the two recursive contact
domains, compatible normalized exponents, and equivalent
recursive canonicality with the same strict order on common suffixes.
\end{theorem}
\begin{proof}
The exact source coefficient presentation supplies the initial
gap and normalized exponents. Literal controlled restriction
handles the inherited slots. The stated finite born generator
memberships give the advanced span on every offending trace.
For the positive branch, use the shared actual mask and tangent
transport. The coefficient records have automatic ordinary gaps,
which give the local normalized exponents and let the recursive
two-pair comparison identify their canonical continuations.
\end{proof}

\begin{theorem}[A submark born factor in actual coefficient geometry]
\label{hir:g1-actual-born-submark-coefficients}
\uses{hir:g1-advancing-supported-coefficient-witnesses,hir:g1-canonical-controlled-gap-counterexample}
\lean{Hironaka.HironakaG1BornSubmarkPoint.actual_born_submark_at_surviving_point,Hironaka.HironakaG1BornSubmarkPoint.actual_coefficient_refresh_removes_born_submark,Hironaka.HironakaG1BornSubmarkPoint.inherited_coefficient_born_submark_at_survivor,Hironaka.HironakaG1BornSubmarkPoint.exact_inherited_coefficient_born_refresh_mismatch,Hironaka.HironakaG1ProperBornStrictPlane.strictPlane_iff_middle_chart,Hironaka.HironakaG1ProperBornStrictPlane.horizontal_order_one_on_strictPlane,Hironaka.HironakaG1ProperBornStrictPlane.quotientGen_mem,Hironaka.HironakaG1ProperBornStrictPlane.controlled_span,Hironaka.HironakaG1ProperBornStrictPlane.strictPlane_survivor,Hironaka.HironakaG1ProperBornStrictPlane.strictPlane_born_order_one,Hironaka.HironakaG1ProperBornStrictPlane.plane_tangent,Hironaka.HironakaG1ProperBornStrictPlane.frequently_strictPlane_born_order_one,Hironaka.HironakaG1ProperBornStrictPlane.strictPlane_born_exponent}
\leanok
On the full carrier, a literal inherited mark-two coefficient
record has a surviving point with born exponent one after an
actual canonical point blowup. The refreshed coefficient ideal
has exponent zero in that same ordered slot. On an actual proper
zero-plane tangent, the raw mark-two controlled ideal of
$(x^3,y^4,z^2)$ also has a surviving point and born exponent one
on the true strict plane.
\end{theorem}
\begin{proof}
In the middle chart the real exceptional equation divides each
generator by the source mark. The controlled derivative-square
contributions vanish to order at least two along the born trace,
whereas a controlled raw generator has generic simple vanishing.
Its refreshed derivative contribution is nonvanishing along the
born slot. For the proper-plane calculation the source $z^2$ gives
the actual tangent. The true strict plane is the vanishing of the
third chart coordinate. There the same actual controlled
generators give a surviving origin and exact order one at
nearby exceptional points.
\end{proof}

\begin{theorem}[Inherited and refreshed coefficients on the literal strict plane]
\label{hir:g1-proper-inherited-refreshed-born-mismatch}
\uses{hir:g1-actual-born-submark-coefficients,hir:g1-advancing-supported-coefficient-witnesses}
\lean{Hironaka.HironakaG1ProperBornDerivativeTrace.baseIdeal_coherent,Hironaka.HironakaG1ProperBornDerivativeTrace.record_eq_padding,Hironaka.HironakaG1ProperBornDerivativeTrace.source_ambient_derivative_trace,Hironaka.HironakaG1ProperBornDerivativeTrace.strictPlane_closedSubmanifold,Hironaka.HironakaG1ProperBornDerivativeTrace.supportedStage_coherent,Hironaka.HironakaG1ProperBornDerivativeTrace.proper_refreshed_born_zero,Hironaka.HironakaG1ProperBornDerivativeTrace.controlled_source_derivative_on_born,Hironaka.HironakaG1ProperBornDerivativeTrace.inheritedStage_ideal,Hironaka.HironakaG1ProperBornDerivativeTrace.proper_inherited_vs_refreshed,Hironaka.HironakaG1ProperBornDerivativeTrace.proper_chronological_inherited_vs_refreshed}
\leanok
For the actual proper zero-plane contact of $(x^3,y^4,z^2)$,
the inherited coefficient record and the refreshed coefficient
record on the same admissible blowup both retain the actual
survivor and the same ordered born boundary. Their exponents
in that slot are one and zero, respectively.
\end{theorem}
\begin{proof}
Write the source as the actual normal padding of the plane
ideal and identify the ambient derivative trace. On the middle
chart divide the entire source coefficient ideal by the actual
exceptional square. The controlled source derivative-square
terms vanish to order at least two along the strict-plane born
trace. The raw controlled trace has exact factor one there.
Refreshing the coefficient ideal differentiates the simple
normal vanishing after control and gives a nonfactor witness.
The controlled restriction identities give these exponent
calculations on the two literal proper marked records.
\end{proof}

\begin{theorem}[Exact keys at the actual positive stage]
\label{hir:g1-actual-positive-stage-keys}
\uses{hir:g1-advancing-supported-coefficient-witnesses,hir:g1-proper-inherited-refreshed-born-mismatch}
\lean{Hironaka.HironakaG1ActualHighestPhaseKeys.later_positive_read_from_canonical,Hironaka.HironakaG1ActualHighestPhaseKeys.padded_maximum_from_actual_highest_keys,Hironaka.HironakaG1ActualHighestPhaseKeys.section_maximum_from_actual_highest_keys,Hironaka.HironakaG1ActualHighestPhaseKeys.mixed_special_maximum_iff_from_canonical}
\leanok
At a positive canonical stage, the finite read is witnessed by the
actual tested suffix. On the highest attained stratum, the padded
maximum uses its special case or the retained birth count and live
leaf. The section maximum uses its special case or the actual local
contact suffix. These criteria decide when the two loci agree in the
mixed special case.
\end{theorem}
\begin{proof}
Split each tested tower at its own regime anchor. The positive
canonical read supplies a local coefficient presentation on the
resulting suffix. Compare the realized values on the highest
attained stratum, with a separate comparison for special points.
\end{proof}

\begin{theorem}[Native realization of the positive selected section phase]
\label{hir:g1-native-positive-section-phase}
\uses{hir:g1-actual-positive-stage-keys}
\lean{Hironaka.HironakaG1NativePositivePhase.selected_section_supported_native_head,Hironaka.HironakaG1NativePositivePhase.exists_actual_positive_selected_native_phase}
\leanok
At a nonspecial point of positive level on an actual maximal-order
section canonical tower, the selected section and prescribed contact
records lift on one actual native suffix. At the head the carrier
is the restriction of the actual strict section. Both native records
are in the invariant domain and canonical on the lifted suffix.
The native transformed selected record is exact, and its cosupport
maps onto the upper selected cosupport. At the given point the upper
value is the native value through the same-bound dimension embedding.
The special branch uses the special top read of the preceding theorem.
\end{theorem}
\begin{proof}
Realize the section at the actual head. Selecting the retained ordered
members preserves transversality. Lift the localized tested suffix on
that section, then use the common carrier and ordered boundary of the
contact comparison to lift the prescribed contact record on the same
tower. Apply same-bound value transport and the whole-record and
cosupport identities of the actual lift.
\end{proof}

\begin{theorem}[Exact refreshed zero continuation and its canonical obstruction]
\label{hir:g1-exact-refreshed-zero-continuation}
\uses{hir:g1-proper-inherited-refreshed-born-mismatch}
\lean{Hironaka.HironakaG1ProperBornCanonicalMismatch.proper_raw_source_domain,Hironaka.HironakaG1ProperBornCanonicalMismatch.proper_point_canonical,Hironaka.HironakaG1ProperBornCanonicalMismatch.canonical_proper_coefficient_mismatch,Hironaka.HironakaG1ProperBornCanonicalMismatch.canonical_inherited_not_refresh_compatible,Hironaka.HironakaG1AdvancingRefreshedCoefficient.refreshed_restricted_coefficient_gap,Hironaka.HironakaG1AdvancingRefreshedCoefficient.advancing_supported_canonical_of_refreshed_coefficient,Hironaka.HironakaG1AdvancingRefreshedCoefficient.actual_zero_head_advances_supported_of_refresh,Hironaka.HironakaG1ChronologicalRefreshedZeroRead.selected_suffix_true_mark_zero_of_refresh}
\leanok
The proper-plane point blowup in the explicit example is canonical.
Its inherited controlled coefficient and refreshed coefficient fail
restricted equivalence on the strict plane. At a general advancing
head, exact restricted equivalence with the true refreshed
coefficient gives the source gap and supported canonicality. The
resulting zero read retains the actual tested suffix, true mark,
endpoint, and every ordered boundary slot.
\end{theorem}
\begin{proof}
Both initial proper-plane records have singleton cosupport. The
actual point blowup is therefore canonical, and its computed born
exponents separate the two restrictions. In the general
continuation, exact refreshed equivalence transfers the ordinary
coefficient gap. Apply the supported chronological zero theorem
and its literal boundary and endpoint identities.
\end{proof}

\begin{theorem}[The actual positive order companion on a native lift]
\label{hir:g1-actual-positive-order-companion-lift}
\uses{hir:g1-native-positive-section-phase,hir:g1-exact-refreshed-zero-continuation}
\lean{Hironaka.HironakaG1ActualNativeCompanionPhase.positive_companion_witness_slot_count_and_special,Hironaka.HironakaG1ActualNativeCompanionPhase.actual_positive_companion_pair_on_native_lift,Hironaka.HironakaG1ActualNativeCompanionPhase.actual_order_head_companion_geometry,Hironaka.HironakaG1NativeCompanionBranchGeometry.actual_masked_companions_maximal_and_special_on_native_lift}
\leanok
A supported common native lift of an original canonical history
has synchronized actual positive-order companion clocks. Its
masked head records have maximal order. Recasting the lift at
the actual cut identifies the lower masks and whole companion
restriction. At the surviving point, the restricted tails have
the same specialness and the same count of incident retained
slots, including those born before the cut. The recursive values
agree through the same-bound dimension embedding.
\end{theorem}
\begin{proof}
Split at the first cut of the common actual order regime.
Use native restriction of the whole transformed records at that
head, then restrict to the actual projected open. The invariant
domain supplies maximal order for the masked companion there.
Its original ordered members give the boundary count. The two
literal recursive reads give the stated dimension transport.
\end{proof}

\begin{theorem}[Actual finite ratio lifts and comparison on the supported cut]
\label{hir:g1-ratio-supported-finite-comparison}
\uses{hir:g1-native-positive-section-phase}
\lean{Hironaka.HironakaG1RatioGraphFiniteTower.withN_blockAmbientMarked_eq,Hironaka.HironakaG1RatioGraphFiniteTower.ratioGraph_supported_domains,Hironaka.HironakaG1RatioGraphFiniteTower.exists_ratioGraph_supported_finite_tower,Hironaka.HironakaG1RatioGraphFiniteTestClass.blockAmbientMarked_locWPEquiv_blockMarked,Hironaka.HironakaG1RatioGraphFiniteTestClass.ratioGraph_sameTestClass,Hironaka.HironakaG1RatioGraphFiniteTestClass.ratioGraph_isGeoTestTower_iff,Hironaka.HironakaG1RatioGraphFiniteTestClass.exists_ratioGraph_finite_tower_of_cylinder,Hironaka.HironakaG1RatioGraphSupportedComparison.exists_ratioGraph_supported_canonical_read,Hironaka.HironakaG1RatioGraphSupportedComparison.supportedEndpoint_order_iff,Hironaka.HironakaG1RatioGraphSupportedComparison.supportedEndpoint_maximum_image_eq,Hironaka.HironakaG1RatioGraphSupportedComparison.exists_supportedEndpoint_valueOrderIso,Hironaka.HironakaG1RatioGraphEndpointAlignment.supportedEndpoint_alignment_iff,Hironaka.HironakaG1RatioGraphEndpointAlignment.supported_alignment_of_shared_native_tower,Hironaka.HironakaG1RatioGraphEndpointAlignment.supported_alignment_of_native_states,Hironaka.HironakaG1SharedNativeLift.nativeLift_heq,Hironaka.HironakaG1SharedNativeLift.sourceSupport_eq_of_commuting_square,Hironaka.HironakaG1SharedNativeLift.nativeLift_heq_of_source_square,Hironaka.HironakaG1SharedNativeLift.unpadded_transport_isTowerIso,Hironaka.HironakaG1SharedNativeLift.nativeLift_heq_of_nonempty_source_square,Hironaka.HironakaG1SharedNativeLift.unpadded_transport_isTowerIso_of_source_square}
\leanok
An actual finite padded ratio test tower gives a literal unpadded
native source-open tower. Its endpoint square commutes and preserves
the true mark and ordered boundary. For the invariant reset to the
supported carrier, canonicality, dimension-embedded values and
maximum loci transport. Alignment of the two supported endpoint
orders is equivalent to alignment on the paired native points.
On a shared native chart clock, common native states with exact
boundary and local codimension two produce such alignment.
A commuting biholomorphic square on the full source sections
determines an isomorphism of their lifted chart towers over that
square. On nonempty supports it also forces the same native dimension.
\end{theorem}
\begin{proof}
Induct on the count of normal coordinates for the supported marked
class and then on the actual ambient tower for the native lift.
The whole pulled-back controlled record identifies the supported
endpoint. Use the original family dimension embedding on the
reset record, and strict order transport on the paired endpoint
cut. Lift uniqueness follows by induction on actual ambient
projections and centres, and the full source square identifies
the strict support and the endpoint injections.
\end{proof}

\begin{theorem}[Ordinary native order evaluation on the actual section]
\label{hir:g1-ordinary-literal-section-evaluator}
\uses{hir:g1-actual-positive-order-companion-lift}
\lean{Hironaka.HironakaG1NativeOrderDispatcher.monomialRead,Hironaka.HironakaG1NativeOrderDispatcher.actual_native_order_dispatch,Hironaka.HironakaG1OrdinaryNativeEvaluatorEndpoint.ActualOrderHeadCoefficients,Hironaka.HironakaG1OrdinaryNativeEvaluatorEndpoint.actual_order_head_companion_coefficients,Hironaka.HironakaG1OrdinaryNativeEvaluatorEndpoint.exists_ordinary_native_order_evaluator,Hironaka.HironakaG1LiteralSectionOrderEndpoint.literal_section_source_geometry,Hironaka.HironakaG1LiteralSectionOrderEndpoint.exists_literal_section_lift_of_padded_canonical,Hironaka.HironakaG1LiteralSectionOrderEndpoint.exists_literal_section_original_evaluator,Hironaka.HironakaG1LiteralSectionOrderEndpoint.cartier_no_shared_nonempty_canonical_step}
\leanok
An ordinary canonical history supported on a transverse native
carrier produces its actual native canonical lift, preserving
the whole current record and exhaustion. The zero order read
uses the whole ordered boundary. The positive order read
transports both masked companions, the actual clock, the live
recursive value and its special top. On this head, the companion,
coefficient and homogenized coefficient records and all
tangential derivative ideals commute with native pullback.

The literal normal section has the required source geometry.
A padded upper canonical history gives a native literal section
test with its strict carrier and full record. On the literal
section's own original-family canonical history one obtains
the complete order evaluator and exhaustion. In the Cartier
example the common first prefix admits no nonempty extension
canonical for both padded and literal upper clocks.
\end{theorem}
\begin{proof}
The ordinary canonical history is a test tower. Lift that
actual history, then compare the original invariant at the
same bound. At the positive order cut, recast the supported
lift at its actual head and use the pullback formulae for the
companion and its tangential ideals. The normal section and
pulled-back original boundary yield its support and
transversality. The Cartier maximum inequality excludes the
same nonempty canonical centre on the two upper clocks.
\end{proof}

\begin{theorem}[Native phase reads on the actual literal and refreshed records]
\label{hir:g1-native-literal-phase-dispatch}
\uses{hir:g1-ordinary-literal-section-evaluator,hir:g1-native-positive-section-phase}
\lean{Hironaka.HironakaG1RefreshedNativeZeroLift.native_supported_canonical_read,Hironaka.HironakaG1RefreshedNativeZeroLift.refreshed_supported_native_read,Hironaka.HironakaG1RefreshedNativeCompanionRead.refreshed_zero_to_native_positive_companion,Hironaka.HironakaG1RefreshedNativeOrderDispatcher.refreshed_supported_native_order_dispatch,Hironaka.HironakaG1NativeSectionEvaluator.native_section_zero_evaluator,Hironaka.HironakaG1NativeSectionPhaseDispatch.native_literal_section_phase_dispatch}
\leanok
For the actual canonical literal section, the masked read
is top at a special point. In the nonspecial case, its actual
zero kept lower record or positive selected member gives
the appropriate native canonical suffix. Each finite branch
has the exact ordered boundary, whole transformed record
and value through the same-bound dimension embedding. On
an advancing supported suffix with exact refreshed-equivalence,
the native lift realizes both recursive order branches on the
actual restricted tower.
\end{theorem}
\begin{proof}
Split the literal step-one read by specialness and the
actual chronological level. At a finite branch compare the
actual selected or kept lower read, and lift its test tail
on the genuine section. For a refreshed supported record,
the ordinary coefficient gap supplies supported canonicality.
The resulting lift satisfies the ordinary order dispatcher.
\end{proof}

\begin{theorem}[Literal shared support and composed native paddings]
\label{hir:g1-literal-shared-support-and-padding}
\uses{hir:g1-ratio-supported-finite-comparison}
\lean{Hironaka.HironakaG1LiteralRatioFullSupport.fullSupport_eq_iff,Hironaka.HironakaG1LiteralRatioFullSupport.pairedCut_subset_fullSupports,Hironaka.HironakaG1LiteralRatioFullSupport.dimension_eq_of_fullSupport_eq,Hironaka.HironakaG1LiteralRatioFullSupport.literalRatio_shared_lifts,Hironaka.HironakaG1LiteralRatioFullSupport.supported_order_on_literal_shared_lift,Hironaka.HironakaG1FiniteNativePaddingComposition.familyDimensionMap_chain,Hironaka.HironakaG1FiniteNativePaddingComposition.familyDimensionMap_chain_offsets,Hironaka.HironakaG1FiniteNativePaddingComposition.realized_padding_chain,Hironaka.HironakaG1FiniteNativePaddingComposition.exists_ratio_supported_nativeChartRealization,Hironaka.HironakaG1FiniteNativePaddingComposition.dimension_eq_of_equal_free_product_profiles}
\leanok
Equality of the two literal full supports is the equality of
their full graph equations on the common window. The paired
cut lies in both supports. If the full supports coincide and
are nonempty, their native dimensions agree, and the shared
lift compares the supported canonical histories under the
native class, dimension, and normalized exponent hypotheses.
For two nested realized supports, the dimension embeddings
compose, and the endpoint cosupports and maximum loci
transport through their composite. A literal ratio tower
supplies an actual first realization. Free products of the
same nonzero Hilbert--Samuel profile preserve equality of
profiles only with equal free dimensions.
\end{theorem}
\begin{proof}
Evaluate the full graph equations on the common window.
Apply lift uniqueness and the same native chart comparison
to the identical full support. In the nested case apply
whole-record and dimension-embedding transport twice and
compose the endpoint maps. For free products, compare the
degree-one profile terms.
\end{proof}

\begin{theorem}[The actual selected record supplies the recursive native evaluator]
\label{hir:g1-selected-section-ordinary-native}
\uses{hir:g1-native-literal-phase-dispatch,hir:g1-ordinary-literal-section-evaluator}
\lean{Hironaka.HironakaG1SelectedSectionOrdinaryRead.actual_selected_section_ordinary_source,Hironaka.HironakaG1SelectedSectionOrdinaryRead.literal_section_recursive_native_dispatch,Hironaka.HironakaG1SelectedSectionOrdinaryRead.literal_section_depth_zero_dispatch}
\leanok
In either finite branch of the actual literal section phase, the
selected ambient lower record carries its original-family canonical
test and actual surviving endpoint on the same restricted suffix.
Its geometric section support and transversality produce one native
lift bearing the whole transformed record and exact outer read.
On that same lift the next recursive read is the full-boundary order
zero letter or the actual positive companion read. The positive
companion includes the exact coefficient records and tangential
differential ideals. At depth zero every surviving point is special.
\end{theorem}
\begin{proof}
Take the selected record from the literal zero kept or positive
member branch. Use its own canonical suffix as the input to the
ordinary native evaluator. The arena and source record remain the
actual ones, hence the outer read and both recursive branches
use exactly the lift that the evaluator produces. At depth zero
the lower invariant domain forbids a surviving finite branch.
\end{proof}

\begin{theorem}[Paired native profiles and recursive companion values]
\label{hir:g1-paired-native-profile-companion-values}
\uses{hir:g1-literal-shared-support-and-padding}
\lean{Hironaka.HironakaG1LiteralRatioPointProfiles.hs_of_actual_openImage,Hironaka.HironakaG1LiteralRatioPointProfiles.hs_blockGraph_point,Hironaka.HironakaG1LiteralRatioPointProfiles.actual_ratioCut_point_profiles,Hironaka.HironakaG1LiteralRatioPointProfiles.native_section_preimage_profile,Hironaka.HironakaG1LiteralRatioPointProfiles.paired_strict_section_profiles,Hironaka.HironakaG1LiteralRatioPointProfiles.paired_lift_endpoint_profiles,Hironaka.HironakaG1PairedLiftCompanionValues.paired_lift_value_of_ambient_value,Hironaka.HironakaG1PairedLiftCompanionValues.paired_same_record_positive_companions,Hironaka.HironakaG1PairedEndpointMaxima.paired_same_record_canonical_max_and_values}
\leanok
The literal reduced ratio cut and both original reader points
have equal Hilbert--Samuel profiles in every degree. The same
equality holds on paired endpoints of an actual common
supported test, including different native dimensions and
marks. Two actual lifts of one supported marked canonical
history have the same companion order cut and actual mark.
Their positive companion values agree through the relative
dimension map. Native canonicality and maximum-locus images
agree on the paired endpoints. For independent records,
equality of the ambient evaluator values suffices for the
relative value comparison.
\end{theorem}
\begin{proof}
Use the literal analytic cut isomorphisms and invariance
of intrinsic vanishing-ideal jets, also on the transformed
strict supports. Transport the supported original values
along both actual lifts. Cancel the common dimension map
at the paired endpoint. These exact order embeddings preserve
maximum membership and compare the occupied canonical cuts.
\end{proof}

\begin{theorem}[Canonical stopping on independently presented histories]
\label{hir:g1-exact-independent-canonical-stop}
\uses{hir:g1-paired-native-profile-companion-values,hir:g1-ordinary-literal-section-evaluator}
\lean{Hironaka.HironakaG1PairedCanonicalStop.RejectedCentre,Hironaka.HironakaG1PairedCanonicalStop.canonical_extension_iff_maximum_agreement,Hironaka.HironakaG1PairedCanonicalStop.canonical_prefix_or_first_rejected_centre,Hironaka.HironakaG1PairedCanonicalStop.literal_section_paired_canonical_stop,Hironaka.HironakaG1PairedCanonicalStop.cartier_exact_paired_stop,Hironaka.HironakaG1PairedEndpointMaxima.paired_maximum_anchor_guard,Hironaka.HironakaG1PairedEndpointMaxima.paired_maximum_no_anchor_guard,Hironaka.HironakaG1PairedEndpointMaxima.paired_occupied_canonical_extension_guard,Hironaka.HironakaG1PointProfileRecursiveObstruction.cartier_controlled_cut_eq,Hironaka.HironakaG1PointProfileRecursiveObstruction.cartier_controlled_cut_profile,Hironaka.HironakaG1PointProfileRecursiveObstruction.cartier_no_stepOne_recode_from_profile,Hironaka.HironakaG1PointProfileRecursiveObstruction.cartier_tied_profiles_strict_padded_order}
\leanok
For independently presented actual histories, the exact
occupied continuation criterion compares the active anchor's
dominance and the full tie set. The no-maximum case has an
equivalent strict competitor criterion. On a padded canonical
tower the literal section is canonical up to its first rejected
centre, where the actual maximum loci diverge. The Cartier
case has a strict rejection just after the common first point
blowup for every occupied padded continuation. Two actual
survivors there have identical whole reduced profiles and
literal section Step One values and strictly different padded
Step One values.
\end{theorem}
\begin{proof}
Read canonicality at the two actual endpoint records. A
nonempty highest fibre is determined by its anchor and ties.
If no highest value is attained, each active point has a
strict competitor. Induct on the actual physical tower to
locate the first rejected occupied cut. In the Cartier chart
the actual cut is a regular curve and its full reduced
profile is computed. The ordered born slots separate the
two exact padded values on that curve.
\end{proof}

\begin{theorem}[The finite closed graph comparison on the original record]
\label{hir:g1-finite-closed-graph-original-comparison}
\uses{hir:g1-selected-section-ordinary-native,hir:g1-paired-native-profile-companion-values,hir:g1-exact-independent-canonical-stop}
\lean{Hironaka.HironakaG1FiniteClosedGraphOriginalComparison.blockAmbientMarked_isInvDomain,Hironaka.HironakaG1FiniteClosedGraphOriginalComparison.exists_closed_graph_actual_native_realization}
\leanok
An ordinary original marked ideal and any finite-dimensional
holomorphic graph give a supported ambient record in the original
invariant domain. From any actual full-cylinder test tower there
is an actual native chart realization. Its source marked ideal
is the original record at its original mark. Its endpoint map
is injective, its projection commutes with the original tower,
and its length is the length of that tower. Native canonicality
is equivalent to ambient canonicality on the graph-supported
record. At the actual surviving endpoints its original values
are embedded by the true finite normal dimension shift, and
the ambient maximum locus is exactly the image of the native
maximum locus, with its exact inverse image.
\end{theorem}
\begin{proof}
Iterate the normal test-class comparison and realize the strict
graph section along the literal tower. The native pullback on
the source graph recovers the entire original marked record.
Canonical comparison of the actual lift transports the original
invariant, with the dimension map determined by the actual
ambient and native dimensions. Its strict monotonicity and the
exact cosupport map give the two maximum-locus identities.
\end{proof}
